\section{一八七五至一九一二账本节点的省份社会深描}

\label{sec:late-qing-ledger-depth-1875-1912}

本组按本卷账本的晚清事件和来源限度展开。改革制度、战争外交、教育、家庭与省份政治并非独立领域；每一项措施都必须经由经费、文书、交通、地方组织和不同社群才会成为经验。

\subsection{省政、财政与地方治理}

账本条目\texttt{EQ-1870-TIANJIN-MASSACRE}记录，同治九年（1870年），清、法国与天津发生“天津教案发生，法国领事、修女和中国民众死伤，引发外交危机”。账本将其背景说明为教会活动、地方谣言和条约保护关系高度紧张，可见结果为责任叙述在中法档案中差异显著，须避免单一归因。这一节点虽可放入改革、外交或革命的宏观年表，却不能脱离省份财政、地方官、商会、学校、军队、家庭与交通网络来理解；同一上谕、条约或战事在不同地域的实际后果并不一致。

核心材料为《清穆宗实录》同治九年相关条；《筹办夷务始末》相关年卷与中外照会、条约文本。它们可以支持年序、制度文本、机构设立或某些地点的行动，却不能自动代表全体居民的收入、教育、身份、性别关系或政治意愿。账本提示“以同治九年（1870年）的同期文书为定位起点；官修记录、奏报或外交文本服务特定机构立场，具体日期、规模、地方执行与对手视角须以独立材料复核。”。因此，官方统计、外交措辞、报刊论战、教会叙述和军报数字必须分别置于生成它们的立场与制度目的之中，不能互相替代。

在地方尺度上，这一事件应被看作资源和信息重新配置的过程：经费能否落实，取决于税源和中介；学校、警政或新军能否运行，取决于人员、土地和日常协作；家庭与社群对改革的经历，则可能因性别、职业、族属、居处与迁徙而不同。现存材料只照亮少数路径，叙事须保留被登记者与未被登记者之间的差距。

账本条目\texttt{EQ-1870-TIANJIN-MASSACRE-DOCUMENTS}记录，同治九年（1870年），清、法国与天津发生“围绕「天津教案发生，法国领事、修女和中国民众死伤，引发外交危机」的照会、条约文本、换约或地方执行报告分别进入外交文书链”。账本将其背景说明为核心事件「天津教案发生，法国领事、修女和中国民众死伤，引发外交危机」要求中央或地方通过文书处理资源、信息或责任。，可见结果为该记录为后续按具体谕旨、奏折、条约、报刊或地方档案拆分事件提供可回查入口。。这一节点虽可放入改革、外交或革命的宏观年表，却不能脱离省份财政、地方官、商会、学校、军队、家庭与交通网络来理解；同一上谕、条约或战事在不同地域的实际后果并不一致。

核心材料为《清穆宗实录》同治九年相关条；《筹办夷务始末》相关年卷与中外照会、条约文本。它们可以支持年序、制度文本、机构设立或某些地点的行动，却不能自动代表全体居民的收入、教育、身份、性别关系或政治意愿。账本提示“此行仅确认相关文书链和可追查的行政动作；不得由其直接推断政策有效、地方服从、民众态度或单一因果。”。因此，官方统计、外交措辞、报刊论战、教会叙述和军报数字必须分别置于生成它们的立场与制度目的之中，不能互相替代。

在地方尺度上，这一事件应被看作资源和信息重新配置的过程：经费能否落实，取决于税源和中介；学校、警政或新军能否运行，取决于人员、土地和日常协作；家庭与社群对改革的经历，则可能因性别、职业、族属、居处与迁徙而不同。现存材料只照亮少数路径，叙事须保留被登记者与未被登记者之间的差距。

账本条目\texttt{EQ-1870-ZHILI-FAMINE-RELIEF}记录，同治九年（1870年），清与华北发生“华北灾荒和赈务问题促使官员、士绅和国际救济网络介入”。账本将其背景说明为旱灾、粮价、交通和地方储备共同影响灾情，可见结果为慈善报告与官府奏报的覆盖范围和数字需核对。这一节点虽可放入改革、外交或革命的宏观年表，却不能脱离省份财政、地方官、商会、学校、军队、家庭与交通网络来理解；同一上谕、条约或战事在不同地域的实际后果并不一致。

核心材料为《清穆宗实录》同治九年相关条；地方志、督抚奏折及地方档案。它们可以支持年序、制度文本、机构设立或某些地点的行动，却不能自动代表全体居民的收入、教育、身份、性别关系或政治意愿。账本提示“以同治九年（1870年）的同期文书为定位起点；官修记录、奏报或外交文本服务特定机构立场，具体日期、规模、地方执行与对手视角须以独立材料复核。”。因此，官方统计、外交措辞、报刊论战、教会叙述和军报数字必须分别置于生成它们的立场与制度目的之中，不能互相替代。

在地方尺度上，这一事件应被看作资源和信息重新配置的过程：经费能否落实，取决于税源和中介；学校、警政或新军能否运行，取决于人员、土地和日常协作；家庭与社群对改革的经历，则可能因性别、职业、族属、居处与迁徙而不同。现存材料只照亮少数路径，叙事须保留被登记者与未被登记者之间的差距。

账本条目\texttt{EQ-1870-ZHILI-FAMINE-RELIEF-DOCUMENTS}记录，同治九年（1870年），清与华北发生“围绕「华北灾荒和赈务问题促使官员、士绅和国际救济网络介入」的地方呈报、案件或赈务记录将社会反应带入官府文书链”。账本将其背景说明为核心事件「华北灾荒和赈务问题促使官员、士绅和国际救济网络介入」要求中央或地方通过文书处理资源、信息或责任。，可见结果为该记录为后续按具体谕旨、奏折、条约、报刊或地方档案拆分事件提供可回查入口。。这一节点虽可放入改革、外交或革命的宏观年表，却不能脱离省份财政、地方官、商会、学校、军队、家庭与交通网络来理解；同一上谕、条约或战事在不同地域的实际后果并不一致。

核心材料为《清穆宗实录》同治九年相关条；地方志、督抚奏折及地方档案。它们可以支持年序、制度文本、机构设立或某些地点的行动，却不能自动代表全体居民的收入、教育、身份、性别关系或政治意愿。账本提示“此行仅确认相关文书链和可追查的行政动作；不得由其直接推断政策有效、地方服从、民众态度或单一因果。”。因此，官方统计、外交措辞、报刊论战、教会叙述和军报数字必须分别置于生成它们的立场与制度目的之中，不能互相替代。

在地方尺度上，这一事件应被看作资源和信息重新配置的过程：经费能否落实，取决于税源和中介；学校、警政或新军能否运行，取决于人员、土地和日常协作；家庭与社群对改革的经历，则可能因性别、职业、族属、居处与迁徙而不同。现存材料只照亮少数路径，叙事须保留被登记者与未被登记者之间的差距。

账本条目\texttt{EQ-1871-MUDAN-INCIDENT}记录，同治十年（1871年），清、琉球、日本与台湾原住民发生“琉球船民在台湾牡丹社一带遇害，事件成为日本对台政策借口”。账本将其背景说明为海难、原住民地域与清朝地方治理边界交错，可见结果为事件叙述须区分原住民行动、清方管辖与日方外交利用。这一节点虽可放入改革、外交或革命的宏观年表，却不能脱离省份财政、地方官、商会、学校、军队、家庭与交通网络来理解；同一上谕、条约或战事在不同地域的实际后果并不一致。

核心材料为《清穆宗实录》同治十年相关条；《筹办夷务始末》相关年卷与中外照会、条约文本。它们可以支持年序、制度文本、机构设立或某些地点的行动，却不能自动代表全体居民的收入、教育、身份、性别关系或政治意愿。账本提示“以同治十年（1871年）的同期文书为定位起点；官修记录、奏报或外交文本服务特定机构立场，具体日期、规模、地方执行与对手视角须以独立材料复核。”。因此，官方统计、外交措辞、报刊论战、教会叙述和军报数字必须分别置于生成它们的立场与制度目的之中，不能互相替代。

在地方尺度上，这一事件应被看作资源和信息重新配置的过程：经费能否落实，取决于税源和中介；学校、警政或新军能否运行，取决于人员、土地和日常协作；家庭与社群对改革的经历，则可能因性别、职业、族属、居处与迁徙而不同。现存材料只照亮少数路径，叙事须保留被登记者与未被登记者之间的差距。

账本条目\texttt{EQ-1871-MUDAN-INCIDENT-DOCUMENTS}记录，同治十年（1871年），清、琉球、日本与台湾原住民发生“围绕「琉球船民在台湾牡丹社一带遇害，事件成为日本对台政策借口」的照会、条约文本、换约或地方执行报告分别进入外交文书链”。账本将其背景说明为核心事件「琉球船民在台湾牡丹社一带遇害，事件成为日本对台政策借口」要求中央或地方通过文书处理资源、信息或责任。，可见结果为该记录为后续按具体谕旨、奏折、条约、报刊或地方档案拆分事件提供可回查入口。。这一节点虽可放入改革、外交或革命的宏观年表，却不能脱离省份财政、地方官、商会、学校、军队、家庭与交通网络来理解；同一上谕、条约或战事在不同地域的实际后果并不一致。

核心材料为《清穆宗实录》同治十年相关条；《筹办夷务始末》相关年卷与中外照会、条约文本。它们可以支持年序、制度文本、机构设立或某些地点的行动，却不能自动代表全体居民的收入、教育、身份、性别关系或政治意愿。账本提示“此行仅确认相关文书链和可追查的行政动作；不得由其直接推断政策有效、地方服从、民众态度或单一因果。”。因此，官方统计、外交措辞、报刊论战、教会叙述和军报数字必须分别置于生成它们的立场与制度目的之中，不能互相替代。

在地方尺度上，这一事件应被看作资源和信息重新配置的过程：经费能否落实，取决于税源和中介；学校、警政或新军能否运行，取决于人员、土地和日常协作；家庭与社群对改革的经历，则可能因性别、职业、族属、居处与迁徙而不同。现存材料只照亮少数路径，叙事须保留被登记者与未被登记者之间的差距。

账本条目\texttt{EQ-1871-QING-STUDENTS-ABROAD}记录，同治十年（1871年），清与美国发生“清廷派遣幼童赴美留学的计划启动，洋务教育进入海外培训阶段”。账本将其背景说明为技术、语言与外交人才需求上升，可见结果为选派和留学经历不代表全国教育制度已转型。这一节点虽可放入改革、外交或革命的宏观年表，却不能脱离省份财政、地方官、商会、学校、军队、家庭与交通网络来理解；同一上谕、条约或战事在不同地域的实际后果并不一致。

核心材料为《清穆宗实录》同治十年相关条；内阁档、文集或报刊相关记录。它们可以支持年序、制度文本、机构设立或某些地点的行动，却不能自动代表全体居民的收入、教育、身份、性别关系或政治意愿。账本提示“以同治十年（1871年）的同期文书为定位起点；官修记录、奏报或外交文本服务特定机构立场，具体日期、规模、地方执行与对手视角须以独立材料复核。”。因此，官方统计、外交措辞、报刊论战、教会叙述和军报数字必须分别置于生成它们的立场与制度目的之中，不能互相替代。

在地方尺度上，这一事件应被看作资源和信息重新配置的过程：经费能否落实，取决于税源和中介；学校、警政或新军能否运行，取决于人员、土地和日常协作；家庭与社群对改革的经历，则可能因性别、职业、族属、居处与迁徙而不同。现存材料只照亮少数路径，叙事须保留被登记者与未被登记者之间的差距。

账本条目\texttt{EQ-1871-QING-STUDENTS-ABROAD-DOCUMENTS}记录，同治十年（1871年），清与美国发生“围绕「清廷派遣幼童赴美留学的计划启动，洋务教育进入海外培训阶段」的禁令、编纂、教育或传播活动留下中央和地方文本记录”。账本将其背景说明为核心事件「清廷派遣幼童赴美留学的计划启动，洋务教育进入海外培训阶段」要求中央或地方通过文书处理资源、信息或责任。，可见结果为该记录为后续按具体谕旨、奏折、条约、报刊或地方档案拆分事件提供可回查入口。。这一节点虽可放入改革、外交或革命的宏观年表，却不能脱离省份财政、地方官、商会、学校、军队、家庭与交通网络来理解；同一上谕、条约或战事在不同地域的实际后果并不一致。

核心材料为《清穆宗实录》同治十年相关条；内阁档、文集或报刊相关记录。它们可以支持年序、制度文本、机构设立或某些地点的行动，却不能自动代表全体居民的收入、教育、身份、性别关系或政治意愿。账本提示“此行仅确认相关文书链和可追查的行政动作；不得由其直接推断政策有效、地方服从、民众态度或单一因果。”。因此，官方统计、外交措辞、报刊论战、教会叙述和军报数字必须分别置于生成它们的立场与制度目的之中，不能互相替代。

在地方尺度上，这一事件应被看作资源和信息重新配置的过程：经费能否落实，取决于税源和中介；学校、警政或新军能否运行，取决于人员、土地和日常协作；家庭与社群对改革的经历，则可能因性别、职业、族属、居处与迁徙而不同。现存材料只照亮少数路径，叙事须保留被登记者与未被登记者之间的差距。

账本条目\texttt{EQ-1872-CHINA-MERCHANTS-STEAM-NAVIGATION}记录，同治十一年（1872年），清与上海商人发生“轮船招商局创办，试图以官督商办方式发展近代航运”。账本将其背景说明为外轮竞争和沿海、长江运输需求推动新企业，可见结果为企业性质、资本来源和经营绩效需依账册分析。这一节点虽可放入改革、外交或革命的宏观年表，却不能脱离省份财政、地方官、商会、学校、军队、家庭与交通网络来理解；同一上谕、条约或战事在不同地域的实际后果并不一致。

核心材料为《清穆宗实录》同治十一年相关条；户部奏销、盐政、河工或海关档案。它们可以支持年序、制度文本、机构设立或某些地点的行动，却不能自动代表全体居民的收入、教育、身份、性别关系或政治意愿。账本提示“以同治十一年（1872年）的同期文书为定位起点；官修记录、奏报或外交文本服务特定机构立场，具体日期、规模、地方执行与对手视角须以独立材料复核。”。因此，官方统计、外交措辞、报刊论战、教会叙述和军报数字必须分别置于生成它们的立场与制度目的之中，不能互相替代。

在地方尺度上，这一事件应被看作资源和信息重新配置的过程：经费能否落实，取决于税源和中介；学校、警政或新军能否运行，取决于人员、土地和日常协作；家庭与社群对改革的经历，则可能因性别、职业、族属、居处与迁徙而不同。现存材料只照亮少数路径，叙事须保留被登记者与未被登记者之间的差距。

账本条目\texttt{EQ-1872-CHINA-MERCHANTS-STEAM-NAVIGATION-DOCUMENTS}记录，同治十一年（1872年），清与上海商人发生“围绕「轮船招商局创办，试图以官督商办方式发展近代航运」的经费、税收、赈济或工程呈报经由户部与地方衙门处理”。账本将其背景说明为核心事件「轮船招商局创办，试图以官督商办方式发展近代航运」要求中央或地方通过文书处理资源、信息或责任。，可见结果为该记录为后续按具体谕旨、奏折、条约、报刊或地方档案拆分事件提供可回查入口。。这一节点虽可放入改革、外交或革命的宏观年表，却不能脱离省份财政、地方官、商会、学校、军队、家庭与交通网络来理解；同一上谕、条约或战事在不同地域的实际后果并不一致。

核心材料为《清穆宗实录》同治十一年相关条；户部奏销、盐政、河工或海关档案。它们可以支持年序、制度文本、机构设立或某些地点的行动，却不能自动代表全体居民的收入、教育、身份、性别关系或政治意愿。账本提示“此行仅确认相关文书链和可追查的行政动作；不得由其直接推断政策有效、地方服从、民众态度或单一因果。”。因此，官方统计、外交措辞、报刊论战、教会叙述和军报数字必须分别置于生成它们的立场与制度目的之中，不能互相替代。

在地方尺度上，这一事件应被看作资源和信息重新配置的过程：经费能否落实，取决于税源和中介；学校、警政或新军能否运行，取决于人员、土地和日常协作；家庭与社群对改革的经历，则可能因性别、职业、族属、居处与迁徙而不同。现存材料只照亮少数路径，叙事须保留被登记者与未被登记者之间的差距。

账本条目\texttt{EQ-1873-TONGZHI-EMPEROR-REIGN}记录，同治十二年（1873年），清发生“同治帝亲政，垂帘体制与皇帝个人政治角色发生调整”。账本将其背景说明为幼帝成年改变摄政安排，可见结果为实际决策仍受太后、军机和重臣网络制约。这一节点虽可放入改革、外交或革命的宏观年表，却不能脱离省份财政、地方官、商会、学校、军队、家庭与交通网络来理解；同一上谕、条约或战事在不同地域的实际后果并不一致。

核心材料为《清穆宗实录》同治十二年相关条；宫中朱批奏折、内阁大库题本与《清史稿》相关志传。它们可以支持年序、制度文本、机构设立或某些地点的行动，却不能自动代表全体居民的收入、教育、身份、性别关系或政治意愿。账本提示“以同治十二年（1873年）的同期文书为定位起点；官修记录、奏报或外交文本服务特定机构立场，具体日期、规模、地方执行与对手视角须以独立材料复核。”。因此，官方统计、外交措辞、报刊论战、教会叙述和军报数字必须分别置于生成它们的立场与制度目的之中，不能互相替代。

在地方尺度上，这一事件应被看作资源和信息重新配置的过程：经费能否落实，取决于税源和中介；学校、警政或新军能否运行，取决于人员、土地和日常协作；家庭与社群对改革的经历，则可能因性别、职业、族属、居处与迁徙而不同。现存材料只照亮少数路径，叙事须保留被登记者与未被登记者之间的差距。

账本条目\texttt{EQ-1873-TONGZHI-EMPEROR-REIGN-DOCUMENTS}记录，同治十二年（1873年），清发生“围绕「同治帝亲政，垂帘体制与皇帝个人政治角色发生调整」的上谕、奏折与部议在中央—地方行政链中流转”。账本将其背景说明为核心事件「同治帝亲政，垂帘体制与皇帝个人政治角色发生调整」要求中央或地方通过文书处理资源、信息或责任。，可见结果为该记录为后续按具体谕旨、奏折、条约、报刊或地方档案拆分事件提供可回查入口。。这一节点虽可放入改革、外交或革命的宏观年表，却不能脱离省份财政、地方官、商会、学校、军队、家庭与交通网络来理解；同一上谕、条约或战事在不同地域的实际后果并不一致。

核心材料为《清穆宗实录》同治十二年相关条；宫中朱批奏折、内阁大库题本与《清史稿》相关志传。它们可以支持年序、制度文本、机构设立或某些地点的行动，却不能自动代表全体居民的收入、教育、身份、性别关系或政治意愿。账本提示“此行仅确认相关文书链和可追查的行政动作；不得由其直接推断政策有效、地方服从、民众态度或单一因果。”。因此，官方统计、外交措辞、报刊论战、教会叙述和军报数字必须分别置于生成它们的立场与制度目的之中，不能互相替代。

在地方尺度上，这一事件应被看作资源和信息重新配置的过程：经费能否落实，取决于税源和中介；学校、警政或新军能否运行，取决于人员、土地和日常协作；家庭与社群对改革的经历，则可能因性别、职业、族属、居处与迁徙而不同。现存材料只照亮少数路径，叙事须保留被登记者与未被登记者之间的差距。

账本条目\texttt{EQ-1873-JAPAN-RYUKYU-DISPUTE}记录，同治十二年（1873年），清、日本与琉球发生“日本围绕琉球与台湾事务加强对清交涉”。账本将其背景说明为日本明治国家扩张与东亚朝贡关系重组，可见结果为册封、属国和主权词汇应置于各方文本语境中。这一节点虽可放入改革、外交或革命的宏观年表，却不能脱离省份财政、地方官、商会、学校、军队、家庭与交通网络来理解；同一上谕、条约或战事在不同地域的实际后果并不一致。

核心材料为《清穆宗实录》同治十二年相关条；《筹办夷务始末》相关年卷与中外照会、条约文本。它们可以支持年序、制度文本、机构设立或某些地点的行动，却不能自动代表全体居民的收入、教育、身份、性别关系或政治意愿。账本提示“以同治十二年（1873年）的同期文书为定位起点；官修记录、奏报或外交文本服务特定机构立场，具体日期、规模、地方执行与对手视角须以独立材料复核。”。因此，官方统计、外交措辞、报刊论战、教会叙述和军报数字必须分别置于生成它们的立场与制度目的之中，不能互相替代。

在地方尺度上，这一事件应被看作资源和信息重新配置的过程：经费能否落实，取决于税源和中介；学校、警政或新军能否运行，取决于人员、土地和日常协作；家庭与社群对改革的经历，则可能因性别、职业、族属、居处与迁徙而不同。现存材料只照亮少数路径，叙事须保留被登记者与未被登记者之间的差距。

账本条目\texttt{EQ-1873-JAPAN-RYUKYU-DISPUTE-DOCUMENTS}记录，同治十二年（1873年），清、日本与琉球发生“围绕「日本围绕琉球与台湾事务加强对清交涉」的照会、条约文本、换约或地方执行报告分别进入外交文书链”。账本将其背景说明为核心事件「日本围绕琉球与台湾事务加强对清交涉」要求中央或地方通过文书处理资源、信息或责任。，可见结果为该记录为后续按具体谕旨、奏折、条约、报刊或地方档案拆分事件提供可回查入口。。这一节点虽可放入改革、外交或革命的宏观年表，却不能脱离省份财政、地方官、商会、学校、军队、家庭与交通网络来理解；同一上谕、条约或战事在不同地域的实际后果并不一致。

核心材料为《清穆宗实录》同治十二年相关条；《筹办夷务始末》相关年卷与中外照会、条约文本。它们可以支持年序、制度文本、机构设立或某些地点的行动，却不能自动代表全体居民的收入、教育、身份、性别关系或政治意愿。账本提示“此行仅确认相关文书链和可追查的行政动作；不得由其直接推断政策有效、地方服从、民众态度或单一因果。”。因此，官方统计、外交措辞、报刊论战、教会叙述和军报数字必须分别置于生成它们的立场与制度目的之中，不能互相替代。

在地方尺度上，这一事件应被看作资源和信息重新配置的过程：经费能否落实，取决于税源和中介；学校、警政或新军能否运行，取决于人员、土地和日常协作；家庭与社群对改革的经历，则可能因性别、职业、族属、居处与迁徙而不同。现存材料只照亮少数路径，叙事须保留被登记者与未被登记者之间的差距。

账本条目\texttt{EQ-1874-JAPANESE-TAIWAN-EXPEDITION}记录，同治十三年（1874年），清与日本发生“日本以牡丹社事件为由出兵台湾南部，清廷被迫应对”。账本将其背景说明为清廷对台湾东部和原住民地区的治理能力有限，可见结果为战后赔款与外交安排刺激清廷加强台湾经营。这一节点虽可放入改革、外交或革命的宏观年表，却不能脱离省份财政、地方官、商会、学校、军队、家庭与交通网络来理解；同一上谕、条约或战事在不同地域的实际后果并不一致。

核心材料为《清穆宗实录》同治十三年相关条；军机处档、战事方略及地方奏报。它们可以支持年序、制度文本、机构设立或某些地点的行动，却不能自动代表全体居民的收入、教育、身份、性别关系或政治意愿。账本提示“以同治十三年（1874年）的同期文书为定位起点；官修记录、奏报或外交文本服务特定机构立场，具体日期、规模、地方执行与对手视角须以独立材料复核。”。因此，官方统计、外交措辞、报刊论战、教会叙述和军报数字必须分别置于生成它们的立场与制度目的之中，不能互相替代。

在地方尺度上，这一事件应被看作资源和信息重新配置的过程：经费能否落实，取决于税源和中介；学校、警政或新军能否运行，取决于人员、土地和日常协作；家庭与社群对改革的经历，则可能因性别、职业、族属、居处与迁徙而不同。现存材料只照亮少数路径，叙事须保留被登记者与未被登记者之间的差距。

账本条目\texttt{EQ-1874-JAPANESE-TAIWAN-EXPEDITION-DOCUMENTS}记录，同治十三年（1874年），清与日本发生“清廷围绕「日本以牡丹社事件为由出兵台湾南部，清廷被迫应对」调遣军队、筹措饷械并处理战报，中央与地方的军政文书形成连续记录”。账本将其背景说明为核心事件「日本以牡丹社事件为由出兵台湾南部，清廷被迫应对」要求中央或地方通过文书处理资源、信息或责任。，可见结果为该记录为后续按具体谕旨、奏折、条约、报刊或地方档案拆分事件提供可回查入口。。这一节点虽可放入改革、外交或革命的宏观年表，却不能脱离省份财政、地方官、商会、学校、军队、家庭与交通网络来理解；同一上谕、条约或战事在不同地域的实际后果并不一致。

核心材料为《清穆宗实录》同治十三年相关条；军机处档、战事方略及地方奏报。它们可以支持年序、制度文本、机构设立或某些地点的行动，却不能自动代表全体居民的收入、教育、身份、性别关系或政治意愿。账本提示“此行仅确认相关文书链和可追查的行政动作；不得由其直接推断政策有效、地方服从、民众态度或单一因果。”。因此，官方统计、外交措辞、报刊论战、教会叙述和军报数字必须分别置于生成它们的立场与制度目的之中，不能互相替代。

在地方尺度上，这一事件应被看作资源和信息重新配置的过程：经费能否落实，取决于税源和中介；学校、警政或新军能否运行，取决于人员、土地和日常协作；家庭与社群对改革的经历，则可能因性别、职业、族属、居处与迁徙而不同。现存材料只照亮少数路径，叙事须保留被登记者与未被登记者之间的差距。

账本条目\texttt{EQ-1874-TONGZHI-DEATH}记录，同治十三年（1874年），清发生“同治帝去世，慈禧选择光绪帝继位，再次形成幼主政治”。账本将其背景说明为同治无嗣使宗法与宫廷权力问题交织，可见结果为慈禧重新垂帘，晚清中枢进入新阶段。这一节点虽可放入改革、外交或革命的宏观年表，却不能脱离省份财政、地方官、商会、学校、军队、家庭与交通网络来理解；同一上谕、条约或战事在不同地域的实际后果并不一致。

核心材料为《清穆宗实录》同治十三年相关条；宫中朱批奏折、内阁大库题本与《清史稿》相关志传。它们可以支持年序、制度文本、机构设立或某些地点的行动，却不能自动代表全体居民的收入、教育、身份、性别关系或政治意愿。账本提示“以同治十三年（1874年）的同期文书为定位起点；官修记录、奏报或外交文本服务特定机构立场，具体日期、规模、地方执行与对手视角须以独立材料复核。”。因此，官方统计、外交措辞、报刊论战、教会叙述和军报数字必须分别置于生成它们的立场与制度目的之中，不能互相替代。

在地方尺度上，这一事件应被看作资源和信息重新配置的过程：经费能否落实，取决于税源和中介；学校、警政或新军能否运行，取决于人员、土地和日常协作；家庭与社群对改革的经历，则可能因性别、职业、族属、居处与迁徙而不同。现存材料只照亮少数路径，叙事须保留被登记者与未被登记者之间的差距。

账本条目\texttt{EQ-1874-TONGZHI-DEATH-DOCUMENTS}记录，同治十三年（1874年），清发生“围绕「同治帝去世，慈禧选择光绪帝继位，再次形成幼主政治」的上谕、奏折与部议在中央—地方行政链中流转”。账本将其背景说明为核心事件「同治帝去世，慈禧选择光绪帝继位，再次形成幼主政治」要求中央或地方通过文书处理资源、信息或责任。，可见结果为该记录为后续按具体谕旨、奏折、条约、报刊或地方档案拆分事件提供可回查入口。。这一节点虽可放入改革、外交或革命的宏观年表，却不能脱离省份财政、地方官、商会、学校、军队、家庭与交通网络来理解；同一上谕、条约或战事在不同地域的实际后果并不一致。

核心材料为《清穆宗实录》同治十三年相关条；宫中朱批奏折、内阁大库题本与《清史稿》相关志传。它们可以支持年序、制度文本、机构设立或某些地点的行动，却不能自动代表全体居民的收入、教育、身份、性别关系或政治意愿。账本提示“此行仅确认相关文书链和可追查的行政动作；不得由其直接推断政策有效、地方服从、民众态度或单一因果。”。因此，官方统计、外交措辞、报刊论战、教会叙述和军报数字必须分别置于生成它们的立场与制度目的之中，不能互相替代。

在地方尺度上，这一事件应被看作资源和信息重新配置的过程：经费能否落实，取决于税源和中介；学校、警政或新军能否运行，取决于人员、土地和日常协作；家庭与社群对改革的经历，则可能因性别、职业、族属、居处与迁徙而不同。现存材料只照亮少数路径，叙事须保留被登记者与未被登记者之间的差距。

账本条目\texttt{EQ-1875-GUANGXU-ACCESSION}记录，光绪元年（1875年），清发生“光绪帝即位，慈禧、慈安继续掌握关键政治影响”。账本将其背景说明为幼主继承使摄政与名分安排再度出现，可见结果为清末改革和对外危机将在这一权力格局下发展。这一节点虽可放入改革、外交或革命的宏观年表，却不能脱离省份财政、地方官、商会、学校、军队、家庭与交通网络来理解；同一上谕、条约或战事在不同地域的实际后果并不一致。

核心材料为《清德宗实录》光绪元年相关条；宫中朱批奏折、内阁大库题本与《清史稿》相关志传。它们可以支持年序、制度文本、机构设立或某些地点的行动，却不能自动代表全体居民的收入、教育、身份、性别关系或政治意愿。账本提示“以光绪元年（1875年）的同期文书为定位起点；官修记录、奏报或外交文本服务特定机构立场，具体日期、规模、地方执行与对手视角须以独立材料复核。”。因此，官方统计、外交措辞、报刊论战、教会叙述和军报数字必须分别置于生成它们的立场与制度目的之中，不能互相替代。

在地方尺度上，这一事件应被看作资源和信息重新配置的过程：经费能否落实，取决于税源和中介；学校、警政或新军能否运行，取决于人员、土地和日常协作；家庭与社群对改革的经历，则可能因性别、职业、族属、居处与迁徙而不同。现存材料只照亮少数路径，叙事须保留被登记者与未被登记者之间的差距。

账本条目\texttt{EQ-1875-GUANGXU-ACCESSION-DOCUMENTS}记录，光绪元年（1875年），清发生“围绕「光绪帝即位，慈禧、慈安继续掌握关键政治影响」的上谕、奏折与部议在中央—地方行政链中流转”。账本将其背景说明为核心事件「光绪帝即位，慈禧、慈安继续掌握关键政治影响」要求中央或地方通过文书处理资源、信息或责任。，可见结果为该记录为后续按具体谕旨、奏折、条约、报刊或地方档案拆分事件提供可回查入口。。这一节点虽可放入改革、外交或革命的宏观年表，却不能脱离省份财政、地方官、商会、学校、军队、家庭与交通网络来理解；同一上谕、条约或战事在不同地域的实际后果并不一致。

核心材料为《清德宗实录》光绪元年相关条；宫中朱批奏折、内阁大库题本与《清史稿》相关志传。它们可以支持年序、制度文本、机构设立或某些地点的行动，却不能自动代表全体居民的收入、教育、身份、性别关系或政治意愿。账本提示“此行仅确认相关文书链和可追查的行政动作；不得由其直接推断政策有效、地方服从、民众态度或单一因果。”。因此，官方统计、外交措辞、报刊论战、教会叙述和军报数字必须分别置于生成它们的立场与制度目的之中，不能互相替代。

在地方尺度上，这一事件应被看作资源和信息重新配置的过程：经费能否落实，取决于税源和中介；学校、警政或新军能否运行，取决于人员、土地和日常协作；家庭与社群对改革的经历，则可能因性别、职业、族属、居处与迁徙而不同。现存材料只照亮少数路径，叙事须保留被登记者与未被登记者之间的差距。

账本条目\texttt{EQ-1875-YUNNAN-MUSLIM-WAR}记录，光绪元年（1875年），清与云南回民社群发生“杜文秀政权覆亡前后，云南回民战争进入惨烈收束阶段”。账本将其背景说明为地方军政、矿业贸易与族群冲突长期累积，可见结果为死亡、迁徙与复垦规模的估计须保持不确定性。这一节点虽可放入改革、外交或革命的宏观年表，却不能脱离省份财政、地方官、商会、学校、军队、家庭与交通网络来理解；同一上谕、条约或战事在不同地域的实际后果并不一致。

核心材料为《清德宗实录》光绪元年相关条；军机处档、战事方略及地方奏报。它们可以支持年序、制度文本、机构设立或某些地点的行动，却不能自动代表全体居民的收入、教育、身份、性别关系或政治意愿。账本提示“以光绪元年（1875年）的同期文书为定位起点；官修记录、奏报或外交文本服务特定机构立场，具体日期、规模、地方执行与对手视角须以独立材料复核。”。因此，官方统计、外交措辞、报刊论战、教会叙述和军报数字必须分别置于生成它们的立场与制度目的之中，不能互相替代。

在地方尺度上，这一事件应被看作资源和信息重新配置的过程：经费能否落实，取决于税源和中介；学校、警政或新军能否运行，取决于人员、土地和日常协作；家庭与社群对改革的经历，则可能因性别、职业、族属、居处与迁徙而不同。现存材料只照亮少数路径，叙事须保留被登记者与未被登记者之间的差距。

账本条目\texttt{EQ-1875-YUNNAN-MUSLIM-WAR-DOCUMENTS}记录，光绪元年（1875年），清与云南回民社群发生“清廷围绕「杜文秀政权覆亡前后，云南回民战争进入惨烈收束阶段」调遣军队、筹措饷械并处理战报，中央与地方的军政文书形成连续记录”。账本将其背景说明为核心事件「杜文秀政权覆亡前后，云南回民战争进入惨烈收束阶段」要求中央或地方通过文书处理资源、信息或责任。，可见结果为该记录为后续按具体谕旨、奏折、条约、报刊或地方档案拆分事件提供可回查入口。。这一节点虽可放入改革、外交或革命的宏观年表，却不能脱离省份财政、地方官、商会、学校、军队、家庭与交通网络来理解；同一上谕、条约或战事在不同地域的实际后果并不一致。

核心材料为《清德宗实录》光绪元年相关条；军机处档、战事方略及地方奏报。它们可以支持年序、制度文本、机构设立或某些地点的行动，却不能自动代表全体居民的收入、教育、身份、性别关系或政治意愿。账本提示“此行仅确认相关文书链和可追查的行政动作；不得由其直接推断政策有效、地方服从、民众态度或单一因果。”。因此，官方统计、外交措辞、报刊论战、教会叙述和军报数字必须分别置于生成它们的立场与制度目的之中，不能互相替代。

在地方尺度上，这一事件应被看作资源和信息重新配置的过程：经费能否落实，取决于税源和中介；学校、警政或新军能否运行，取决于人员、土地和日常协作；家庭与社群对改革的经历，则可能因性别、职业、族属、居处与迁徙而不同。现存材料只照亮少数路径，叙事须保留被登记者与未被登记者之间的差距。

账本条目\texttt{EQ-1876-CHEFOO-CONVENTION}记录，光绪二年（1876年），清与英国发生“《烟台条约》签订，扩大英国在华通商和内地活动权利”。账本将其背景说明为马嘉理事件与英国外交压力推动谈判，可见结果为条约权利的地区落实和地方反应需要逐案追踪。这一节点虽可放入改革、外交或革命的宏观年表，却不能脱离省份财政、地方官、商会、学校、军队、家庭与交通网络来理解；同一上谕、条约或战事在不同地域的实际后果并不一致。

核心材料为《清德宗实录》光绪二年相关条；《筹办夷务始末》相关年卷与中外照会、条约文本。它们可以支持年序、制度文本、机构设立或某些地点的行动，却不能自动代表全体居民的收入、教育、身份、性别关系或政治意愿。账本提示“以光绪二年（1876年）的同期文书为定位起点；官修记录、奏报或外交文本服务特定机构立场，具体日期、规模、地方执行与对手视角须以独立材料复核。”。因此，官方统计、外交措辞、报刊论战、教会叙述和军报数字必须分别置于生成它们的立场与制度目的之中，不能互相替代。

在地方尺度上，这一事件应被看作资源和信息重新配置的过程：经费能否落实，取决于税源和中介；学校、警政或新军能否运行，取决于人员、土地和日常协作；家庭与社群对改革的经历，则可能因性别、职业、族属、居处与迁徙而不同。现存材料只照亮少数路径，叙事须保留被登记者与未被登记者之间的差距。

账本条目\texttt{EQ-1876-CHEFOO-CONVENTION-DOCUMENTS}记录，光绪二年（1876年），清与英国发生“围绕「《烟台条约》签订，扩大英国在华通商和内地活动权利」的照会、条约文本、换约或地方执行报告分别进入外交文书链”。账本将其背景说明为核心事件「《烟台条约》签订，扩大英国在华通商和内地活动权利」要求中央或地方通过文书处理资源、信息或责任。，可见结果为该记录为后续按具体谕旨、奏折、条约、报刊或地方档案拆分事件提供可回查入口。。这一节点虽可放入改革、外交或革命的宏观年表，却不能脱离省份财政、地方官、商会、学校、军队、家庭与交通网络来理解；同一上谕、条约或战事在不同地域的实际后果并不一致。

核心材料为《清德宗实录》光绪二年相关条；《筹办夷务始末》相关年卷与中外照会、条约文本。它们可以支持年序、制度文本、机构设立或某些地点的行动，却不能自动代表全体居民的收入、教育、身份、性别关系或政治意愿。账本提示“此行仅确认相关文书链和可追查的行政动作；不得由其直接推断政策有效、地方服从、民众态度或单一因果。”。因此，官方统计、外交措辞、报刊论战、教会叙述和军报数字必须分别置于生成它们的立场与制度目的之中，不能互相替代。

在地方尺度上，这一事件应被看作资源和信息重新配置的过程：经费能否落实，取决于税源和中介；学校、警政或新军能否运行，取决于人员、土地和日常协作；家庭与社群对改革的经历，则可能因性别、职业、族属、居处与迁徙而不同。现存材料只照亮少数路径，叙事须保留被登记者与未被登记者之间的差距。

账本条目\texttt{EQ-1876-NORTH-CHINA-FAMINE}记录，光绪二年（1876年），清与华北发生“丁戊奇荒在华北及西北造成广泛饥荒、迁徙和救济行动”。账本将其背景说明为旱灾、粮价、交通和行政能力共同塑造灾情，可见结果为伤亡总数争议极大，应按地区、材料与估算口径说明。这一节点虽可放入改革、外交或革命的宏观年表，却不能脱离省份财政、地方官、商会、学校、军队、家庭与交通网络来理解；同一上谕、条约或战事在不同地域的实际后果并不一致。

核心材料为《清德宗实录》光绪二年相关条；地方志、督抚奏折及地方档案。它们可以支持年序、制度文本、机构设立或某些地点的行动，却不能自动代表全体居民的收入、教育、身份、性别关系或政治意愿。账本提示“以光绪二年（1876年）的同期文书为定位起点；官修记录、奏报或外交文本服务特定机构立场，具体日期、规模、地方执行与对手视角须以独立材料复核。”。因此，官方统计、外交措辞、报刊论战、教会叙述和军报数字必须分别置于生成它们的立场与制度目的之中，不能互相替代。

在地方尺度上，这一事件应被看作资源和信息重新配置的过程：经费能否落实，取决于税源和中介；学校、警政或新军能否运行，取决于人员、土地和日常协作；家庭与社群对改革的经历，则可能因性别、职业、族属、居处与迁徙而不同。现存材料只照亮少数路径，叙事须保留被登记者与未被登记者之间的差距。

账本条目\texttt{EQ-1876-NORTH-CHINA-FAMINE-DOCUMENTS}记录，光绪二年（1876年），清与华北发生“围绕「丁戊奇荒在华北及西北造成广泛饥荒、迁徙和救济行动」的地方呈报、案件或赈务记录将社会反应带入官府文书链”。账本将其背景说明为核心事件「丁戊奇荒在华北及西北造成广泛饥荒、迁徙和救济行动」要求中央或地方通过文书处理资源、信息或责任。，可见结果为该记录为后续按具体谕旨、奏折、条约、报刊或地方档案拆分事件提供可回查入口。。这一节点虽可放入改革、外交或革命的宏观年表，却不能脱离省份财政、地方官、商会、学校、军队、家庭与交通网络来理解；同一上谕、条约或战事在不同地域的实际后果并不一致。

核心材料为《清德宗实录》光绪二年相关条；地方志、督抚奏折及地方档案。它们可以支持年序、制度文本、机构设立或某些地点的行动，却不能自动代表全体居民的收入、教育、身份、性别关系或政治意愿。账本提示“此行仅确认相关文书链和可追查的行政动作；不得由其直接推断政策有效、地方服从、民众态度或单一因果。”。因此，官方统计、外交措辞、报刊论战、教会叙述和军报数字必须分别置于生成它们的立场与制度目的之中，不能互相替代。

在地方尺度上，这一事件应被看作资源和信息重新配置的过程：经费能否落实，取决于税源和中介；学校、警政或新军能否运行，取决于人员、土地和日常协作；家庭与社群对改革的经历，则可能因性别、职业、族属、居处与迁徙而不同。现存材料只照亮少数路径，叙事须保留被登记者与未被登记者之间的差距。

账本条目\texttt{EQ-1877-FAMINE-RELIEF-NETWORKS}记录，光绪三年（1877年），清、士绅与国际救济者发生“赈济网络在灾区展开，官办、绅办和传教团体的救助并存”。账本将其背景说明为灾害规模超出常规仓储和地方财政能力，可见结果为救助记录反映组织覆盖，不等于所有灾民获救。这一节点虽可放入改革、外交或革命的宏观年表，却不能脱离省份财政、地方官、商会、学校、军队、家庭与交通网络来理解；同一上谕、条约或战事在不同地域的实际后果并不一致。

核心材料为《清德宗实录》光绪三年相关条；地方志、督抚奏折及地方档案。它们可以支持年序、制度文本、机构设立或某些地点的行动，却不能自动代表全体居民的收入、教育、身份、性别关系或政治意愿。账本提示“以光绪三年（1877年）的同期文书为定位起点；官修记录、奏报或外交文本服务特定机构立场，具体日期、规模、地方执行与对手视角须以独立材料复核。”。因此，官方统计、外交措辞、报刊论战、教会叙述和军报数字必须分别置于生成它们的立场与制度目的之中，不能互相替代。

在地方尺度上，这一事件应被看作资源和信息重新配置的过程：经费能否落实，取决于税源和中介；学校、警政或新军能否运行，取决于人员、土地和日常协作；家庭与社群对改革的经历，则可能因性别、职业、族属、居处与迁徙而不同。现存材料只照亮少数路径，叙事须保留被登记者与未被登记者之间的差距。

账本条目\texttt{EQ-1877-FAMINE-RELIEF-NETWORKS-DOCUMENTS}记录，光绪三年（1877年），清、士绅与国际救济者发生“围绕「赈济网络在灾区展开，官办、绅办和传教团体的救助并存」的地方呈报、案件或赈务记录将社会反应带入官府文书链”。账本将其背景说明为核心事件「赈济网络在灾区展开，官办、绅办和传教团体的救助并存」要求中央或地方通过文书处理资源、信息或责任。，可见结果为该记录为后续按具体谕旨、奏折、条约、报刊或地方档案拆分事件提供可回查入口。。这一节点虽可放入改革、外交或革命的宏观年表，却不能脱离省份财政、地方官、商会、学校、军队、家庭与交通网络来理解；同一上谕、条约或战事在不同地域的实际后果并不一致。

核心材料为《清德宗实录》光绪三年相关条；地方志、督抚奏折及地方档案。它们可以支持年序、制度文本、机构设立或某些地点的行动，却不能自动代表全体居民的收入、教育、身份、性别关系或政治意愿。账本提示“此行仅确认相关文书链和可追查的行政动作；不得由其直接推断政策有效、地方服从、民众态度或单一因果。”。因此，官方统计、外交措辞、报刊论战、教会叙述和军报数字必须分别置于生成它们的立场与制度目的之中，不能互相替代。

在地方尺度上，这一事件应被看作资源和信息重新配置的过程：经费能否落实，取决于税源和中介；学校、警政或新军能否运行，取决于人员、土地和日常协作；家庭与社群对改革的经历，则可能因性别、职业、族属、居处与迁徙而不同。现存材料只照亮少数路径，叙事须保留被登记者与未被登记者之间的差距。

账本条目\texttt{EQ-1878-Zuo-ZONGTANG-XINJIANG}记录，光绪四年（1878年），清与阿古柏政权发生“左宗棠军收复新疆大部，阿古柏政权崩解”。账本将其背景说明为清廷集中湘军、饷源和西北运输资源进行反攻，可见结果为军事收复与后续行政、绿洲社会重建必须分开叙述。这一节点虽可放入改革、外交或革命的宏观年表，却不能脱离省份财政、地方官、商会、学校、军队、家庭与交通网络来理解；同一上谕、条约或战事在不同地域的实际后果并不一致。

核心材料为《清德宗实录》光绪四年相关条；军机处档、战事方略及地方奏报。它们可以支持年序、制度文本、机构设立或某些地点的行动，却不能自动代表全体居民的收入、教育、身份、性别关系或政治意愿。账本提示“以光绪四年（1878年）的同期文书为定位起点；官修记录、奏报或外交文本服务特定机构立场，具体日期、规模、地方执行与对手视角须以独立材料复核。”。因此，官方统计、外交措辞、报刊论战、教会叙述和军报数字必须分别置于生成它们的立场与制度目的之中，不能互相替代。

在地方尺度上，这一事件应被看作资源和信息重新配置的过程：经费能否落实，取决于税源和中介；学校、警政或新军能否运行，取决于人员、土地和日常协作；家庭与社群对改革的经历，则可能因性别、职业、族属、居处与迁徙而不同。现存材料只照亮少数路径，叙事须保留被登记者与未被登记者之间的差距。

账本条目\texttt{EQ-1878-Zuo-ZONGTANG-XINJIANG-DOCUMENTS}记录，光绪四年（1878年），清与阿古柏政权发生“清廷围绕「左宗棠军收复新疆大部，阿古柏政权崩解」调遣军队、筹措饷械并处理战报，中央与地方的军政文书形成连续记录”。账本将其背景说明为核心事件「左宗棠军收复新疆大部，阿古柏政权崩解」要求中央或地方通过文书处理资源、信息或责任。，可见结果为该记录为后续按具体谕旨、奏折、条约、报刊或地方档案拆分事件提供可回查入口。。这一节点虽可放入改革、外交或革命的宏观年表，却不能脱离省份财政、地方官、商会、学校、军队、家庭与交通网络来理解；同一上谕、条约或战事在不同地域的实际后果并不一致。

核心材料为《清德宗实录》光绪四年相关条；军机处档、战事方略及地方奏报。它们可以支持年序、制度文本、机构设立或某些地点的行动，却不能自动代表全体居民的收入、教育、身份、性别关系或政治意愿。账本提示“此行仅确认相关文书链和可追查的行政动作；不得由其直接推断政策有效、地方服从、民众态度或单一因果。”。因此，官方统计、外交措辞、报刊论战、教会叙述和军报数字必须分别置于生成它们的立场与制度目的之中，不能互相替代。

在地方尺度上，这一事件应被看作资源和信息重新配置的过程：经费能否落实，取决于税源和中介；学校、警政或新军能否运行，取决于人员、土地和日常协作；家庭与社群对改革的经历，则可能因性别、职业、族属、居处与迁徙而不同。现存材料只照亮少数路径，叙事须保留被登记者与未被登记者之间的差距。

账本条目\texttt{EQ-1878-XINJIANG-PROVINCE-DEBATE}记录，光绪四年（1878年），清发生“朝廷讨论新疆设省、军政与财政安排”。账本将其背景说明为战后长期驻防成本和俄国边境压力要求重建制度，可见结果为设省方案反映中央目标，地方实施尚需持续资源。这一节点虽可放入改革、外交或革命的宏观年表，却不能脱离省份财政、地方官、商会、学校、军队、家庭与交通网络来理解；同一上谕、条约或战事在不同地域的实际后果并不一致。

核心材料为《清德宗实录》光绪四年相关条；宫中朱批奏折、内阁大库题本与《清史稿》相关志传。它们可以支持年序、制度文本、机构设立或某些地点的行动，却不能自动代表全体居民的收入、教育、身份、性别关系或政治意愿。账本提示“以光绪四年（1878年）的同期文书为定位起点；官修记录、奏报或外交文本服务特定机构立场，具体日期、规模、地方执行与对手视角须以独立材料复核。”。因此，官方统计、外交措辞、报刊论战、教会叙述和军报数字必须分别置于生成它们的立场与制度目的之中，不能互相替代。

在地方尺度上，这一事件应被看作资源和信息重新配置的过程：经费能否落实，取决于税源和中介；学校、警政或新军能否运行，取决于人员、土地和日常协作；家庭与社群对改革的经历，则可能因性别、职业、族属、居处与迁徙而不同。现存材料只照亮少数路径，叙事须保留被登记者与未被登记者之间的差距。

账本条目\texttt{EQ-1878-XINJIANG-PROVINCE-DEBATE-DOCUMENTS}记录，光绪四年（1878年），清发生“围绕「朝廷讨论新疆设省、军政与财政安排」的上谕、奏折与部议在中央—地方行政链中流转”。账本将其背景说明为核心事件「朝廷讨论新疆设省、军政与财政安排」要求中央或地方通过文书处理资源、信息或责任。，可见结果为该记录为后续按具体谕旨、奏折、条约、报刊或地方档案拆分事件提供可回查入口。。这一节点虽可放入改革、外交或革命的宏观年表，却不能脱离省份财政、地方官、商会、学校、军队、家庭与交通网络来理解；同一上谕、条约或战事在不同地域的实际后果并不一致。

核心材料为《清德宗实录》光绪四年相关条；宫中朱批奏折、内阁大库题本与《清史稿》相关志传。它们可以支持年序、制度文本、机构设立或某些地点的行动，却不能自动代表全体居民的收入、教育、身份、性别关系或政治意愿。账本提示“此行仅确认相关文书链和可追查的行政动作；不得由其直接推断政策有效、地方服从、民众态度或单一因果。”。因此，官方统计、外交措辞、报刊论战、教会叙述和军报数字必须分别置于生成它们的立场与制度目的之中，不能互相替代。

在地方尺度上，这一事件应被看作资源和信息重新配置的过程：经费能否落实，取决于税源和中介；学校、警政或新军能否运行，取决于人员、土地和日常协作；家庭与社群对改革的经历，则可能因性别、职业、族属、居处与迁徙而不同。现存材料只照亮少数路径，叙事须保留被登记者与未被登记者之间的差距。

账本条目\texttt{EQ-1879-RYUKYU-ANNEXATION}记录，光绪五年（1879年），日本、琉球与清发生“日本吞并琉球，清廷的交涉未能恢复旧有册封关系”。账本将其背景说明为日本国家建构与东亚秩序转变相互作用，可见结果为清方外交文书与琉球、日本材料对权利表述不同。这一节点虽可放入改革、外交或革命的宏观年表，却不能脱离省份财政、地方官、商会、学校、军队、家庭与交通网络来理解；同一上谕、条约或战事在不同地域的实际后果并不一致。

核心材料为《清德宗实录》光绪五年相关条；《筹办夷务始末》相关年卷与中外照会、条约文本。它们可以支持年序、制度文本、机构设立或某些地点的行动，却不能自动代表全体居民的收入、教育、身份、性别关系或政治意愿。账本提示“以光绪五年（1879年）的同期文书为定位起点；官修记录、奏报或外交文本服务特定机构立场，具体日期、规模、地方执行与对手视角须以独立材料复核。”。因此，官方统计、外交措辞、报刊论战、教会叙述和军报数字必须分别置于生成它们的立场与制度目的之中，不能互相替代。

在地方尺度上，这一事件应被看作资源和信息重新配置的过程：经费能否落实，取决于税源和中介；学校、警政或新军能否运行，取决于人员、土地和日常协作；家庭与社群对改革的经历，则可能因性别、职业、族属、居处与迁徙而不同。现存材料只照亮少数路径，叙事须保留被登记者与未被登记者之间的差距。

账本条目\texttt{EQ-1879-RYUKYU-ANNEXATION-DOCUMENTS}记录，光绪五年（1879年），日本、琉球与清发生“围绕「日本吞并琉球，清廷的交涉未能恢复旧有册封关系」的照会、条约文本、换约或地方执行报告分别进入外交文书链”。账本将其背景说明为核心事件「日本吞并琉球，清廷的交涉未能恢复旧有册封关系」要求中央或地方通过文书处理资源、信息或责任。，可见结果为该记录为后续按具体谕旨、奏折、条约、报刊或地方档案拆分事件提供可回查入口。。这一节点虽可放入改革、外交或革命的宏观年表，却不能脱离省份财政、地方官、商会、学校、军队、家庭与交通网络来理解；同一上谕、条约或战事在不同地域的实际后果并不一致。

核心材料为《清德宗实录》光绪五年相关条；《筹办夷务始末》相关年卷与中外照会、条约文本。它们可以支持年序、制度文本、机构设立或某些地点的行动，却不能自动代表全体居民的收入、教育、身份、性别关系或政治意愿。账本提示“此行仅确认相关文书链和可追查的行政动作；不得由其直接推断政策有效、地方服从、民众态度或单一因果。”。因此，官方统计、外交措辞、报刊论战、教会叙述和军报数字必须分别置于生成它们的立场与制度目的之中，不能互相替代。

在地方尺度上，这一事件应被看作资源和信息重新配置的过程：经费能否落实，取决于税源和中介；学校、警政或新军能否运行，取决于人员、土地和日常协作；家庭与社群对改革的经历，则可能因性别、职业、族属、居处与迁徙而不同。现存材料只照亮少数路径，叙事须保留被登记者与未被登记者之间的差距。

账本条目\texttt{EQ-1880-ILI-NEGOTIATIONS}记录，光绪六年（1880年），清与俄罗斯发生“清俄围绕伊犁归还和边界条件持续谈判”。账本将其背景说明为俄军占据伊犁后，清廷在军事与财政限制下寻求外交解决，可见结果为边界线、赔款和贸易条款须逐项区分。这一节点虽可放入改革、外交或革命的宏观年表，却不能脱离省份财政、地方官、商会、学校、军队、家庭与交通网络来理解；同一上谕、条约或战事在不同地域的实际后果并不一致。

核心材料为《清德宗实录》光绪六年相关条；《筹办夷务始末》相关年卷与中外照会、条约文本。它们可以支持年序、制度文本、机构设立或某些地点的行动，却不能自动代表全体居民的收入、教育、身份、性别关系或政治意愿。账本提示“以光绪六年（1880年）的同期文书为定位起点；官修记录、奏报或外交文本服务特定机构立场，具体日期、规模、地方执行与对手视角须以独立材料复核。”。因此，官方统计、外交措辞、报刊论战、教会叙述和军报数字必须分别置于生成它们的立场与制度目的之中，不能互相替代。

在地方尺度上，这一事件应被看作资源和信息重新配置的过程：经费能否落实，取决于税源和中介；学校、警政或新军能否运行，取决于人员、土地和日常协作；家庭与社群对改革的经历，则可能因性别、职业、族属、居处与迁徙而不同。现存材料只照亮少数路径，叙事须保留被登记者与未被登记者之间的差距。

账本条目\texttt{EQ-1880-ILI-NEGOTIATIONS-DOCUMENTS}记录，光绪六年（1880年），清与俄罗斯发生“围绕「清俄围绕伊犁归还和边界条件持续谈判」的照会、条约文本、换约或地方执行报告分别进入外交文书链”。账本将其背景说明为核心事件「清俄围绕伊犁归还和边界条件持续谈判」要求中央或地方通过文书处理资源、信息或责任。，可见结果为该记录为后续按具体谕旨、奏折、条约、报刊或地方档案拆分事件提供可回查入口。。这一节点虽可放入改革、外交或革命的宏观年表，却不能脱离省份财政、地方官、商会、学校、军队、家庭与交通网络来理解；同一上谕、条约或战事在不同地域的实际后果并不一致。

核心材料为《清德宗实录》光绪六年相关条；《筹办夷务始末》相关年卷与中外照会、条约文本。它们可以支持年序、制度文本、机构设立或某些地点的行动，却不能自动代表全体居民的收入、教育、身份、性别关系或政治意愿。账本提示“此行仅确认相关文书链和可追查的行政动作；不得由其直接推断政策有效、地方服从、民众态度或单一因果。”。因此，官方统计、外交措辞、报刊论战、教会叙述和军报数字必须分别置于生成它们的立场与制度目的之中，不能互相替代。

在地方尺度上，这一事件应被看作资源和信息重新配置的过程：经费能否落实，取决于税源和中介；学校、警政或新军能否运行，取决于人员、土地和日常协作；家庭与社群对改革的经历，则可能因性别、职业、族属、居处与迁徙而不同。现存材料只照亮少数路径，叙事须保留被登记者与未被登记者之间的差距。

账本条目\texttt{EQ-1881-SAINT-PETERSBURG-TREATY}记录，光绪七年（1881年），清与俄罗斯发生“《中俄伊犁条约》签订，伊犁大部归还清廷并附带领土、赔款和贸易条件”。账本将其背景说明为长期谈判与西北军事实力恢复共同作用，可见结果为条约地图与实际边界管理不是同一问题。这一节点虽可放入改革、外交或革命的宏观年表，却不能脱离省份财政、地方官、商会、学校、军队、家庭与交通网络来理解；同一上谕、条约或战事在不同地域的实际后果并不一致。

核心材料为《清德宗实录》光绪七年相关条；《筹办夷务始末》相关年卷与中外照会、条约文本。它们可以支持年序、制度文本、机构设立或某些地点的行动，却不能自动代表全体居民的收入、教育、身份、性别关系或政治意愿。账本提示“以光绪七年（1881年）的同期文书为定位起点；官修记录、奏报或外交文本服务特定机构立场，具体日期、规模、地方执行与对手视角须以独立材料复核。”。因此，官方统计、外交措辞、报刊论战、教会叙述和军报数字必须分别置于生成它们的立场与制度目的之中，不能互相替代。

在地方尺度上，这一事件应被看作资源和信息重新配置的过程：经费能否落实，取决于税源和中介；学校、警政或新军能否运行，取决于人员、土地和日常协作；家庭与社群对改革的经历，则可能因性别、职业、族属、居处与迁徙而不同。现存材料只照亮少数路径，叙事须保留被登记者与未被登记者之间的差距。

账本条目\texttt{EQ-1881-SAINT-PETERSBURG-TREATY-DOCUMENTS}记录，光绪七年（1881年），清与俄罗斯发生“围绕「《中俄伊犁条约》签订，伊犁大部归还清廷并附带领土、赔款和贸易条件」的照会、条约文本、换约或地方执行报告分别进入外交文书链”。账本将其背景说明为核心事件「《中俄伊犁条约》签订，伊犁大部归还清廷并附带领土、赔款和贸易条件」要求中央或地方通过文书处理资源、信息或责任。，可见结果为该记录为后续按具体谕旨、奏折、条约、报刊或地方档案拆分事件提供可回查入口。。这一节点虽可放入改革、外交或革命的宏观年表，却不能脱离省份财政、地方官、商会、学校、军队、家庭与交通网络来理解；同一上谕、条约或战事在不同地域的实际后果并不一致。

核心材料为《清德宗实录》光绪七年相关条；《筹办夷务始末》相关年卷与中外照会、条约文本。它们可以支持年序、制度文本、机构设立或某些地点的行动，却不能自动代表全体居民的收入、教育、身份、性别关系或政治意愿。账本提示“此行仅确认相关文书链和可追查的行政动作；不得由其直接推断政策有效、地方服从、民众态度或单一因果。”。因此，官方统计、外交措辞、报刊论战、教会叙述和军报数字必须分别置于生成它们的立场与制度目的之中，不能互相替代。

在地方尺度上，这一事件应被看作资源和信息重新配置的过程：经费能否落实，取决于税源和中介；学校、警政或新军能否运行，取决于人员、土地和日常协作；家庭与社群对改革的经历，则可能因性别、职业、族属、居处与迁徙而不同。现存材料只照亮少数路径，叙事须保留被登记者与未被登记者之间的差距。

账本条目\texttt{EQ-1882-IMO-INCIDENT}记录，光绪八年（1882年），清、朝鲜与日本发生“朝鲜壬午兵变后清军入朝，日本也扩大影响”。账本将其背景说明为朝鲜内政、军制改革与列强竞争交织，可见结果为清、朝、日文献中的“保护”与控制含义不同。这一节点虽可放入改革、外交或革命的宏观年表，却不能脱离省份财政、地方官、商会、学校、军队、家庭与交通网络来理解；同一上谕、条约或战事在不同地域的实际后果并不一致。

核心材料为《清德宗实录》光绪八年相关条；《筹办夷务始末》相关年卷与中外照会、条约文本。它们可以支持年序、制度文本、机构设立或某些地点的行动，却不能自动代表全体居民的收入、教育、身份、性别关系或政治意愿。账本提示“以光绪八年（1882年）的同期文书为定位起点；官修记录、奏报或外交文本服务特定机构立场，具体日期、规模、地方执行与对手视角须以独立材料复核。”。因此，官方统计、外交措辞、报刊论战、教会叙述和军报数字必须分别置于生成它们的立场与制度目的之中，不能互相替代。

在地方尺度上，这一事件应被看作资源和信息重新配置的过程：经费能否落实，取决于税源和中介；学校、警政或新军能否运行，取决于人员、土地和日常协作；家庭与社群对改革的经历，则可能因性别、职业、族属、居处与迁徙而不同。现存材料只照亮少数路径，叙事须保留被登记者与未被登记者之间的差距。

账本条目\texttt{EQ-1882-IMO-INCIDENT-DOCUMENTS}记录，光绪八年（1882年），清、朝鲜与日本发生“围绕「朝鲜壬午兵变后清军入朝，日本也扩大影响」的照会、条约文本、换约或地方执行报告分别进入外交文书链”。账本将其背景说明为核心事件「朝鲜壬午兵变后清军入朝，日本也扩大影响」要求中央或地方通过文书处理资源、信息或责任。，可见结果为该记录为后续按具体谕旨、奏折、条约、报刊或地方档案拆分事件提供可回查入口。。这一节点虽可放入改革、外交或革命的宏观年表，却不能脱离省份财政、地方官、商会、学校、军队、家庭与交通网络来理解；同一上谕、条约或战事在不同地域的实际后果并不一致。

核心材料为《清德宗实录》光绪八年相关条；《筹办夷务始末》相关年卷与中外照会、条约文本。它们可以支持年序、制度文本、机构设立或某些地点的行动，却不能自动代表全体居民的收入、教育、身份、性别关系或政治意愿。账本提示“此行仅确认相关文书链和可追查的行政动作；不得由其直接推断政策有效、地方服从、民众态度或单一因果。”。因此，官方统计、外交措辞、报刊论战、教会叙述和军报数字必须分别置于生成它们的立场与制度目的之中，不能互相替代。

在地方尺度上，这一事件应被看作资源和信息重新配置的过程：经费能否落实，取决于税源和中介；学校、警政或新军能否运行，取决于人员、土地和日常协作；家庭与社群对改革的经历，则可能因性别、职业、族属、居处与迁徙而不同。现存材料只照亮少数路径，叙事须保留被登记者与未被登记者之间的差距。

账本条目\texttt{EQ-1883-SINO-FRENCH-WAR-BEGIN}记录，光绪九年（1883年），清、法国与越南发生“中法因越南北部和宗藩关系冲突升级为战争”。账本将其背景说明为法国殖民扩张与清廷在越南的传统关系碰撞，可见结果为战事跨越越南、台湾和华南，不能只从北京外交看。这一节点虽可放入改革、外交或革命的宏观年表，却不能脱离省份财政、地方官、商会、学校、军队、家庭与交通网络来理解；同一上谕、条约或战事在不同地域的实际后果并不一致。

核心材料为《清德宗实录》光绪九年相关条；军机处档、战事方略及地方奏报。它们可以支持年序、制度文本、机构设立或某些地点的行动，却不能自动代表全体居民的收入、教育、身份、性别关系或政治意愿。账本提示“以光绪九年（1883年）的同期文书为定位起点；官修记录、奏报或外交文本服务特定机构立场，具体日期、规模、地方执行与对手视角须以独立材料复核。”。因此，官方统计、外交措辞、报刊论战、教会叙述和军报数字必须分别置于生成它们的立场与制度目的之中，不能互相替代。

在地方尺度上，这一事件应被看作资源和信息重新配置的过程：经费能否落实，取决于税源和中介；学校、警政或新军能否运行，取决于人员、土地和日常协作；家庭与社群对改革的经历，则可能因性别、职业、族属、居处与迁徙而不同。现存材料只照亮少数路径，叙事须保留被登记者与未被登记者之间的差距。

账本条目\texttt{EQ-1883-SINO-FRENCH-WAR-BEGIN-DOCUMENTS}记录，光绪九年（1883年），清、法国与越南发生“清廷围绕「中法因越南北部和宗藩关系冲突升级为战争」调遣军队、筹措饷械并处理战报，中央与地方的军政文书形成连续记录”。账本将其背景说明为核心事件「中法因越南北部和宗藩关系冲突升级为战争」要求中央或地方通过文书处理资源、信息或责任。，可见结果为该记录为后续按具体谕旨、奏折、条约、报刊或地方档案拆分事件提供可回查入口。。这一节点虽可放入改革、外交或革命的宏观年表，却不能脱离省份财政、地方官、商会、学校、军队、家庭与交通网络来理解；同一上谕、条约或战事在不同地域的实际后果并不一致。

核心材料为《清德宗实录》光绪九年相关条；军机处档、战事方略及地方奏报。它们可以支持年序、制度文本、机构设立或某些地点的行动，却不能自动代表全体居民的收入、教育、身份、性别关系或政治意愿。账本提示“此行仅确认相关文书链和可追查的行政动作；不得由其直接推断政策有效、地方服从、民众态度或单一因果。”。因此，官方统计、外交措辞、报刊论战、教会叙述和军报数字必须分别置于生成它们的立场与制度目的之中，不能互相替代。

在地方尺度上，这一事件应被看作资源和信息重新配置的过程：经费能否落实，取决于税源和中介；学校、警政或新军能否运行，取决于人员、土地和日常协作；家庭与社群对改革的经历，则可能因性别、职业、族属、居处与迁徙而不同。现存材料只照亮少数路径，叙事须保留被登记者与未被登记者之间的差距。

账本条目\texttt{EQ-1884-FUZHOU-NAVAL-BATTLE}记录，光绪十年（1884年），清与法国发生“马江海战中福建水师遭法国舰队重创”。账本将其背景说明为清军舰队训练、指挥与武器体系面临现代海战挑战，可见结果为伤亡和舰只损失应依据中法记录和船政资料核验。这一节点虽可放入改革、外交或革命的宏观年表，却不能脱离省份财政、地方官、商会、学校、军队、家庭与交通网络来理解；同一上谕、条约或战事在不同地域的实际后果并不一致。

核心材料为《清德宗实录》光绪十年相关条；军机处档、战事方略及地方奏报。它们可以支持年序、制度文本、机构设立或某些地点的行动，却不能自动代表全体居民的收入、教育、身份、性别关系或政治意愿。账本提示“以光绪十年（1884年）的同期文书为定位起点；官修记录、奏报或外交文本服务特定机构立场，具体日期、规模、地方执行与对手视角须以独立材料复核。”。因此，官方统计、外交措辞、报刊论战、教会叙述和军报数字必须分别置于生成它们的立场与制度目的之中，不能互相替代。

在地方尺度上，这一事件应被看作资源和信息重新配置的过程：经费能否落实，取决于税源和中介；学校、警政或新军能否运行，取决于人员、土地和日常协作；家庭与社群对改革的经历，则可能因性别、职业、族属、居处与迁徙而不同。现存材料只照亮少数路径，叙事须保留被登记者与未被登记者之间的差距。

账本条目\texttt{EQ-1884-FUZHOU-NAVAL-BATTLE-DOCUMENTS}记录，光绪十年（1884年），清与法国发生“清廷围绕「马江海战中福建水师遭法国舰队重创」调遣军队、筹措饷械并处理战报，中央与地方的军政文书形成连续记录”。账本将其背景说明为核心事件「马江海战中福建水师遭法国舰队重创」要求中央或地方通过文书处理资源、信息或责任。，可见结果为该记录为后续按具体谕旨、奏折、条约、报刊或地方档案拆分事件提供可回查入口。。这一节点虽可放入改革、外交或革命的宏观年表，却不能脱离省份财政、地方官、商会、学校、军队、家庭与交通网络来理解；同一上谕、条约或战事在不同地域的实际后果并不一致。

核心材料为《清德宗实录》光绪十年相关条；军机处档、战事方略及地方奏报。它们可以支持年序、制度文本、机构设立或某些地点的行动，却不能自动代表全体居民的收入、教育、身份、性别关系或政治意愿。账本提示“此行仅确认相关文书链和可追查的行政动作；不得由其直接推断政策有效、地方服从、民众态度或单一因果。”。因此，官方统计、外交措辞、报刊论战、教会叙述和军报数字必须分别置于生成它们的立场与制度目的之中，不能互相替代。

在地方尺度上，这一事件应被看作资源和信息重新配置的过程：经费能否落实，取决于税源和中介；学校、警政或新军能否运行，取决于人员、土地和日常协作；家庭与社群对改革的经历，则可能因性别、职业、族属、居处与迁徙而不同。现存材料只照亮少数路径，叙事须保留被登记者与未被登记者之间的差距。

账本条目\texttt{EQ-1884-TAIWAN-BLOCKADE}记录，光绪十年（1884年），清、法国与台湾发生“法军进攻台湾北部并封锁，台湾成为中法战争重要战场”。账本将其背景说明为台湾港口、煤矿和海峡航路具有战略价值，可见结果为地方抵抗、商业中断与原住民地区情况需分区考察。这一节点虽可放入改革、外交或革命的宏观年表，却不能脱离省份财政、地方官、商会、学校、军队、家庭与交通网络来理解；同一上谕、条约或战事在不同地域的实际后果并不一致。

核心材料为《清德宗实录》光绪十年相关条；军机处档、战事方略及地方奏报。它们可以支持年序、制度文本、机构设立或某些地点的行动，却不能自动代表全体居民的收入、教育、身份、性别关系或政治意愿。账本提示“以光绪十年（1884年）的同期文书为定位起点；官修记录、奏报或外交文本服务特定机构立场，具体日期、规模、地方执行与对手视角须以独立材料复核。”。因此，官方统计、外交措辞、报刊论战、教会叙述和军报数字必须分别置于生成它们的立场与制度目的之中，不能互相替代。

在地方尺度上，这一事件应被看作资源和信息重新配置的过程：经费能否落实，取决于税源和中介；学校、警政或新军能否运行，取决于人员、土地和日常协作；家庭与社群对改革的经历，则可能因性别、职业、族属、居处与迁徙而不同。现存材料只照亮少数路径，叙事须保留被登记者与未被登记者之间的差距。

\subsection{教育、家庭与公共文化}

账本条目\texttt{EQ-1884-TAIWAN-BLOCKADE-DOCUMENTS}记录，光绪十年（1884年），清、法国与台湾发生“清廷围绕「法军进攻台湾北部并封锁，台湾成为中法战争重要战场」调遣军队、筹措饷械并处理战报，中央与地方的军政文书形成连续记录”。账本将其背景说明为核心事件「法军进攻台湾北部并封锁，台湾成为中法战争重要战场」要求中央或地方通过文书处理资源、信息或责任。，可见结果为该记录为后续按具体谕旨、奏折、条约、报刊或地方档案拆分事件提供可回查入口。。这一节点虽可放入改革、外交或革命的宏观年表，却不能脱离省份财政、地方官、商会、学校、军队、家庭与交通网络来理解；同一上谕、条约或战事在不同地域的实际后果并不一致。

核心材料为《清德宗实录》光绪十年相关条；军机处档、战事方略及地方奏报。它们可以支持年序、制度文本、机构设立或某些地点的行动，却不能自动代表全体居民的收入、教育、身份、性别关系或政治意愿。账本提示“此行仅确认相关文书链和可追查的行政动作；不得由其直接推断政策有效、地方服从、民众态度或单一因果。”。因此，官方统计、外交措辞、报刊论战、教会叙述和军报数字必须分别置于生成它们的立场与制度目的之中，不能互相替代。

在地方尺度上，这一事件应被看作资源和信息重新配置的过程：经费能否落实，取决于税源和中介；学校、警政或新军能否运行，取决于人员、土地和日常协作；家庭与社群对改革的经历，则可能因性别、职业、族属、居处与迁徙而不同。现存材料只照亮少数路径，叙事须保留被登记者与未被登记者之间的差距。

账本条目\texttt{EQ-1885-TIANJIN-CHINA-FRANCE-TREATY}记录，光绪十一年（1885年），清与法国发生“《中法新约》确认越南脱离清朝宗藩体系，战争结束”。账本将其背景说明为军事与财政压力推动双方议和，可见结果为条约确认与边界勘界、华南地方影响应分别处理。这一节点虽可放入改革、外交或革命的宏观年表，却不能脱离省份财政、地方官、商会、学校、军队、家庭与交通网络来理解；同一上谕、条约或战事在不同地域的实际后果并不一致。

核心材料为《清德宗实录》光绪十一年相关条；《筹办夷务始末》相关年卷与中外照会、条约文本。它们可以支持年序、制度文本、机构设立或某些地点的行动，却不能自动代表全体居民的收入、教育、身份、性别关系或政治意愿。账本提示“以光绪十一年（1885年）的同期文书为定位起点；官修记录、奏报或外交文本服务特定机构立场，具体日期、规模、地方执行与对手视角须以独立材料复核。”。因此，官方统计、外交措辞、报刊论战、教会叙述和军报数字必须分别置于生成它们的立场与制度目的之中，不能互相替代。

在地方尺度上，这一事件应被看作资源和信息重新配置的过程：经费能否落实，取决于税源和中介；学校、警政或新军能否运行，取决于人员、土地和日常协作；家庭与社群对改革的经历，则可能因性别、职业、族属、居处与迁徙而不同。现存材料只照亮少数路径，叙事须保留被登记者与未被登记者之间的差距。

账本条目\texttt{EQ-1885-TIANJIN-CHINA-FRANCE-TREATY-DOCUMENTS}记录，光绪十一年（1885年），清与法国发生“围绕「《中法新约》确认越南脱离清朝宗藩体系，战争结束」的照会、条约文本、换约或地方执行报告分别进入外交文书链”。账本将其背景说明为核心事件「《中法新约》确认越南脱离清朝宗藩体系，战争结束」要求中央或地方通过文书处理资源、信息或责任。，可见结果为该记录为后续按具体谕旨、奏折、条约、报刊或地方档案拆分事件提供可回查入口。。这一节点虽可放入改革、外交或革命的宏观年表，却不能脱离省份财政、地方官、商会、学校、军队、家庭与交通网络来理解；同一上谕、条约或战事在不同地域的实际后果并不一致。

核心材料为《清德宗实录》光绪十一年相关条；《筹办夷务始末》相关年卷与中外照会、条约文本。它们可以支持年序、制度文本、机构设立或某些地点的行动，却不能自动代表全体居民的收入、教育、身份、性别关系或政治意愿。账本提示“此行仅确认相关文书链和可追查的行政动作；不得由其直接推断政策有效、地方服从、民众态度或单一因果。”。因此，官方统计、外交措辞、报刊论战、教会叙述和军报数字必须分别置于生成它们的立场与制度目的之中，不能互相替代。

在地方尺度上，这一事件应被看作资源和信息重新配置的过程：经费能否落实，取决于税源和中介；学校、警政或新军能否运行，取决于人员、土地和日常协作；家庭与社群对改革的经历，则可能因性别、职业、族属、居处与迁徙而不同。现存材料只照亮少数路径，叙事须保留被登记者与未被登记者之间的差距。

账本条目\texttt{EQ-1885-TAIWAN-PROVINCE}记录，光绪十一年（1885年），清与台湾发生“清廷决定将台湾建省，加强行政、军防和基础设施经营”。账本将其背景说明为中法战争显示台湾海防与行政地位的重要性，可见结果为建省不等于全岛被均质治理，东部和原住民地区尤需谨慎。这一节点虽可放入改革、外交或革命的宏观年表，却不能脱离省份财政、地方官、商会、学校、军队、家庭与交通网络来理解；同一上谕、条约或战事在不同地域的实际后果并不一致。

核心材料为《清德宗实录》光绪十一年相关条；宫中朱批奏折、内阁大库题本与《清史稿》相关志传。它们可以支持年序、制度文本、机构设立或某些地点的行动，却不能自动代表全体居民的收入、教育、身份、性别关系或政治意愿。账本提示“以光绪十一年（1885年）的同期文书为定位起点；官修记录、奏报或外交文本服务特定机构立场，具体日期、规模、地方执行与对手视角须以独立材料复核。”。因此，官方统计、外交措辞、报刊论战、教会叙述和军报数字必须分别置于生成它们的立场与制度目的之中，不能互相替代。

在地方尺度上，这一事件应被看作资源和信息重新配置的过程：经费能否落实，取决于税源和中介；学校、警政或新军能否运行，取决于人员、土地和日常协作；家庭与社群对改革的经历，则可能因性别、职业、族属、居处与迁徙而不同。现存材料只照亮少数路径，叙事须保留被登记者与未被登记者之间的差距。

账本条目\texttt{EQ-1885-TAIWAN-PROVINCE-DOCUMENTS}记录，光绪十一年（1885年），清与台湾发生“围绕「清廷决定将台湾建省，加强行政、军防和基础设施经营」的上谕、奏折与部议在中央—地方行政链中流转”。账本将其背景说明为核心事件「清廷决定将台湾建省，加强行政、军防和基础设施经营」要求中央或地方通过文书处理资源、信息或责任。，可见结果为该记录为后续按具体谕旨、奏折、条约、报刊或地方档案拆分事件提供可回查入口。。这一节点虽可放入改革、外交或革命的宏观年表，却不能脱离省份财政、地方官、商会、学校、军队、家庭与交通网络来理解；同一上谕、条约或战事在不同地域的实际后果并不一致。

核心材料为《清德宗实录》光绪十一年相关条；宫中朱批奏折、内阁大库题本与《清史稿》相关志传。它们可以支持年序、制度文本、机构设立或某些地点的行动，却不能自动代表全体居民的收入、教育、身份、性别关系或政治意愿。账本提示“此行仅确认相关文书链和可追查的行政动作；不得由其直接推断政策有效、地方服从、民众态度或单一因果。”。因此，官方统计、外交措辞、报刊论战、教会叙述和军报数字必须分别置于生成它们的立场与制度目的之中，不能互相替代。

在地方尺度上，这一事件应被看作资源和信息重新配置的过程：经费能否落实，取决于税源和中介；学校、警政或新军能否运行，取决于人员、土地和日常协作；家庭与社群对改革的经历，则可能因性别、职业、族属、居处与迁徙而不同。现存材料只照亮少数路径，叙事须保留被登记者与未被登记者之间的差距。

账本条目\texttt{EQ-1886-BEIYANG-NAVY-EXPANSION}记录，光绪十二年（1886年），清与北洋发生“北洋海军和沿海防务建设持续推进”。账本将其背景说明为中法战争和日本扩张促使海军现代化，可见结果为购舰数量不等于训练、维护、弹药和指挥能力。这一节点虽可放入改革、外交或革命的宏观年表，却不能脱离省份财政、地方官、商会、学校、军队、家庭与交通网络来理解；同一上谕、条约或战事在不同地域的实际后果并不一致。

核心材料为《清德宗实录》光绪十二年相关条；军机处档、战事方略及地方奏报。它们可以支持年序、制度文本、机构设立或某些地点的行动，却不能自动代表全体居民的收入、教育、身份、性别关系或政治意愿。账本提示“以光绪十二年（1886年）的同期文书为定位起点；官修记录、奏报或外交文本服务特定机构立场，具体日期、规模、地方执行与对手视角须以独立材料复核。”。因此，官方统计、外交措辞、报刊论战、教会叙述和军报数字必须分别置于生成它们的立场与制度目的之中，不能互相替代。

在地方尺度上，这一事件应被看作资源和信息重新配置的过程：经费能否落实，取决于税源和中介；学校、警政或新军能否运行，取决于人员、土地和日常协作；家庭与社群对改革的经历，则可能因性别、职业、族属、居处与迁徙而不同。现存材料只照亮少数路径，叙事须保留被登记者与未被登记者之间的差距。

账本条目\texttt{EQ-1886-BEIYANG-NAVY-EXPANSION-DOCUMENTS}记录，光绪十二年（1886年），清与北洋发生“清廷围绕「北洋海军和沿海防务建设持续推进」调遣军队、筹措饷械并处理战报，中央与地方的军政文书形成连续记录”。账本将其背景说明为核心事件「北洋海军和沿海防务建设持续推进」要求中央或地方通过文书处理资源、信息或责任。，可见结果为该记录为后续按具体谕旨、奏折、条约、报刊或地方档案拆分事件提供可回查入口。。这一节点虽可放入改革、外交或革命的宏观年表，却不能脱离省份财政、地方官、商会、学校、军队、家庭与交通网络来理解；同一上谕、条约或战事在不同地域的实际后果并不一致。

核心材料为《清德宗实录》光绪十二年相关条；军机处档、战事方略及地方奏报。它们可以支持年序、制度文本、机构设立或某些地点的行动，却不能自动代表全体居民的收入、教育、身份、性别关系或政治意愿。账本提示“此行仅确认相关文书链和可追查的行政动作；不得由其直接推断政策有效、地方服从、民众态度或单一因果。”。因此，官方统计、外交措辞、报刊论战、教会叙述和军报数字必须分别置于生成它们的立场与制度目的之中，不能互相替代。

在地方尺度上，这一事件应被看作资源和信息重新配置的过程：经费能否落实，取决于税源和中介；学校、警政或新军能否运行，取决于人员、土地和日常协作；家庭与社群对改革的经历，则可能因性别、职业、族属、居处与迁徙而不同。现存材料只照亮少数路径，叙事须保留被登记者与未被登记者之间的差距。

账本条目\texttt{EQ-1887-TAIWAN-RAILWAY-PLANNING}记录，光绪十三年（1887年），清与台湾发生“台湾巡抚刘铭传推动铁路、电报、矿务和新式行政项目”。账本将其背景说明为建省后需要连接港口、城市和军防节点，可见结果为项目经费、路段使用与社会影响需具体核验。这一节点虽可放入改革、外交或革命的宏观年表，却不能脱离省份财政、地方官、商会、学校、军队、家庭与交通网络来理解；同一上谕、条约或战事在不同地域的实际后果并不一致。

核心材料为《清德宗实录》光绪十三年相关条；户部奏销、盐政、河工或海关档案。它们可以支持年序、制度文本、机构设立或某些地点的行动，却不能自动代表全体居民的收入、教育、身份、性别关系或政治意愿。账本提示“以光绪十三年（1887年）的同期文书为定位起点；官修记录、奏报或外交文本服务特定机构立场，具体日期、规模、地方执行与对手视角须以独立材料复核。”。因此，官方统计、外交措辞、报刊论战、教会叙述和军报数字必须分别置于生成它们的立场与制度目的之中，不能互相替代。

在地方尺度上，这一事件应被看作资源和信息重新配置的过程：经费能否落实，取决于税源和中介；学校、警政或新军能否运行，取决于人员、土地和日常协作；家庭与社群对改革的经历，则可能因性别、职业、族属、居处与迁徙而不同。现存材料只照亮少数路径，叙事须保留被登记者与未被登记者之间的差距。

账本条目\texttt{EQ-1887-TAIWAN-RAILWAY-PLANNING-DOCUMENTS}记录，光绪十三年（1887年），清与台湾发生“围绕「台湾巡抚刘铭传推动铁路、电报、矿务和新式行政项目」的经费、税收、赈济或工程呈报经由户部与地方衙门处理”。账本将其背景说明为核心事件「台湾巡抚刘铭传推动铁路、电报、矿务和新式行政项目」要求中央或地方通过文书处理资源、信息或责任。，可见结果为该记录为后续按具体谕旨、奏折、条约、报刊或地方档案拆分事件提供可回查入口。。这一节点虽可放入改革、外交或革命的宏观年表，却不能脱离省份财政、地方官、商会、学校、军队、家庭与交通网络来理解；同一上谕、条约或战事在不同地域的实际后果并不一致。

核心材料为《清德宗实录》光绪十三年相关条；户部奏销、盐政、河工或海关档案。它们可以支持年序、制度文本、机构设立或某些地点的行动，却不能自动代表全体居民的收入、教育、身份、性别关系或政治意愿。账本提示“此行仅确认相关文书链和可追查的行政动作；不得由其直接推断政策有效、地方服从、民众态度或单一因果。”。因此，官方统计、外交措辞、报刊论战、教会叙述和军报数字必须分别置于生成它们的立场与制度目的之中，不能互相替代。

在地方尺度上，这一事件应被看作资源和信息重新配置的过程：经费能否落实，取决于税源和中介；学校、警政或新军能否运行，取决于人员、土地和日常协作；家庭与社群对改革的经历，则可能因性别、职业、族属、居处与迁徙而不同。现存材料只照亮少数路径，叙事须保留被登记者与未被登记者之间的差距。

账本条目\texttt{EQ-1888-NAVAL-AND-ARMY-REFORM}记录，光绪十四年（1888年），清发生“清廷讨论海军、陆军训练和军工体系的经费与制度问题”。账本将其背景说明为列强战争压力暴露既有军制不足，可见结果为改革奏议与实际编制、操练之间常有距离。这一节点虽可放入改革、外交或革命的宏观年表，却不能脱离省份财政、地方官、商会、学校、军队、家庭与交通网络来理解；同一上谕、条约或战事在不同地域的实际后果并不一致。

核心材料为《清德宗实录》光绪十四年相关条；宫中朱批奏折、内阁大库题本与《清史稿》相关志传。它们可以支持年序、制度文本、机构设立或某些地点的行动，却不能自动代表全体居民的收入、教育、身份、性别关系或政治意愿。账本提示“以光绪十四年（1888年）的同期文书为定位起点；官修记录、奏报或外交文本服务特定机构立场，具体日期、规模、地方执行与对手视角须以独立材料复核。”。因此，官方统计、外交措辞、报刊论战、教会叙述和军报数字必须分别置于生成它们的立场与制度目的之中，不能互相替代。

在地方尺度上，这一事件应被看作资源和信息重新配置的过程：经费能否落实，取决于税源和中介；学校、警政或新军能否运行，取决于人员、土地和日常协作；家庭与社群对改革的经历，则可能因性别、职业、族属、居处与迁徙而不同。现存材料只照亮少数路径，叙事须保留被登记者与未被登记者之间的差距。

账本条目\texttt{EQ-1888-NAVAL-AND-ARMY-REFORM-DOCUMENTS}记录，光绪十四年（1888年），清发生“围绕「清廷讨论海军、陆军训练和军工体系的经费与制度问题」的上谕、奏折与部议在中央—地方行政链中流转”。账本将其背景说明为核心事件「清廷讨论海军、陆军训练和军工体系的经费与制度问题」要求中央或地方通过文书处理资源、信息或责任。，可见结果为该记录为后续按具体谕旨、奏折、条约、报刊或地方档案拆分事件提供可回查入口。。这一节点虽可放入改革、外交或革命的宏观年表，却不能脱离省份财政、地方官、商会、学校、军队、家庭与交通网络来理解；同一上谕、条约或战事在不同地域的实际后果并不一致。

核心材料为《清德宗实录》光绪十四年相关条；宫中朱批奏折、内阁大库题本与《清史稿》相关志传。它们可以支持年序、制度文本、机构设立或某些地点的行动，却不能自动代表全体居民的收入、教育、身份、性别关系或政治意愿。账本提示“此行仅确认相关文书链和可追查的行政动作；不得由其直接推断政策有效、地方服从、民众态度或单一因果。”。因此，官方统计、外交措辞、报刊论战、教会叙述和军报数字必须分别置于生成它们的立场与制度目的之中，不能互相替代。

在地方尺度上，这一事件应被看作资源和信息重新配置的过程：经费能否落实，取决于税源和中介；学校、警政或新军能否运行，取决于人员、土地和日常协作；家庭与社群对改革的经历，则可能因性别、职业、族属、居处与迁徙而不同。现存材料只照亮少数路径，叙事须保留被登记者与未被登记者之间的差距。

账本条目\texttt{EQ-1889-GUANGXU-MARRIAGE}记录，光绪十五年（1889年），清发生“光绪帝大婚并开始形式上的亲政安排”。账本将其背景说明为皇帝成年要求完成宫廷礼制程序，可见结果为慈禧仍对重要决策保有深刻影响。这一节点虽可放入改革、外交或革命的宏观年表，却不能脱离省份财政、地方官、商会、学校、军队、家庭与交通网络来理解；同一上谕、条约或战事在不同地域的实际后果并不一致。

核心材料为《清德宗实录》光绪十五年相关条；宫中朱批奏折、内阁大库题本与《清史稿》相关志传。它们可以支持年序、制度文本、机构设立或某些地点的行动，却不能自动代表全体居民的收入、教育、身份、性别关系或政治意愿。账本提示“以光绪十五年（1889年）的同期文书为定位起点；官修记录、奏报或外交文本服务特定机构立场，具体日期、规模、地方执行与对手视角须以独立材料复核。”。因此，官方统计、外交措辞、报刊论战、教会叙述和军报数字必须分别置于生成它们的立场与制度目的之中，不能互相替代。

在地方尺度上，这一事件应被看作资源和信息重新配置的过程：经费能否落实，取决于税源和中介；学校、警政或新军能否运行，取决于人员、土地和日常协作；家庭与社群对改革的经历，则可能因性别、职业、族属、居处与迁徙而不同。现存材料只照亮少数路径，叙事须保留被登记者与未被登记者之间的差距。

账本条目\texttt{EQ-1889-GUANGXU-MARRIAGE-DOCUMENTS}记录，光绪十五年（1889年），清发生“围绕「光绪帝大婚并开始形式上的亲政安排」的上谕、奏折与部议在中央—地方行政链中流转”。账本将其背景说明为核心事件「光绪帝大婚并开始形式上的亲政安排」要求中央或地方通过文书处理资源、信息或责任。，可见结果为该记录为后续按具体谕旨、奏折、条约、报刊或地方档案拆分事件提供可回查入口。。这一节点虽可放入改革、外交或革命的宏观年表，却不能脱离省份财政、地方官、商会、学校、军队、家庭与交通网络来理解；同一上谕、条约或战事在不同地域的实际后果并不一致。

核心材料为《清德宗实录》光绪十五年相关条；宫中朱批奏折、内阁大库题本与《清史稿》相关志传。它们可以支持年序、制度文本、机构设立或某些地点的行动，却不能自动代表全体居民的收入、教育、身份、性别关系或政治意愿。账本提示“此行仅确认相关文书链和可追查的行政动作；不得由其直接推断政策有效、地方服从、民众态度或单一因果。”。因此，官方统计、外交措辞、报刊论战、教会叙述和军报数字必须分别置于生成它们的立场与制度目的之中，不能互相替代。

在地方尺度上，这一事件应被看作资源和信息重新配置的过程：经费能否落实，取决于税源和中介；学校、警政或新军能否运行，取决于人员、土地和日常协作；家庭与社群对改革的经历，则可能因性别、职业、族属、居处与迁徙而不同。现存材料只照亮少数路径，叙事须保留被登记者与未被登记者之间的差距。

账本条目\texttt{EQ-1890-KOREAN-POLITICAL-CONFLICT}记录，光绪十六年（1890年），清、朝鲜与日本发生“朝鲜内政纷争持续牵动清日两国在半岛的影响力竞争”。账本将其背景说明为壬午事变后清日均扩大驻军与交涉，可见结果为朝鲜行动者不是外部竞争的被动对象，应纳入叙述。这一节点虽可放入改革、外交或革命的宏观年表，却不能脱离省份财政、地方官、商会、学校、军队、家庭与交通网络来理解；同一上谕、条约或战事在不同地域的实际后果并不一致。

核心材料为《清德宗实录》光绪十六年相关条；《筹办夷务始末》相关年卷与中外照会、条约文本。它们可以支持年序、制度文本、机构设立或某些地点的行动，却不能自动代表全体居民的收入、教育、身份、性别关系或政治意愿。账本提示“以光绪十六年（1890年）的同期文书为定位起点；官修记录、奏报或外交文本服务特定机构立场，具体日期、规模、地方执行与对手视角须以独立材料复核。”。因此，官方统计、外交措辞、报刊论战、教会叙述和军报数字必须分别置于生成它们的立场与制度目的之中，不能互相替代。

在地方尺度上，这一事件应被看作资源和信息重新配置的过程：经费能否落实，取决于税源和中介；学校、警政或新军能否运行，取决于人员、土地和日常协作；家庭与社群对改革的经历，则可能因性别、职业、族属、居处与迁徙而不同。现存材料只照亮少数路径，叙事须保留被登记者与未被登记者之间的差距。

账本条目\texttt{EQ-1890-KOREAN-POLITICAL-CONFLICT-DOCUMENTS}记录，光绪十六年（1890年），清、朝鲜与日本发生“围绕「朝鲜内政纷争持续牵动清日两国在半岛的影响力竞争」的照会、条约文本、换约或地方执行报告分别进入外交文书链”。账本将其背景说明为核心事件「朝鲜内政纷争持续牵动清日两国在半岛的影响力竞争」要求中央或地方通过文书处理资源、信息或责任。，可见结果为该记录为后续按具体谕旨、奏折、条约、报刊或地方档案拆分事件提供可回查入口。。这一节点虽可放入改革、外交或革命的宏观年表，却不能脱离省份财政、地方官、商会、学校、军队、家庭与交通网络来理解；同一上谕、条约或战事在不同地域的实际后果并不一致。

核心材料为《清德宗实录》光绪十六年相关条；《筹办夷务始末》相关年卷与中外照会、条约文本。它们可以支持年序、制度文本、机构设立或某些地点的行动，却不能自动代表全体居民的收入、教育、身份、性别关系或政治意愿。账本提示“此行仅确认相关文书链和可追查的行政动作；不得由其直接推断政策有效、地方服从、民众态度或单一因果。”。因此，官方统计、外交措辞、报刊论战、教会叙述和军报数字必须分别置于生成它们的立场与制度目的之中，不能互相替代。

在地方尺度上，这一事件应被看作资源和信息重新配置的过程：经费能否落实，取决于税源和中介；学校、警政或新军能否运行，取决于人员、土地和日常协作；家庭与社群对改革的经历，则可能因性别、职业、族属、居处与迁徙而不同。现存材料只照亮少数路径，叙事须保留被登记者与未被登记者之间的差距。

账本条目\texttt{EQ-1891-MISSIONARY-CONFLICTS}记录，光绪十七年（1891年），清与各地教民、传教士发生“各地教案和宗教纠纷继续引发地方治理与外交压力”。账本将其背景说明为土地、诉讼、保护权和谣言与宗教身份交织，可见结果为“教案”包含不同类型纠纷，不能单一化。这一节点虽可放入改革、外交或革命的宏观年表，却不能脱离省份财政、地方官、商会、学校、军队、家庭与交通网络来理解；同一上谕、条约或战事在不同地域的实际后果并不一致。

核心材料为《清德宗实录》光绪十七年相关条；地方志、督抚奏折及地方档案。它们可以支持年序、制度文本、机构设立或某些地点的行动，却不能自动代表全体居民的收入、教育、身份、性别关系或政治意愿。账本提示“以光绪十七年（1891年）的同期文书为定位起点；官修记录、奏报或外交文本服务特定机构立场，具体日期、规模、地方执行与对手视角须以独立材料复核。”。因此，官方统计、外交措辞、报刊论战、教会叙述和军报数字必须分别置于生成它们的立场与制度目的之中，不能互相替代。

在地方尺度上，这一事件应被看作资源和信息重新配置的过程：经费能否落实，取决于税源和中介；学校、警政或新军能否运行，取决于人员、土地和日常协作；家庭与社群对改革的经历，则可能因性别、职业、族属、居处与迁徙而不同。现存材料只照亮少数路径，叙事须保留被登记者与未被登记者之间的差距。

账本条目\texttt{EQ-1891-MISSIONARY-CONFLICTS-DOCUMENTS}记录，光绪十七年（1891年），清与各地教民、传教士发生“围绕「各地教案和宗教纠纷继续引发地方治理与外交压力」的地方呈报、案件或赈务记录将社会反应带入官府文书链”。账本将其背景说明为核心事件「各地教案和宗教纠纷继续引发地方治理与外交压力」要求中央或地方通过文书处理资源、信息或责任。，可见结果为该记录为后续按具体谕旨、奏折、条约、报刊或地方档案拆分事件提供可回查入口。。这一节点虽可放入改革、外交或革命的宏观年表，却不能脱离省份财政、地方官、商会、学校、军队、家庭与交通网络来理解；同一上谕、条约或战事在不同地域的实际后果并不一致。

核心材料为《清德宗实录》光绪十七年相关条；地方志、督抚奏折及地方档案。它们可以支持年序、制度文本、机构设立或某些地点的行动，却不能自动代表全体居民的收入、教育、身份、性别关系或政治意愿。账本提示“此行仅确认相关文书链和可追查的行政动作；不得由其直接推断政策有效、地方服从、民众态度或单一因果。”。因此，官方统计、外交措辞、报刊论战、教会叙述和军报数字必须分别置于生成它们的立场与制度目的之中，不能互相替代。

在地方尺度上，这一事件应被看作资源和信息重新配置的过程：经费能否落实，取决于税源和中介；学校、警政或新军能否运行，取决于人员、土地和日常协作；家庭与社群对改革的经历，则可能因性别、职业、族属、居处与迁徙而不同。现存材料只照亮少数路径，叙事须保留被登记者与未被登记者之间的差距。

账本条目\texttt{EQ-1892-RAILWAY-DEBATE}记录，光绪十八年（1892年），清发生“官员、商人和士绅围绕铁路修筑、主权与资金展开争论”。账本将其背景说明为交通技术、外债和地方利益引发新问题，可见结果为赞成或反对文本不等于实际线路建设和公众意见。这一节点虽可放入改革、外交或革命的宏观年表，却不能脱离省份财政、地方官、商会、学校、军队、家庭与交通网络来理解；同一上谕、条约或战事在不同地域的实际后果并不一致。

核心材料为《清德宗实录》光绪十八年相关条；户部奏销、盐政、河工或海关档案。它们可以支持年序、制度文本、机构设立或某些地点的行动，却不能自动代表全体居民的收入、教育、身份、性别关系或政治意愿。账本提示“以光绪十八年（1892年）的同期文书为定位起点；官修记录、奏报或外交文本服务特定机构立场，具体日期、规模、地方执行与对手视角须以独立材料复核。”。因此，官方统计、外交措辞、报刊论战、教会叙述和军报数字必须分别置于生成它们的立场与制度目的之中，不能互相替代。

在地方尺度上，这一事件应被看作资源和信息重新配置的过程：经费能否落实，取决于税源和中介；学校、警政或新军能否运行，取决于人员、土地和日常协作；家庭与社群对改革的经历，则可能因性别、职业、族属、居处与迁徙而不同。现存材料只照亮少数路径，叙事须保留被登记者与未被登记者之间的差距。

账本条目\texttt{EQ-1892-RAILWAY-DEBATE-DOCUMENTS}记录，光绪十八年（1892年），清发生“围绕「官员、商人和士绅围绕铁路修筑、主权与资金展开争论」的经费、税收、赈济或工程呈报经由户部与地方衙门处理”。账本将其背景说明为核心事件「官员、商人和士绅围绕铁路修筑、主权与资金展开争论」要求中央或地方通过文书处理资源、信息或责任。，可见结果为该记录为后续按具体谕旨、奏折、条约、报刊或地方档案拆分事件提供可回查入口。。这一节点虽可放入改革、外交或革命的宏观年表，却不能脱离省份财政、地方官、商会、学校、军队、家庭与交通网络来理解；同一上谕、条约或战事在不同地域的实际后果并不一致。

核心材料为《清德宗实录》光绪十八年相关条；户部奏销、盐政、河工或海关档案。它们可以支持年序、制度文本、机构设立或某些地点的行动，却不能自动代表全体居民的收入、教育、身份、性别关系或政治意愿。账本提示“此行仅确认相关文书链和可追查的行政动作；不得由其直接推断政策有效、地方服从、民众态度或单一因果。”。因此，官方统计、外交措辞、报刊论战、教会叙述和军报数字必须分别置于生成它们的立场与制度目的之中，不能互相替代。

在地方尺度上，这一事件应被看作资源和信息重新配置的过程：经费能否落实，取决于税源和中介；学校、警政或新军能否运行，取决于人员、土地和日常协作；家庭与社群对改革的经历，则可能因性别、职业、族属、居处与迁徙而不同。现存材料只照亮少数路径，叙事须保留被登记者与未被登记者之间的差距。

账本条目\texttt{EQ-1893-KOREAN-REFORM-CRISIS}记录，光绪十九年（1893年），朝鲜、清与日本发生“朝鲜改革与政治危机使清日关系持续紧张”。账本将其背景说明为半岛内政和驻军安排缺乏稳定共识，可见结果为清日文书各自服务战略目的，需与朝鲜材料互读。这一节点虽可放入改革、外交或革命的宏观年表，却不能脱离省份财政、地方官、商会、学校、军队、家庭与交通网络来理解；同一上谕、条约或战事在不同地域的实际后果并不一致。

核心材料为《清德宗实录》光绪十九年相关条；《筹办夷务始末》相关年卷与中外照会、条约文本。它们可以支持年序、制度文本、机构设立或某些地点的行动，却不能自动代表全体居民的收入、教育、身份、性别关系或政治意愿。账本提示“以光绪十九年（1893年）的同期文书为定位起点；官修记录、奏报或外交文本服务特定机构立场，具体日期、规模、地方执行与对手视角须以独立材料复核。”。因此，官方统计、外交措辞、报刊论战、教会叙述和军报数字必须分别置于生成它们的立场与制度目的之中，不能互相替代。

在地方尺度上，这一事件应被看作资源和信息重新配置的过程：经费能否落实，取决于税源和中介；学校、警政或新军能否运行，取决于人员、土地和日常协作；家庭与社群对改革的经历，则可能因性别、职业、族属、居处与迁徙而不同。现存材料只照亮少数路径，叙事须保留被登记者与未被登记者之间的差距。

账本条目\texttt{EQ-1893-KOREAN-REFORM-CRISIS-DOCUMENTS}记录，光绪十九年（1893年），朝鲜、清与日本发生“围绕「朝鲜改革与政治危机使清日关系持续紧张」的照会、条约文本、换约或地方执行报告分别进入外交文书链”。账本将其背景说明为核心事件「朝鲜改革与政治危机使清日关系持续紧张」要求中央或地方通过文书处理资源、信息或责任。，可见结果为该记录为后续按具体谕旨、奏折、条约、报刊或地方档案拆分事件提供可回查入口。。这一节点虽可放入改革、外交或革命的宏观年表，却不能脱离省份财政、地方官、商会、学校、军队、家庭与交通网络来理解；同一上谕、条约或战事在不同地域的实际后果并不一致。

核心材料为《清德宗实录》光绪十九年相关条；《筹办夷务始末》相关年卷与中外照会、条约文本。它们可以支持年序、制度文本、机构设立或某些地点的行动，却不能自动代表全体居民的收入、教育、身份、性别关系或政治意愿。账本提示“此行仅确认相关文书链和可追查的行政动作；不得由其直接推断政策有效、地方服从、民众态度或单一因果。”。因此，官方统计、外交措辞、报刊论战、教会叙述和军报数字必须分别置于生成它们的立场与制度目的之中，不能互相替代。

在地方尺度上，这一事件应被看作资源和信息重新配置的过程：经费能否落实，取决于税源和中介；学校、警政或新军能否运行，取决于人员、土地和日常协作；家庭与社群对改革的经历，则可能因性别、职业、族属、居处与迁徙而不同。现存材料只照亮少数路径，叙事须保留被登记者与未被登记者之间的差距。

账本条目\texttt{EQ-1894-DONGHAK-PEASANT-WAR}记录，光绪二十年（1894年），朝鲜、清与日本发生“东学农民战争促使清日两国派兵入朝，危机迅速国际化”。账本将其背景说明为朝鲜社会改革诉求与政府镇压失败交织，可见结果为农民运动不应只作为清日战争的前奏叙述。这一节点虽可放入改革、外交或革命的宏观年表，却不能脱离省份财政、地方官、商会、学校、军队、家庭与交通网络来理解；同一上谕、条约或战事在不同地域的实际后果并不一致。

核心材料为《清德宗实录》光绪二十年相关条；军机处档、战事方略及地方奏报。它们可以支持年序、制度文本、机构设立或某些地点的行动，却不能自动代表全体居民的收入、教育、身份、性别关系或政治意愿。账本提示“以光绪二十年（1894年）的同期文书为定位起点；官修记录、奏报或外交文本服务特定机构立场，具体日期、规模、地方执行与对手视角须以独立材料复核。”。因此，官方统计、外交措辞、报刊论战、教会叙述和军报数字必须分别置于生成它们的立场与制度目的之中，不能互相替代。

在地方尺度上，这一事件应被看作资源和信息重新配置的过程：经费能否落实，取决于税源和中介；学校、警政或新军能否运行，取决于人员、土地和日常协作；家庭与社群对改革的经历，则可能因性别、职业、族属、居处与迁徙而不同。现存材料只照亮少数路径，叙事须保留被登记者与未被登记者之间的差距。

账本条目\texttt{EQ-1894-DONGHAK-PEASANT-WAR-DOCUMENTS}记录，光绪二十年（1894年），朝鲜、清与日本发生“清廷围绕「东学农民战争促使清日两国派兵入朝，危机迅速国际化」调遣军队、筹措饷械并处理战报，中央与地方的军政文书形成连续记录”。账本将其背景说明为核心事件「东学农民战争促使清日两国派兵入朝，危机迅速国际化」要求中央或地方通过文书处理资源、信息或责任。，可见结果为该记录为后续按具体谕旨、奏折、条约、报刊或地方档案拆分事件提供可回查入口。。这一节点虽可放入改革、外交或革命的宏观年表，却不能脱离省份财政、地方官、商会、学校、军队、家庭与交通网络来理解；同一上谕、条约或战事在不同地域的实际后果并不一致。

核心材料为《清德宗实录》光绪二十年相关条；军机处档、战事方略及地方奏报。它们可以支持年序、制度文本、机构设立或某些地点的行动，却不能自动代表全体居民的收入、教育、身份、性别关系或政治意愿。账本提示“此行仅确认相关文书链和可追查的行政动作；不得由其直接推断政策有效、地方服从、民众态度或单一因果。”。因此，官方统计、外交措辞、报刊论战、教会叙述和军报数字必须分别置于生成它们的立场与制度目的之中，不能互相替代。

在地方尺度上，这一事件应被看作资源和信息重新配置的过程：经费能否落实，取决于税源和中介；学校、警政或新军能否运行，取决于人员、土地和日常协作；家庭与社群对改革的经历，则可能因性别、职业、族属、居处与迁徙而不同。现存材料只照亮少数路径，叙事须保留被登记者与未被登记者之间的差距。

账本条目\texttt{EQ-1894-FIRST-SINO-JAPANESE-WAR}记录，光绪二十年（1894年），清与日本发生“甲午战争爆发，清日围绕朝鲜和区域秩序展开海陆战”。账本将其背景说明为驻军竞争与外交谈判破裂升级为战争，可见结果为战场损失、指挥责任与民众经验须多源核验。这一节点虽可放入改革、外交或革命的宏观年表，却不能脱离省份财政、地方官、商会、学校、军队、家庭与交通网络来理解；同一上谕、条约或战事在不同地域的实际后果并不一致。

核心材料为《清德宗实录》光绪二十年相关条；军机处档、战事方略及地方奏报。它们可以支持年序、制度文本、机构设立或某些地点的行动，却不能自动代表全体居民的收入、教育、身份、性别关系或政治意愿。账本提示“以光绪二十年（1894年）的同期文书为定位起点；官修记录、奏报或外交文本服务特定机构立场，具体日期、规模、地方执行与对手视角须以独立材料复核。”。因此，官方统计、外交措辞、报刊论战、教会叙述和军报数字必须分别置于生成它们的立场与制度目的之中，不能互相替代。

在地方尺度上，这一事件应被看作资源和信息重新配置的过程：经费能否落实，取决于税源和中介；学校、警政或新军能否运行，取决于人员、土地和日常协作；家庭与社群对改革的经历，则可能因性别、职业、族属、居处与迁徙而不同。现存材料只照亮少数路径，叙事须保留被登记者与未被登记者之间的差距。

账本条目\texttt{EQ-1894-FIRST-SINO-JAPANESE-WAR-DOCUMENTS}记录，光绪二十年（1894年），清与日本发生“清廷围绕「甲午战争爆发，清日围绕朝鲜和区域秩序展开海陆战」调遣军队、筹措饷械并处理战报，中央与地方的军政文书形成连续记录”。账本将其背景说明为核心事件「甲午战争爆发，清日围绕朝鲜和区域秩序展开海陆战」要求中央或地方通过文书处理资源、信息或责任。，可见结果为该记录为后续按具体谕旨、奏折、条约、报刊或地方档案拆分事件提供可回查入口。。这一节点虽可放入改革、外交或革命的宏观年表，却不能脱离省份财政、地方官、商会、学校、军队、家庭与交通网络来理解；同一上谕、条约或战事在不同地域的实际后果并不一致。

核心材料为《清德宗实录》光绪二十年相关条；军机处档、战事方略及地方奏报。它们可以支持年序、制度文本、机构设立或某些地点的行动，却不能自动代表全体居民的收入、教育、身份、性别关系或政治意愿。账本提示“此行仅确认相关文书链和可追查的行政动作；不得由其直接推断政策有效、地方服从、民众态度或单一因果。”。因此，官方统计、外交措辞、报刊论战、教会叙述和军报数字必须分别置于生成它们的立场与制度目的之中，不能互相替代。

在地方尺度上，这一事件应被看作资源和信息重新配置的过程：经费能否落实，取决于税源和中介；学校、警政或新军能否运行，取决于人员、土地和日常协作；家庭与社群对改革的经历，则可能因性别、职业、族属、居处与迁徙而不同。现存材料只照亮少数路径，叙事须保留被登记者与未被登记者之间的差距。

账本条目\texttt{EQ-1894-YALU-NAVAL-BATTLE}记录，光绪二十年（1894年），清与日本发生“黄海海战后北洋舰队受创，海军力量对比进一步恶化”。账本将其背景说明为舰艇性能、训练、战术和指挥因素共同作用，可见结果为单一“装备落后”解释无法涵盖战役复杂性。这一节点虽可放入改革、外交或革命的宏观年表，却不能脱离省份财政、地方官、商会、学校、军队、家庭与交通网络来理解；同一上谕、条约或战事在不同地域的实际后果并不一致。

核心材料为《清德宗实录》光绪二十年相关条；军机处档、战事方略及地方奏报。它们可以支持年序、制度文本、机构设立或某些地点的行动，却不能自动代表全体居民的收入、教育、身份、性别关系或政治意愿。账本提示“以光绪二十年（1894年）的同期文书为定位起点；官修记录、奏报或外交文本服务特定机构立场，具体日期、规模、地方执行与对手视角须以独立材料复核。”。因此，官方统计、外交措辞、报刊论战、教会叙述和军报数字必须分别置于生成它们的立场与制度目的之中，不能互相替代。

在地方尺度上，这一事件应被看作资源和信息重新配置的过程：经费能否落实，取决于税源和中介；学校、警政或新军能否运行，取决于人员、土地和日常协作；家庭与社群对改革的经历，则可能因性别、职业、族属、居处与迁徙而不同。现存材料只照亮少数路径，叙事须保留被登记者与未被登记者之间的差距。

账本条目\texttt{EQ-1894-YALU-NAVAL-BATTLE-DOCUMENTS}记录，光绪二十年（1894年），清与日本发生“清廷围绕「黄海海战后北洋舰队受创，海军力量对比进一步恶化」调遣军队、筹措饷械并处理战报，中央与地方的军政文书形成连续记录”。账本将其背景说明为核心事件「黄海海战后北洋舰队受创，海军力量对比进一步恶化」要求中央或地方通过文书处理资源、信息或责任。，可见结果为该记录为后续按具体谕旨、奏折、条约、报刊或地方档案拆分事件提供可回查入口。。这一节点虽可放入改革、外交或革命的宏观年表，却不能脱离省份财政、地方官、商会、学校、军队、家庭与交通网络来理解；同一上谕、条约或战事在不同地域的实际后果并不一致。

核心材料为《清德宗实录》光绪二十年相关条；军机处档、战事方略及地方奏报。它们可以支持年序、制度文本、机构设立或某些地点的行动，却不能自动代表全体居民的收入、教育、身份、性别关系或政治意愿。账本提示“此行仅确认相关文书链和可追查的行政动作；不得由其直接推断政策有效、地方服从、民众态度或单一因果。”。因此，官方统计、外交措辞、报刊论战、教会叙述和军报数字必须分别置于生成它们的立场与制度目的之中，不能互相替代。

在地方尺度上，这一事件应被看作资源和信息重新配置的过程：经费能否落实，取决于税源和中介；学校、警政或新军能否运行，取决于人员、土地和日常协作；家庭与社群对改革的经历，则可能因性别、职业、族属、居处与迁徙而不同。现存材料只照亮少数路径，叙事须保留被登记者与未被登记者之间的差距。

账本条目\texttt{EQ-1895-SHIMONOSEKI-TREATY}记录，光绪二十一年（1895年），清与日本发生“《马关条约》签订，清割让台湾、澎湖并承认朝鲜独立，赔款与通商条件加重”。账本将其背景说明为清军战败和日本军事压力迫使议和，可见结果为条约的签署、割让交接和地方抵抗是不同过程。这一节点虽可放入改革、外交或革命的宏观年表，却不能脱离省份财政、地方官、商会、学校、军队、家庭与交通网络来理解；同一上谕、条约或战事在不同地域的实际后果并不一致。

核心材料为《清德宗实录》光绪二十一年相关条；《筹办夷务始末》相关年卷与中外照会、条约文本。它们可以支持年序、制度文本、机构设立或某些地点的行动，却不能自动代表全体居民的收入、教育、身份、性别关系或政治意愿。账本提示“以光绪二十一年（1895年）的同期文书为定位起点；官修记录、奏报或外交文本服务特定机构立场，具体日期、规模、地方执行与对手视角须以独立材料复核。”。因此，官方统计、外交措辞、报刊论战、教会叙述和军报数字必须分别置于生成它们的立场与制度目的之中，不能互相替代。

在地方尺度上，这一事件应被看作资源和信息重新配置的过程：经费能否落实，取决于税源和中介；学校、警政或新军能否运行，取决于人员、土地和日常协作；家庭与社群对改革的经历，则可能因性别、职业、族属、居处与迁徙而不同。现存材料只照亮少数路径，叙事须保留被登记者与未被登记者之间的差距。

账本条目\texttt{EQ-1895-SHIMONOSEKI-TREATY-DOCUMENTS}记录，光绪二十一年（1895年），清与日本发生“围绕「《马关条约》签订，清割让台湾、澎湖并承认朝鲜独立，赔款与通商条件加重」的照会、条约文本、换约或地方执行报告分别进入外交文书链”。账本将其背景说明为核心事件「《马关条约》签订，清割让台湾、澎湖并承认朝鲜独立，赔款与通商条件加重」要求中央或地方通过文书处理资源、信息或责任。，可见结果为该记录为后续按具体谕旨、奏折、条约、报刊或地方档案拆分事件提供可回查入口。。这一节点虽可放入改革、外交或革命的宏观年表，却不能脱离省份财政、地方官、商会、学校、军队、家庭与交通网络来理解；同一上谕、条约或战事在不同地域的实际后果并不一致。

核心材料为《清德宗实录》光绪二十一年相关条；《筹办夷务始末》相关年卷与中外照会、条约文本。它们可以支持年序、制度文本、机构设立或某些地点的行动，却不能自动代表全体居民的收入、教育、身份、性别关系或政治意愿。账本提示“此行仅确认相关文书链和可追查的行政动作；不得由其直接推断政策有效、地方服从、民众态度或单一因果。”。因此，官方统计、外交措辞、报刊论战、教会叙述和军报数字必须分别置于生成它们的立场与制度目的之中，不能互相替代。

在地方尺度上，这一事件应被看作资源和信息重新配置的过程：经费能否落实，取决于税源和中介；学校、警政或新军能否运行，取决于人员、土地和日常协作；家庭与社群对改革的经历，则可能因性别、职业、族属、居处与迁徙而不同。现存材料只照亮少数路径，叙事须保留被登记者与未被登记者之间的差距。

账本条目\texttt{EQ-1895-TAIWAN-REPUBLIC-RESISTANCE}记录，光绪二十一年（1895年），台湾地方社会、清与日本发生“台湾割让后地方力量建立台湾民主国并抵抗日军登陆”。账本将其背景说明为条约割让与地方社会的政治、军事反应不一致，可见结果为抵抗组织与普通居民处境不能由单一民族叙事概括。这一节点虽可放入改革、外交或革命的宏观年表，却不能脱离省份财政、地方官、商会、学校、军队、家庭与交通网络来理解；同一上谕、条约或战事在不同地域的实际后果并不一致。

核心材料为《清德宗实录》光绪二十一年相关条；军机处档、战事方略及地方奏报。它们可以支持年序、制度文本、机构设立或某些地点的行动，却不能自动代表全体居民的收入、教育、身份、性别关系或政治意愿。账本提示“以光绪二十一年（1895年）的同期文书为定位起点；官修记录、奏报或外交文本服务特定机构立场，具体日期、规模、地方执行与对手视角须以独立材料复核。”。因此，官方统计、外交措辞、报刊论战、教会叙述和军报数字必须分别置于生成它们的立场与制度目的之中，不能互相替代。

在地方尺度上，这一事件应被看作资源和信息重新配置的过程：经费能否落实，取决于税源和中介；学校、警政或新军能否运行，取决于人员、土地和日常协作；家庭与社群对改革的经历，则可能因性别、职业、族属、居处与迁徙而不同。现存材料只照亮少数路径，叙事须保留被登记者与未被登记者之间的差距。

账本条目\texttt{EQ-1895-TAIWAN-REPUBLIC-RESISTANCE-DOCUMENTS}记录，光绪二十一年（1895年），台湾地方社会、清与日本发生“清廷围绕「台湾割让后地方力量建立台湾民主国并抵抗日军登陆」调遣军队、筹措饷械并处理战报，中央与地方的军政文书形成连续记录”。账本将其背景说明为核心事件「台湾割让后地方力量建立台湾民主国并抵抗日军登陆」要求中央或地方通过文书处理资源、信息或责任。，可见结果为该记录为后续按具体谕旨、奏折、条约、报刊或地方档案拆分事件提供可回查入口。。这一节点虽可放入改革、外交或革命的宏观年表，却不能脱离省份财政、地方官、商会、学校、军队、家庭与交通网络来理解；同一上谕、条约或战事在不同地域的实际后果并不一致。

核心材料为《清德宗实录》光绪二十一年相关条；军机处档、战事方略及地方奏报。它们可以支持年序、制度文本、机构设立或某些地点的行动，却不能自动代表全体居民的收入、教育、身份、性别关系或政治意愿。账本提示“此行仅确认相关文书链和可追查的行政动作；不得由其直接推断政策有效、地方服从、民众态度或单一因果。”。因此，官方统计、外交措辞、报刊论战、教会叙述和军报数字必须分别置于生成它们的立场与制度目的之中，不能互相替代。

在地方尺度上，这一事件应被看作资源和信息重新配置的过程：经费能否落实，取决于税源和中介；学校、警政或新军能否运行，取决于人员、土地和日常协作；家庭与社群对改革的经历，则可能因性别、职业、族属、居处与迁徙而不同。现存材料只照亮少数路径，叙事须保留被登记者与未被登记者之间的差距。

账本条目\texttt{EQ-1895-TRIPLE-INTERVENTION}记录，光绪二十一年（1895年），俄、法、德与日本、清发生“三国干涉还辽改变战后列强竞争和清廷外交处境”。账本将其背景说明为列强在东北和东亚的利益冲突加剧，可见结果为干涉并非出于中立调停，需置于帝国竞争中理解。这一节点虽可放入改革、外交或革命的宏观年表，却不能脱离省份财政、地方官、商会、学校、军队、家庭与交通网络来理解；同一上谕、条约或战事在不同地域的实际后果并不一致。

核心材料为《清德宗实录》光绪二十一年相关条；《筹办夷务始末》相关年卷与中外照会、条约文本。它们可以支持年序、制度文本、机构设立或某些地点的行动，却不能自动代表全体居民的收入、教育、身份、性别关系或政治意愿。账本提示“以光绪二十一年（1895年）的同期文书为定位起点；官修记录、奏报或外交文本服务特定机构立场，具体日期、规模、地方执行与对手视角须以独立材料复核。”。因此，官方统计、外交措辞、报刊论战、教会叙述和军报数字必须分别置于生成它们的立场与制度目的之中，不能互相替代。

在地方尺度上，这一事件应被看作资源和信息重新配置的过程：经费能否落实，取决于税源和中介；学校、警政或新军能否运行，取决于人员、土地和日常协作；家庭与社群对改革的经历，则可能因性别、职业、族属、居处与迁徙而不同。现存材料只照亮少数路径，叙事须保留被登记者与未被登记者之间的差距。

账本条目\texttt{EQ-1895-TRIPLE-INTERVENTION-DOCUMENTS}记录，光绪二十一年（1895年），俄、法、德与日本、清发生“围绕「三国干涉还辽改变战后列强竞争和清廷外交处境」的照会、条约文本、换约或地方执行报告分别进入外交文书链”。账本将其背景说明为核心事件「三国干涉还辽改变战后列强竞争和清廷外交处境」要求中央或地方通过文书处理资源、信息或责任。，可见结果为该记录为后续按具体谕旨、奏折、条约、报刊或地方档案拆分事件提供可回查入口。。这一节点虽可放入改革、外交或革命的宏观年表，却不能脱离省份财政、地方官、商会、学校、军队、家庭与交通网络来理解；同一上谕、条约或战事在不同地域的实际后果并不一致。

核心材料为《清德宗实录》光绪二十一年相关条；《筹办夷务始末》相关年卷与中外照会、条约文本。它们可以支持年序、制度文本、机构设立或某些地点的行动，却不能自动代表全体居民的收入、教育、身份、性别关系或政治意愿。账本提示“此行仅确认相关文书链和可追查的行政动作；不得由其直接推断政策有效、地方服从、民众态度或单一因果。”。因此，官方统计、外交措辞、报刊论战、教会叙述和军报数字必须分别置于生成它们的立场与制度目的之中，不能互相替代。

在地方尺度上，这一事件应被看作资源和信息重新配置的过程：经费能否落实，取决于税源和中介；学校、警政或新军能否运行，取决于人员、土地和日常协作；家庭与社群对改革的经历，则可能因性别、职业、族属、居处与迁徙而不同。现存材料只照亮少数路径，叙事须保留被登记者与未被登记者之间的差距。

账本条目\texttt{EQ-1896-LI-LOBANOV-TREATY}记录，光绪二十二年（1896年），清与俄罗斯发生“《中俄密约》及东清铁路安排加强俄国在东北的影响”。账本将其背景说明为清廷寻求对日后的外交支撑，俄国谋求远东利益，可见结果为秘密条约的文本、执行与地方影响需分层说明。这一节点虽可放入改革、外交或革命的宏观年表，却不能脱离省份财政、地方官、商会、学校、军队、家庭与交通网络来理解；同一上谕、条约或战事在不同地域的实际后果并不一致。

核心材料为《清德宗实录》光绪二十二年相关条；《筹办夷务始末》相关年卷与中外照会、条约文本。它们可以支持年序、制度文本、机构设立或某些地点的行动，却不能自动代表全体居民的收入、教育、身份、性别关系或政治意愿。账本提示“以光绪二十二年（1896年）的同期文书为定位起点；官修记录、奏报或外交文本服务特定机构立场，具体日期、规模、地方执行与对手视角须以独立材料复核。”。因此，官方统计、外交措辞、报刊论战、教会叙述和军报数字必须分别置于生成它们的立场与制度目的之中，不能互相替代。

在地方尺度上，这一事件应被看作资源和信息重新配置的过程：经费能否落实，取决于税源和中介；学校、警政或新军能否运行，取决于人员、土地和日常协作；家庭与社群对改革的经历，则可能因性别、职业、族属、居处与迁徙而不同。现存材料只照亮少数路径，叙事须保留被登记者与未被登记者之间的差距。

账本条目\texttt{EQ-1896-LI-LOBANOV-TREATY-DOCUMENTS}记录，光绪二十二年（1896年），清与俄罗斯发生“围绕「《中俄密约》及东清铁路安排加强俄国在东北的影响」的照会、条约文本、换约或地方执行报告分别进入外交文书链”。账本将其背景说明为核心事件「《中俄密约》及东清铁路安排加强俄国在东北的影响」要求中央或地方通过文书处理资源、信息或责任。，可见结果为该记录为后续按具体谕旨、奏折、条约、报刊或地方档案拆分事件提供可回查入口。。这一节点虽可放入改革、外交或革命的宏观年表，却不能脱离省份财政、地方官、商会、学校、军队、家庭与交通网络来理解；同一上谕、条约或战事在不同地域的实际后果并不一致。

核心材料为《清德宗实录》光绪二十二年相关条；《筹办夷务始末》相关年卷与中外照会、条约文本。它们可以支持年序、制度文本、机构设立或某些地点的行动，却不能自动代表全体居民的收入、教育、身份、性别关系或政治意愿。账本提示“此行仅确认相关文书链和可追查的行政动作；不得由其直接推断政策有效、地方服从、民众态度或单一因果。”。因此，官方统计、外交措辞、报刊论战、教会叙述和军报数字必须分别置于生成它们的立场与制度目的之中，不能互相替代。

在地方尺度上，这一事件应被看作资源和信息重新配置的过程：经费能否落实，取决于税源和中介；学校、警政或新军能否运行，取决于人员、土地和日常协作；家庭与社群对改革的经历，则可能因性别、职业、族属、居处与迁徙而不同。现存材料只照亮少数路径，叙事须保留被登记者与未被登记者之间的差距。

账本条目\texttt{EQ-1896-MODERN-EDUCATION-EXPANSION}记录，光绪二十二年（1896年），清发生“甲午战后新式学堂、译书与留学讨论加速”。账本将其背景说明为战败促使士大夫和官员反思军政、教育与技术，可见结果为教育项目的地域覆盖和入学规模不宜夸大。这一节点虽可放入改革、外交或革命的宏观年表，却不能脱离省份财政、地方官、商会、学校、军队、家庭与交通网络来理解；同一上谕、条约或战事在不同地域的实际后果并不一致。

核心材料为《清德宗实录》光绪二十二年相关条；内阁档、文集或报刊相关记录。它们可以支持年序、制度文本、机构设立或某些地点的行动，却不能自动代表全体居民的收入、教育、身份、性别关系或政治意愿。账本提示“以光绪二十二年（1896年）的同期文书为定位起点；官修记录、奏报或外交文本服务特定机构立场，具体日期、规模、地方执行与对手视角须以独立材料复核。”。因此，官方统计、外交措辞、报刊论战、教会叙述和军报数字必须分别置于生成它们的立场与制度目的之中，不能互相替代。

在地方尺度上，这一事件应被看作资源和信息重新配置的过程：经费能否落实，取决于税源和中介；学校、警政或新军能否运行，取决于人员、土地和日常协作；家庭与社群对改革的经历，则可能因性别、职业、族属、居处与迁徙而不同。现存材料只照亮少数路径，叙事须保留被登记者与未被登记者之间的差距。

账本条目\texttt{EQ-1896-MODERN-EDUCATION-EXPANSION-DOCUMENTS}记录，光绪二十二年（1896年），清发生“围绕「甲午战后新式学堂、译书与留学讨论加速」的禁令、编纂、教育或传播活动留下中央和地方文本记录”。账本将其背景说明为核心事件「甲午战后新式学堂、译书与留学讨论加速」要求中央或地方通过文书处理资源、信息或责任。，可见结果为该记录为后续按具体谕旨、奏折、条约、报刊或地方档案拆分事件提供可回查入口。。这一节点虽可放入改革、外交或革命的宏观年表，却不能脱离省份财政、地方官、商会、学校、军队、家庭与交通网络来理解；同一上谕、条约或战事在不同地域的实际后果并不一致。

核心材料为《清德宗实录》光绪二十二年相关条；内阁档、文集或报刊相关记录。它们可以支持年序、制度文本、机构设立或某些地点的行动，却不能自动代表全体居民的收入、教育、身份、性别关系或政治意愿。账本提示“此行仅确认相关文书链和可追查的行政动作；不得由其直接推断政策有效、地方服从、民众态度或单一因果。”。因此，官方统计、外交措辞、报刊论战、教会叙述和军报数字必须分别置于生成它们的立场与制度目的之中，不能互相替代。

在地方尺度上，这一事件应被看作资源和信息重新配置的过程：经费能否落实，取决于税源和中介；学校、警政或新军能否运行，取决于人员、土地和日常协作；家庭与社群对改革的经历，则可能因性别、职业、族属、居处与迁徙而不同。现存材料只照亮少数路径，叙事须保留被登记者与未被登记者之间的差距。

账本条目\texttt{EQ-1897-JIAOZHOU-BAY-OCCUPATION}记录，光绪二十三年（1897年），清与德国发生“德国以巨野教案为由占领胶州湾，清廷面临新的租借压力”。账本将其背景说明为列强借教案和军舰扩张在华据点，可见结果为教案、军事占领和租借谈判应分别记录。这一节点虽可放入改革、外交或革命的宏观年表，却不能脱离省份财政、地方官、商会、学校、军队、家庭与交通网络来理解；同一上谕、条约或战事在不同地域的实际后果并不一致。

核心材料为《清德宗实录》光绪二十三年相关条；《筹办夷务始末》相关年卷与中外照会、条约文本。它们可以支持年序、制度文本、机构设立或某些地点的行动，却不能自动代表全体居民的收入、教育、身份、性别关系或政治意愿。账本提示“以光绪二十三年（1897年）的同期文书为定位起点；官修记录、奏报或外交文本服务特定机构立场，具体日期、规模、地方执行与对手视角须以独立材料复核。”。因此，官方统计、外交措辞、报刊论战、教会叙述和军报数字必须分别置于生成它们的立场与制度目的之中，不能互相替代。

在地方尺度上，这一事件应被看作资源和信息重新配置的过程：经费能否落实，取决于税源和中介；学校、警政或新军能否运行，取决于人员、土地和日常协作；家庭与社群对改革的经历，则可能因性别、职业、族属、居处与迁徙而不同。现存材料只照亮少数路径，叙事须保留被登记者与未被登记者之间的差距。

账本条目\texttt{EQ-1897-JIAOZHOU-BAY-OCCUPATION-DOCUMENTS}记录，光绪二十三年（1897年），清与德国发生“围绕「德国以巨野教案为由占领胶州湾，清廷面临新的租借压力」的照会、条约文本、换约或地方执行报告分别进入外交文书链”。账本将其背景说明为核心事件「德国以巨野教案为由占领胶州湾，清廷面临新的租借压力」要求中央或地方通过文书处理资源、信息或责任。，可见结果为该记录为后续按具体谕旨、奏折、条约、报刊或地方档案拆分事件提供可回查入口。。这一节点虽可放入改革、外交或革命的宏观年表，却不能脱离省份财政、地方官、商会、学校、军队、家庭与交通网络来理解；同一上谕、条约或战事在不同地域的实际后果并不一致。

核心材料为《清德宗实录》光绪二十三年相关条；《筹办夷务始末》相关年卷与中外照会、条约文本。它们可以支持年序、制度文本、机构设立或某些地点的行动，却不能自动代表全体居民的收入、教育、身份、性别关系或政治意愿。账本提示“此行仅确认相关文书链和可追查的行政动作；不得由其直接推断政策有效、地方服从、民众态度或单一因果。”。因此，官方统计、外交措辞、报刊论战、教会叙述和军报数字必须分别置于生成它们的立场与制度目的之中，不能互相替代。

在地方尺度上，这一事件应被看作资源和信息重新配置的过程：经费能否落实，取决于税源和中介；学校、警政或新军能否运行，取决于人员、土地和日常协作；家庭与社群对改革的经历，则可能因性别、职业、族属、居处与迁徙而不同。现存材料只照亮少数路径，叙事须保留被登记者与未被登记者之间的差距。

账本条目\texttt{EQ-1898-LEASED-TERRITORIES}记录，光绪二十四年（1898年），清与列强发生“俄、英、法、德相继取得租借地或势力范围安排”。账本将其背景说明为甲午后清廷军事弱势与列强竞争叠加，可见结果为“势力范围”并不自动等于完整主权转移。这一节点虽可放入改革、外交或革命的宏观年表，却不能脱离省份财政、地方官、商会、学校、军队、家庭与交通网络来理解；同一上谕、条约或战事在不同地域的实际后果并不一致。

核心材料为《清德宗实录》光绪二十四年相关条；《筹办夷务始末》相关年卷与中外照会、条约文本。它们可以支持年序、制度文本、机构设立或某些地点的行动，却不能自动代表全体居民的收入、教育、身份、性别关系或政治意愿。账本提示“以光绪二十四年（1898年）的同期文书为定位起点；官修记录、奏报或外交文本服务特定机构立场，具体日期、规模、地方执行与对手视角须以独立材料复核。”。因此，官方统计、外交措辞、报刊论战、教会叙述和军报数字必须分别置于生成它们的立场与制度目的之中，不能互相替代。

在地方尺度上，这一事件应被看作资源和信息重新配置的过程：经费能否落实，取决于税源和中介；学校、警政或新军能否运行，取决于人员、土地和日常协作；家庭与社群对改革的经历，则可能因性别、职业、族属、居处与迁徙而不同。现存材料只照亮少数路径，叙事须保留被登记者与未被登记者之间的差距。

账本条目\texttt{EQ-1898-LEASED-TERRITORIES-DOCUMENTS}记录，光绪二十四年（1898年），清与列强发生“围绕「俄、英、法、德相继取得租借地或势力范围安排」的照会、条约文本、换约或地方执行报告分别进入外交文书链”。账本将其背景说明为核心事件「俄、英、法、德相继取得租借地或势力范围安排」要求中央或地方通过文书处理资源、信息或责任。，可见结果为该记录为后续按具体谕旨、奏折、条约、报刊或地方档案拆分事件提供可回查入口。。这一节点虽可放入改革、外交或革命的宏观年表，却不能脱离省份财政、地方官、商会、学校、军队、家庭与交通网络来理解；同一上谕、条约或战事在不同地域的实际后果并不一致。

核心材料为《清德宗实录》光绪二十四年相关条；《筹办夷务始末》相关年卷与中外照会、条约文本。它们可以支持年序、制度文本、机构设立或某些地点的行动，却不能自动代表全体居民的收入、教育、身份、性别关系或政治意愿。账本提示“此行仅确认相关文书链和可追查的行政动作；不得由其直接推断政策有效、地方服从、民众态度或单一因果。”。因此，官方统计、外交措辞、报刊论战、教会叙述和军报数字必须分别置于生成它们的立场与制度目的之中，不能互相替代。

在地方尺度上，这一事件应被看作资源和信息重新配置的过程：经费能否落实，取决于税源和中介；学校、警政或新军能否运行，取决于人员、土地和日常协作；家庭与社群对改革的经历，则可能因性别、职业、族属、居处与迁徙而不同。现存材料只照亮少数路径，叙事须保留被登记者与未被登记者之间的差距。

账本条目\texttt{EQ-1898-HUNDRED-DAYS-REFORM}记录，光绪二十四年（1898年），清发生“光绪帝颁布变法诏令，推动官制、教育、经济与军事改革”。账本将其背景说明为甲午战败和维新思潮使改革诉求集中爆发，可见结果为诏令与实际实施之间时间短、阻力大，须区别方案与结果。这一节点虽可放入改革、外交或革命的宏观年表，却不能脱离省份财政、地方官、商会、学校、军队、家庭与交通网络来理解；同一上谕、条约或战事在不同地域的实际后果并不一致。

核心材料为《清德宗实录》光绪二十四年相关条；宫中朱批奏折、内阁大库题本与《清史稿》相关志传。它们可以支持年序、制度文本、机构设立或某些地点的行动，却不能自动代表全体居民的收入、教育、身份、性别关系或政治意愿。账本提示“以光绪二十四年（1898年）的同期文书为定位起点；官修记录、奏报或外交文本服务特定机构立场，具体日期、规模、地方执行与对手视角须以独立材料复核。”。因此，官方统计、外交措辞、报刊论战、教会叙述和军报数字必须分别置于生成它们的立场与制度目的之中，不能互相替代。

在地方尺度上，这一事件应被看作资源和信息重新配置的过程：经费能否落实，取决于税源和中介；学校、警政或新军能否运行，取决于人员、土地和日常协作；家庭与社群对改革的经历，则可能因性别、职业、族属、居处与迁徙而不同。现存材料只照亮少数路径，叙事须保留被登记者与未被登记者之间的差距。

账本条目\texttt{EQ-1898-HUNDRED-DAYS-REFORM-DOCUMENTS}记录，光绪二十四年（1898年），清发生“围绕「光绪帝颁布变法诏令，推动官制、教育、经济与军事改革」的上谕、奏折与部议在中央—地方行政链中流转”。账本将其背景说明为核心事件「光绪帝颁布变法诏令，推动官制、教育、经济与军事改革」要求中央或地方通过文书处理资源、信息或责任。，可见结果为该记录为后续按具体谕旨、奏折、条约、报刊或地方档案拆分事件提供可回查入口。。这一节点虽可放入改革、外交或革命的宏观年表，却不能脱离省份财政、地方官、商会、学校、军队、家庭与交通网络来理解；同一上谕、条约或战事在不同地域的实际后果并不一致。

核心材料为《清德宗实录》光绪二十四年相关条；宫中朱批奏折、内阁大库题本与《清史稿》相关志传。它们可以支持年序、制度文本、机构设立或某些地点的行动，却不能自动代表全体居民的收入、教育、身份、性别关系或政治意愿。账本提示“此行仅确认相关文书链和可追查的行政动作；不得由其直接推断政策有效、地方服从、民众态度或单一因果。”。因此，官方统计、外交措辞、报刊论战、教会叙述和军报数字必须分别置于生成它们的立场与制度目的之中，不能互相替代。

在地方尺度上，这一事件应被看作资源和信息重新配置的过程：经费能否落实，取决于税源和中介；学校、警政或新军能否运行，取决于人员、土地和日常协作；家庭与社群对改革的经历，则可能因性别、职业、族属、居处与迁徙而不同。现存材料只照亮少数路径，叙事须保留被登记者与未被登记者之间的差距。

账本条目\texttt{EQ-1898-WUXU-COUP}记录，光绪二十四年（1898年），清发生“慈禧发动政变，光绪帝被幽禁，戊戌变法主要措施被撤销或改写”。账本将其背景说明为改革触及官僚利益、皇权关系与军事权力，可见结果为政变叙述中的人物动机和密谋细节需谨慎处理。这一节点虽可放入改革、外交或革命的宏观年表，却不能脱离省份财政、地方官、商会、学校、军队、家庭与交通网络来理解；同一上谕、条约或战事在不同地域的实际后果并不一致。

核心材料为《清德宗实录》光绪二十四年相关条；宫中朱批奏折、内阁大库题本与《清史稿》相关志传。它们可以支持年序、制度文本、机构设立或某些地点的行动，却不能自动代表全体居民的收入、教育、身份、性别关系或政治意愿。账本提示“以光绪二十四年（1898年）的同期文书为定位起点；官修记录、奏报或外交文本服务特定机构立场，具体日期、规模、地方执行与对手视角须以独立材料复核。”。因此，官方统计、外交措辞、报刊论战、教会叙述和军报数字必须分别置于生成它们的立场与制度目的之中，不能互相替代。

在地方尺度上，这一事件应被看作资源和信息重新配置的过程：经费能否落实，取决于税源和中介；学校、警政或新军能否运行，取决于人员、土地和日常协作；家庭与社群对改革的经历，则可能因性别、职业、族属、居处与迁徙而不同。现存材料只照亮少数路径，叙事须保留被登记者与未被登记者之间的差距。

账本条目\texttt{EQ-1898-WUXU-COUP-DOCUMENTS}记录，光绪二十四年（1898年），清发生“围绕「慈禧发动政变，光绪帝被幽禁，戊戌变法主要措施被撤销或改写」的上谕、奏折与部议在中央—地方行政链中流转”。账本将其背景说明为核心事件「慈禧发动政变，光绪帝被幽禁，戊戌变法主要措施被撤销或改写」要求中央或地方通过文书处理资源、信息或责任。，可见结果为该记录为后续按具体谕旨、奏折、条约、报刊或地方档案拆分事件提供可回查入口。。这一节点虽可放入改革、外交或革命的宏观年表，却不能脱离省份财政、地方官、商会、学校、军队、家庭与交通网络来理解；同一上谕、条约或战事在不同地域的实际后果并不一致。

核心材料为《清德宗实录》光绪二十四年相关条；宫中朱批奏折、内阁大库题本与《清史稿》相关志传。它们可以支持年序、制度文本、机构设立或某些地点的行动，却不能自动代表全体居民的收入、教育、身份、性别关系或政治意愿。账本提示“此行仅确认相关文书链和可追查的行政动作；不得由其直接推断政策有效、地方服从、民众态度或单一因果。”。因此，官方统计、外交措辞、报刊论战、教会叙述和军报数字必须分别置于生成它们的立场与制度目的之中，不能互相替代。

在地方尺度上，这一事件应被看作资源和信息重新配置的过程：经费能否落实，取决于税源和中介；学校、警政或新军能否运行，取决于人员、土地和日常协作；家庭与社群对改革的经历，则可能因性别、职业、族属、居处与迁徙而不同。现存材料只照亮少数路径，叙事须保留被登记者与未被登记者之间的差距。

\subsection{边疆、外交与全球交换}

账本条目\texttt{EQ-1899-BOXER-UPRISING}记录，光绪二十五年（1899年），清、义和团与华北社会发生“义和团在华北扩展，反教、反洋冲突加剧”。账本将其背景说明为灾荒、地方武术结社、教案和列强压力相互作用，可见结果为“义和团”内部成分复杂，不能以单一意识形态概括。这一节点虽可放入改革、外交或革命的宏观年表，却不能脱离省份财政、地方官、商会、学校、军队、家庭与交通网络来理解；同一上谕、条约或战事在不同地域的实际后果并不一致。

核心材料为《清德宗实录》光绪二十五年相关条；地方志、督抚奏折及地方档案。它们可以支持年序、制度文本、机构设立或某些地点的行动，却不能自动代表全体居民的收入、教育、身份、性别关系或政治意愿。账本提示“以光绪二十五年（1899年）的同期文书为定位起点；官修记录、奏报或外交文本服务特定机构立场，具体日期、规模、地方执行与对手视角须以独立材料复核。”。因此，官方统计、外交措辞、报刊论战、教会叙述和军报数字必须分别置于生成它们的立场与制度目的之中，不能互相替代。

在地方尺度上，这一事件应被看作资源和信息重新配置的过程：经费能否落实，取决于税源和中介；学校、警政或新军能否运行，取决于人员、土地和日常协作；家庭与社群对改革的经历，则可能因性别、职业、族属、居处与迁徙而不同。现存材料只照亮少数路径，叙事须保留被登记者与未被登记者之间的差距。

账本条目\texttt{EQ-1899-BOXER-UPRISING-DOCUMENTS}记录，光绪二十五年（1899年），清、义和团与华北社会发生“围绕「义和团在华北扩展，反教、反洋冲突加剧」的地方呈报、案件或赈务记录将社会反应带入官府文书链”。账本将其背景说明为核心事件「义和团在华北扩展，反教、反洋冲突加剧」要求中央或地方通过文书处理资源、信息或责任。，可见结果为该记录为后续按具体谕旨、奏折、条约、报刊或地方档案拆分事件提供可回查入口。。这一节点虽可放入改革、外交或革命的宏观年表，却不能脱离省份财政、地方官、商会、学校、军队、家庭与交通网络来理解；同一上谕、条约或战事在不同地域的实际后果并不一致。

核心材料为《清德宗实录》光绪二十五年相关条；地方志、督抚奏折及地方档案。它们可以支持年序、制度文本、机构设立或某些地点的行动，却不能自动代表全体居民的收入、教育、身份、性别关系或政治意愿。账本提示“此行仅确认相关文书链和可追查的行政动作；不得由其直接推断政策有效、地方服从、民众态度或单一因果。”。因此，官方统计、外交措辞、报刊论战、教会叙述和军报数字必须分别置于生成它们的立场与制度目的之中，不能互相替代。

在地方尺度上，这一事件应被看作资源和信息重新配置的过程：经费能否落实，取决于税源和中介；学校、警政或新军能否运行，取决于人员、土地和日常协作；家庭与社群对改革的经历，则可能因性别、职业、族属、居处与迁徙而不同。现存材料只照亮少数路径，叙事须保留被登记者与未被登记者之间的差距。

账本条目\texttt{EQ-1900-BOXER-WAR}记录，光绪二十六年（1900年），清、义和团与八国联军发生“慈禧对列强宣战，八国联军入侵并占领北京”。账本将其背景说明为朝廷内部政策摇摆、义和团动员与外交危机相互激化，可见结果为战争责任、暴力与劫掠需同时纳入中外和地方材料。这一节点虽可放入改革、外交或革命的宏观年表，却不能脱离省份财政、地方官、商会、学校、军队、家庭与交通网络来理解；同一上谕、条约或战事在不同地域的实际后果并不一致。

核心材料为《清德宗实录》光绪二十六年相关条；军机处档、战事方略及地方奏报。它们可以支持年序、制度文本、机构设立或某些地点的行动，却不能自动代表全体居民的收入、教育、身份、性别关系或政治意愿。账本提示“以光绪二十六年（1900年）的同期文书为定位起点；官修记录、奏报或外交文本服务特定机构立场，具体日期、规模、地方执行与对手视角须以独立材料复核。”。因此，官方统计、外交措辞、报刊论战、教会叙述和军报数字必须分别置于生成它们的立场与制度目的之中，不能互相替代。

在地方尺度上，这一事件应被看作资源和信息重新配置的过程：经费能否落实，取决于税源和中介；学校、警政或新军能否运行，取决于人员、土地和日常协作；家庭与社群对改革的经历，则可能因性别、职业、族属、居处与迁徙而不同。现存材料只照亮少数路径，叙事须保留被登记者与未被登记者之间的差距。

账本条目\texttt{EQ-1900-BOXER-WAR-DOCUMENTS}记录，光绪二十六年（1900年），清、义和团与八国联军发生“清廷围绕「慈禧对列强宣战，八国联军入侵并占领北京」调遣军队、筹措饷械并处理战报，中央与地方的军政文书形成连续记录”。账本将其背景说明为核心事件「慈禧对列强宣战，八国联军入侵并占领北京」要求中央或地方通过文书处理资源、信息或责任。，可见结果为该记录为后续按具体谕旨、奏折、条约、报刊或地方档案拆分事件提供可回查入口。。这一节点虽可放入改革、外交或革命的宏观年表，却不能脱离省份财政、地方官、商会、学校、军队、家庭与交通网络来理解；同一上谕、条约或战事在不同地域的实际后果并不一致。

核心材料为《清德宗实录》光绪二十六年相关条；军机处档、战事方略及地方奏报。它们可以支持年序、制度文本、机构设立或某些地点的行动，却不能自动代表全体居民的收入、教育、身份、性别关系或政治意愿。账本提示“此行仅确认相关文书链和可追查的行政动作；不得由其直接推断政策有效、地方服从、民众态度或单一因果。”。因此，官方统计、外交措辞、报刊论战、教会叙述和军报数字必须分别置于生成它们的立场与制度目的之中，不能互相替代。

在地方尺度上，这一事件应被看作资源和信息重新配置的过程：经费能否落实，取决于税源和中介；学校、警政或新军能否运行，取决于人员、土地和日常协作；家庭与社群对改革的经历，则可能因性别、职业、族属、居处与迁徙而不同。现存材料只照亮少数路径，叙事须保留被登记者与未被登记者之间的差距。

账本条目\texttt{EQ-1900-SOUTHEAST-MUTUAL-PROTECTION}记录，光绪二十六年（1900年），清与东南督抚发生“东南互保使部分督抚与列强协商维持地方秩序”。账本将其背景说明为中央战争命令与地方对商贸、治安的现实考量冲突，可见结果为“互保”并非整齐统一的地方独立，须逐省考察。这一节点虽可放入改革、外交或革命的宏观年表，却不能脱离省份财政、地方官、商会、学校、军队、家庭与交通网络来理解；同一上谕、条约或战事在不同地域的实际后果并不一致。

核心材料为《清德宗实录》光绪二十六年相关条；宫中朱批奏折、内阁大库题本与《清史稿》相关志传。它们可以支持年序、制度文本、机构设立或某些地点的行动，却不能自动代表全体居民的收入、教育、身份、性别关系或政治意愿。账本提示“以光绪二十六年（1900年）的同期文书为定位起点；官修记录、奏报或外交文本服务特定机构立场，具体日期、规模、地方执行与对手视角须以独立材料复核。”。因此，官方统计、外交措辞、报刊论战、教会叙述和军报数字必须分别置于生成它们的立场与制度目的之中，不能互相替代。

在地方尺度上，这一事件应被看作资源和信息重新配置的过程：经费能否落实，取决于税源和中介；学校、警政或新军能否运行，取决于人员、土地和日常协作；家庭与社群对改革的经历，则可能因性别、职业、族属、居处与迁徙而不同。现存材料只照亮少数路径，叙事须保留被登记者与未被登记者之间的差距。

账本条目\texttt{EQ-1900-SOUTHEAST-MUTUAL-PROTECTION-DOCUMENTS}记录，光绪二十六年（1900年），清与东南督抚发生“围绕「东南互保使部分督抚与列强协商维持地方秩序」的上谕、奏折与部议在中央—地方行政链中流转”。账本将其背景说明为核心事件「东南互保使部分督抚与列强协商维持地方秩序」要求中央或地方通过文书处理资源、信息或责任。，可见结果为该记录为后续按具体谕旨、奏折、条约、报刊或地方档案拆分事件提供可回查入口。。这一节点虽可放入改革、外交或革命的宏观年表，却不能脱离省份财政、地方官、商会、学校、军队、家庭与交通网络来理解；同一上谕、条约或战事在不同地域的实际后果并不一致。

核心材料为《清德宗实录》光绪二十六年相关条；宫中朱批奏折、内阁大库题本与《清史稿》相关志传。它们可以支持年序、制度文本、机构设立或某些地点的行动，却不能自动代表全体居民的收入、教育、身份、性别关系或政治意愿。账本提示“此行仅确认相关文书链和可追查的行政动作；不得由其直接推断政策有效、地方服从、民众态度或单一因果。”。因此，官方统计、外交措辞、报刊论战、教会叙述和军报数字必须分别置于生成它们的立场与制度目的之中，不能互相替代。

在地方尺度上，这一事件应被看作资源和信息重新配置的过程：经费能否落实，取决于税源和中介；学校、警政或新军能否运行，取决于人员、土地和日常协作；家庭与社群对改革的经历，则可能因性别、职业、族属、居处与迁徙而不同。现存材料只照亮少数路径，叙事须保留被登记者与未被登记者之间的差距。

账本条目\texttt{EQ-1901-BOXER-PROTOCOL}记录，光绪二十七年（1901年），清与列强发生“《辛丑条约》签订，赔款、驻军和外交条件进一步限制清廷”。账本将其背景说明为北京陷落和议和谈判使清廷处于弱势，可见结果为赔款总额、支付机制与地方税收后果应分开分析。这一节点虽可放入改革、外交或革命的宏观年表，却不能脱离省份财政、地方官、商会、学校、军队、家庭与交通网络来理解；同一上谕、条约或战事在不同地域的实际后果并不一致。

核心材料为《清德宗实录》光绪二十七年相关条；《筹办夷务始末》相关年卷与中外照会、条约文本。它们可以支持年序、制度文本、机构设立或某些地点的行动，却不能自动代表全体居民的收入、教育、身份、性别关系或政治意愿。账本提示“以光绪二十七年（1901年）的同期文书为定位起点；官修记录、奏报或外交文本服务特定机构立场，具体日期、规模、地方执行与对手视角须以独立材料复核。”。因此，官方统计、外交措辞、报刊论战、教会叙述和军报数字必须分别置于生成它们的立场与制度目的之中，不能互相替代。

在地方尺度上，这一事件应被看作资源和信息重新配置的过程：经费能否落实，取决于税源和中介；学校、警政或新军能否运行，取决于人员、土地和日常协作；家庭与社群对改革的经历，则可能因性别、职业、族属、居处与迁徙而不同。现存材料只照亮少数路径，叙事须保留被登记者与未被登记者之间的差距。

账本条目\texttt{EQ-1901-BOXER-PROTOCOL-DOCUMENTS}记录，光绪二十七年（1901年），清与列强发生“围绕「《辛丑条约》签订，赔款、驻军和外交条件进一步限制清廷」的照会、条约文本、换约或地方执行报告分别进入外交文书链”。账本将其背景说明为核心事件「《辛丑条约》签订，赔款、驻军和外交条件进一步限制清廷」要求中央或地方通过文书处理资源、信息或责任。，可见结果为该记录为后续按具体谕旨、奏折、条约、报刊或地方档案拆分事件提供可回查入口。。这一节点虽可放入改革、外交或革命的宏观年表，却不能脱离省份财政、地方官、商会、学校、军队、家庭与交通网络来理解；同一上谕、条约或战事在不同地域的实际后果并不一致。

核心材料为《清德宗实录》光绪二十七年相关条；《筹办夷务始末》相关年卷与中外照会、条约文本。它们可以支持年序、制度文本、机构设立或某些地点的行动，却不能自动代表全体居民的收入、教育、身份、性别关系或政治意愿。账本提示“此行仅确认相关文书链和可追查的行政动作；不得由其直接推断政策有效、地方服从、民众态度或单一因果。”。因此，官方统计、外交措辞、报刊论战、教会叙述和军报数字必须分别置于生成它们的立场与制度目的之中，不能互相替代。

在地方尺度上，这一事件应被看作资源和信息重新配置的过程：经费能否落实，取决于税源和中介；学校、警政或新军能否运行，取决于人员、土地和日常协作；家庭与社群对改革的经历，则可能因性别、职业、族属、居处与迁徙而不同。现存材料只照亮少数路径，叙事须保留被登记者与未被登记者之间的差距。

账本条目\texttt{EQ-1901-ZONGLI-YAMEN-REPLACED}记录，光绪二十七年（1901年），清发生“总理衙门改为外务部，晚清外交机构进一步制度化”。账本将其背景说明为战后改革要求重组对外事务处理，可见结果为机构改名不等于外交权力和信息网络立即现代化。这一节点虽可放入改革、外交或革命的宏观年表，却不能脱离省份财政、地方官、商会、学校、军队、家庭与交通网络来理解；同一上谕、条约或战事在不同地域的实际后果并不一致。

核心材料为《清德宗实录》光绪二十七年相关条；宫中朱批奏折、内阁大库题本与《清史稿》相关志传。它们可以支持年序、制度文本、机构设立或某些地点的行动，却不能自动代表全体居民的收入、教育、身份、性别关系或政治意愿。账本提示“以光绪二十七年（1901年）的同期文书为定位起点；官修记录、奏报或外交文本服务特定机构立场，具体日期、规模、地方执行与对手视角须以独立材料复核。”。因此，官方统计、外交措辞、报刊论战、教会叙述和军报数字必须分别置于生成它们的立场与制度目的之中，不能互相替代。

在地方尺度上，这一事件应被看作资源和信息重新配置的过程：经费能否落实，取决于税源和中介；学校、警政或新军能否运行，取决于人员、土地和日常协作；家庭与社群对改革的经历，则可能因性别、职业、族属、居处与迁徙而不同。现存材料只照亮少数路径，叙事须保留被登记者与未被登记者之间的差距。

账本条目\texttt{EQ-1901-ZONGLI-YAMEN-REPLACED-DOCUMENTS}记录，光绪二十七年（1901年），清发生“围绕「总理衙门改为外务部，晚清外交机构进一步制度化」的上谕、奏折与部议在中央—地方行政链中流转”。账本将其背景说明为核心事件「总理衙门改为外务部，晚清外交机构进一步制度化」要求中央或地方通过文书处理资源、信息或责任。，可见结果为该记录为后续按具体谕旨、奏折、条约、报刊或地方档案拆分事件提供可回查入口。。这一节点虽可放入改革、外交或革命的宏观年表，却不能脱离省份财政、地方官、商会、学校、军队、家庭与交通网络来理解；同一上谕、条约或战事在不同地域的实际后果并不一致。

核心材料为《清德宗实录》光绪二十七年相关条；宫中朱批奏折、内阁大库题本与《清史稿》相关志传。它们可以支持年序、制度文本、机构设立或某些地点的行动，却不能自动代表全体居民的收入、教育、身份、性别关系或政治意愿。账本提示“此行仅确认相关文书链和可追查的行政动作；不得由其直接推断政策有效、地方服从、民众态度或单一因果。”。因此，官方统计、外交措辞、报刊论战、教会叙述和军报数字必须分别置于生成它们的立场与制度目的之中，不能互相替代。

在地方尺度上，这一事件应被看作资源和信息重新配置的过程：经费能否落实，取决于税源和中介；学校、警政或新军能否运行，取决于人员、土地和日常协作；家庭与社群对改革的经历，则可能因性别、职业、族属、居处与迁徙而不同。现存材料只照亮少数路径，叙事须保留被登记者与未被登记者之间的差距。

账本条目\texttt{EQ-1902-CIXI-RETURN-BEIJING}记录，光绪二十八年（1902年），清发生“慈禧与光绪帝返回北京，战后朝廷开始推出新政措施”。账本将其背景说明为辛丑议和结束流亡状态，可见结果为新政既回应危机，也服务朝廷重建合法性。这一节点虽可放入改革、外交或革命的宏观年表，却不能脱离省份财政、地方官、商会、学校、军队、家庭与交通网络来理解；同一上谕、条约或战事在不同地域的实际后果并不一致。

核心材料为《清德宗实录》光绪二十八年相关条；宫中朱批奏折、内阁大库题本与《清史稿》相关志传。它们可以支持年序、制度文本、机构设立或某些地点的行动，却不能自动代表全体居民的收入、教育、身份、性别关系或政治意愿。账本提示“以光绪二十八年（1902年）的同期文书为定位起点；官修记录、奏报或外交文本服务特定机构立场，具体日期、规模、地方执行与对手视角须以独立材料复核。”。因此，官方统计、外交措辞、报刊论战、教会叙述和军报数字必须分别置于生成它们的立场与制度目的之中，不能互相替代。

在地方尺度上，这一事件应被看作资源和信息重新配置的过程：经费能否落实，取决于税源和中介；学校、警政或新军能否运行，取决于人员、土地和日常协作；家庭与社群对改革的经历，则可能因性别、职业、族属、居处与迁徙而不同。现存材料只照亮少数路径，叙事须保留被登记者与未被登记者之间的差距。

账本条目\texttt{EQ-1902-CIXI-RETURN-BEIJING-DOCUMENTS}记录，光绪二十八年（1902年），清发生“围绕「慈禧与光绪帝返回北京，战后朝廷开始推出新政措施」的上谕、奏折与部议在中央—地方行政链中流转”。账本将其背景说明为核心事件「慈禧与光绪帝返回北京，战后朝廷开始推出新政措施」要求中央或地方通过文书处理资源、信息或责任。，可见结果为该记录为后续按具体谕旨、奏折、条约、报刊或地方档案拆分事件提供可回查入口。。这一节点虽可放入改革、外交或革命的宏观年表，却不能脱离省份财政、地方官、商会、学校、军队、家庭与交通网络来理解；同一上谕、条约或战事在不同地域的实际后果并不一致。

核心材料为《清德宗实录》光绪二十八年相关条；宫中朱批奏折、内阁大库题本与《清史稿》相关志传。它们可以支持年序、制度文本、机构设立或某些地点的行动，却不能自动代表全体居民的收入、教育、身份、性别关系或政治意愿。账本提示“此行仅确认相关文书链和可追查的行政动作；不得由其直接推断政策有效、地方服从、民众态度或单一因果。”。因此，官方统计、外交措辞、报刊论战、教会叙述和军报数字必须分别置于生成它们的立场与制度目的之中，不能互相替代。

在地方尺度上，这一事件应被看作资源和信息重新配置的过程：经费能否落实，取决于税源和中介；学校、警政或新军能否运行，取决于人员、土地和日常协作；家庭与社群对改革的经历，则可能因性别、职业、族属、居处与迁徙而不同。现存材料只照亮少数路径，叙事须保留被登记者与未被登记者之间的差距。

账本条目\texttt{EQ-1902-NEW-POLICIES-EDUCATION}记录，光绪二十八年（1902年），清发生“清末新政在学制、军制、商政和地方行政方面逐步启动”。账本将其背景说明为战败与财政外交压力显示旧制难以维持，可见结果为改革法令、试办与地方成效必须分别记录。这一节点虽可放入改革、外交或革命的宏观年表，却不能脱离省份财政、地方官、商会、学校、军队、家庭与交通网络来理解；同一上谕、条约或战事在不同地域的实际后果并不一致。

核心材料为《清德宗实录》光绪二十八年相关条；内阁档、文集或报刊相关记录。它们可以支持年序、制度文本、机构设立或某些地点的行动，却不能自动代表全体居民的收入、教育、身份、性别关系或政治意愿。账本提示“以光绪二十八年（1902年）的同期文书为定位起点；官修记录、奏报或外交文本服务特定机构立场，具体日期、规模、地方执行与对手视角须以独立材料复核。”。因此，官方统计、外交措辞、报刊论战、教会叙述和军报数字必须分别置于生成它们的立场与制度目的之中，不能互相替代。

在地方尺度上，这一事件应被看作资源和信息重新配置的过程：经费能否落实，取决于税源和中介；学校、警政或新军能否运行，取决于人员、土地和日常协作；家庭与社群对改革的经历，则可能因性别、职业、族属、居处与迁徙而不同。现存材料只照亮少数路径，叙事须保留被登记者与未被登记者之间的差距。

账本条目\texttt{EQ-1902-NEW-POLICIES-EDUCATION-DOCUMENTS}记录，光绪二十八年（1902年），清发生“围绕「清末新政在学制、军制、商政和地方行政方面逐步启动」的禁令、编纂、教育或传播活动留下中央和地方文本记录”。账本将其背景说明为核心事件「清末新政在学制、军制、商政和地方行政方面逐步启动」要求中央或地方通过文书处理资源、信息或责任。，可见结果为该记录为后续按具体谕旨、奏折、条约、报刊或地方档案拆分事件提供可回查入口。。这一节点虽可放入改革、外交或革命的宏观年表，却不能脱离省份财政、地方官、商会、学校、军队、家庭与交通网络来理解；同一上谕、条约或战事在不同地域的实际后果并不一致。

核心材料为《清德宗实录》光绪二十八年相关条；内阁档、文集或报刊相关记录。它们可以支持年序、制度文本、机构设立或某些地点的行动，却不能自动代表全体居民的收入、教育、身份、性别关系或政治意愿。账本提示“此行仅确认相关文书链和可追查的行政动作；不得由其直接推断政策有效、地方服从、民众态度或单一因果。”。因此，官方统计、外交措辞、报刊论战、教会叙述和军报数字必须分别置于生成它们的立场与制度目的之中，不能互相替代。

在地方尺度上，这一事件应被看作资源和信息重新配置的过程：经费能否落实，取决于税源和中介；学校、警政或新军能否运行，取决于人员、土地和日常协作；家庭与社群对改革的经历，则可能因性别、职业、族属、居处与迁徙而不同。现存材料只照亮少数路径，叙事须保留被登记者与未被登记者之间的差距。

账本条目\texttt{EQ-1903-BANNER-POSTS-REFORM}记录，光绪二十九年（1903年），清八旗发生“清廷废除部分旗人官职垄断，调整旗籍制度与职务资格”。账本将其背景说明为八旗财政危机与官僚改革需求上升，可见结果为制度变化对不同旗籍人群的影响不均。这一节点虽可放入改革、外交或革命的宏观年表，却不能脱离省份财政、地方官、商会、学校、军队、家庭与交通网络来理解；同一上谕、条约或战事在不同地域的实际后果并不一致。

核心材料为《清德宗实录》光绪二十九年相关条；地方志、督抚奏折及地方档案。它们可以支持年序、制度文本、机构设立或某些地点的行动，却不能自动代表全体居民的收入、教育、身份、性别关系或政治意愿。账本提示“以光绪二十九年（1903年）的同期文书为定位起点；官修记录、奏报或外交文本服务特定机构立场，具体日期、规模、地方执行与对手视角须以独立材料复核。”。因此，官方统计、外交措辞、报刊论战、教会叙述和军报数字必须分别置于生成它们的立场与制度目的之中，不能互相替代。

在地方尺度上，这一事件应被看作资源和信息重新配置的过程：经费能否落实，取决于税源和中介；学校、警政或新军能否运行，取决于人员、土地和日常协作；家庭与社群对改革的经历，则可能因性别、职业、族属、居处与迁徙而不同。现存材料只照亮少数路径，叙事须保留被登记者与未被登记者之间的差距。

账本条目\texttt{EQ-1903-BANNER-POSTS-REFORM-DOCUMENTS}记录，光绪二十九年（1903年），清八旗发生“围绕「清廷废除部分旗人官职垄断，调整旗籍制度与职务资格」的地方呈报、案件或赈务记录将社会反应带入官府文书链”。账本将其背景说明为核心事件「清廷废除部分旗人官职垄断，调整旗籍制度与职务资格」要求中央或地方通过文书处理资源、信息或责任。，可见结果为该记录为后续按具体谕旨、奏折、条约、报刊或地方档案拆分事件提供可回查入口。。这一节点虽可放入改革、外交或革命的宏观年表，却不能脱离省份财政、地方官、商会、学校、军队、家庭与交通网络来理解；同一上谕、条约或战事在不同地域的实际后果并不一致。

核心材料为《清德宗实录》光绪二十九年相关条；地方志、督抚奏折及地方档案。它们可以支持年序、制度文本、机构设立或某些地点的行动，却不能自动代表全体居民的收入、教育、身份、性别关系或政治意愿。账本提示“此行仅确认相关文书链和可追查的行政动作；不得由其直接推断政策有效、地方服从、民众态度或单一因果。”。因此，官方统计、外交措辞、报刊论战、教会叙述和军报数字必须分别置于生成它们的立场与制度目的之中，不能互相替代。

在地方尺度上，这一事件应被看作资源和信息重新配置的过程：经费能否落实，取决于税源和中介；学校、警政或新军能否运行，取决于人员、土地和日常协作；家庭与社群对改革的经历，则可能因性别、职业、族属、居处与迁徙而不同。现存材料只照亮少数路径，叙事须保留被登记者与未被登记者之间的差距。

账本条目\texttt{EQ-1903-SHANGHAI-COMMERCIAL-LAW}记录，光绪二十九年（1903年），清与通商口岸发生“商法、公司与工商管理的近代制度讨论在口岸和中央之间推进”。账本将其背景说明为国内商业、外贸和新式企业需要新的法律语言，可见结果为法典公布不等于商事审判和公司治理普遍落实。这一节点虽可放入改革、外交或革命的宏观年表，却不能脱离省份财政、地方官、商会、学校、军队、家庭与交通网络来理解；同一上谕、条约或战事在不同地域的实际后果并不一致。

核心材料为《清德宗实录》光绪二十九年相关条；户部奏销、盐政、河工或海关档案。它们可以支持年序、制度文本、机构设立或某些地点的行动，却不能自动代表全体居民的收入、教育、身份、性别关系或政治意愿。账本提示“以光绪二十九年（1903年）的同期文书为定位起点；官修记录、奏报或外交文本服务特定机构立场，具体日期、规模、地方执行与对手视角须以独立材料复核。”。因此，官方统计、外交措辞、报刊论战、教会叙述和军报数字必须分别置于生成它们的立场与制度目的之中，不能互相替代。

在地方尺度上，这一事件应被看作资源和信息重新配置的过程：经费能否落实，取决于税源和中介；学校、警政或新军能否运行，取决于人员、土地和日常协作；家庭与社群对改革的经历，则可能因性别、职业、族属、居处与迁徙而不同。现存材料只照亮少数路径，叙事须保留被登记者与未被登记者之间的差距。

账本条目\texttt{EQ-1903-SHANGHAI-COMMERCIAL-LAW-DOCUMENTS}记录，光绪二十九年（1903年），清与通商口岸发生“围绕「商法、公司与工商管理的近代制度讨论在口岸和中央之间推进」的经费、税收、赈济或工程呈报经由户部与地方衙门处理”。账本将其背景说明为核心事件「商法、公司与工商管理的近代制度讨论在口岸和中央之间推进」要求中央或地方通过文书处理资源、信息或责任。，可见结果为该记录为后续按具体谕旨、奏折、条约、报刊或地方档案拆分事件提供可回查入口。。这一节点虽可放入改革、外交或革命的宏观年表，却不能脱离省份财政、地方官、商会、学校、军队、家庭与交通网络来理解；同一上谕、条约或战事在不同地域的实际后果并不一致。

核心材料为《清德宗实录》光绪二十九年相关条；户部奏销、盐政、河工或海关档案。它们可以支持年序、制度文本、机构设立或某些地点的行动，却不能自动代表全体居民的收入、教育、身份、性别关系或政治意愿。账本提示“此行仅确认相关文书链和可追查的行政动作；不得由其直接推断政策有效、地方服从、民众态度或单一因果。”。因此，官方统计、外交措辞、报刊论战、教会叙述和军报数字必须分别置于生成它们的立场与制度目的之中，不能互相替代。

在地方尺度上，这一事件应被看作资源和信息重新配置的过程：经费能否落实，取决于税源和中介；学校、警政或新军能否运行，取决于人员、土地和日常协作；家庭与社群对改革的经历，则可能因性别、职业、族属、居处与迁徙而不同。现存材料只照亮少数路径，叙事须保留被登记者与未被登记者之间的差距。

账本条目\texttt{EQ-1904-RUSSO-JAPANESE-WAR}记录，光绪三十年（1904年），俄国、日本与清东北发生“日俄战争在中国东北爆发，清廷名义中立而东北遭受战场化”。账本将其背景说明为俄日争夺满洲铁路、港口和势力范围，可见结果为中国地方社会的损失与清廷的有限行动不能被列强叙事遮蔽。这一节点虽可放入改革、外交或革命的宏观年表，却不能脱离省份财政、地方官、商会、学校、军队、家庭与交通网络来理解；同一上谕、条约或战事在不同地域的实际后果并不一致。

核心材料为《清德宗实录》光绪三十年相关条；军机处档、战事方略及地方奏报。它们可以支持年序、制度文本、机构设立或某些地点的行动，却不能自动代表全体居民的收入、教育、身份、性别关系或政治意愿。账本提示“以光绪三十年（1904年）的同期文书为定位起点；官修记录、奏报或外交文本服务特定机构立场，具体日期、规模、地方执行与对手视角须以独立材料复核。”。因此，官方统计、外交措辞、报刊论战、教会叙述和军报数字必须分别置于生成它们的立场与制度目的之中，不能互相替代。

在地方尺度上，这一事件应被看作资源和信息重新配置的过程：经费能否落实，取决于税源和中介；学校、警政或新军能否运行，取决于人员、土地和日常协作；家庭与社群对改革的经历，则可能因性别、职业、族属、居处与迁徙而不同。现存材料只照亮少数路径，叙事须保留被登记者与未被登记者之间的差距。

账本条目\texttt{EQ-1904-RUSSO-JAPANESE-WAR-DOCUMENTS}记录，光绪三十年（1904年），俄国、日本与清东北发生“清廷围绕「日俄战争在中国东北爆发，清廷名义中立而东北遭受战场化」调遣军队、筹措饷械并处理战报，中央与地方的军政文书形成连续记录”。账本将其背景说明为核心事件「日俄战争在中国东北爆发，清廷名义中立而东北遭受战场化」要求中央或地方通过文书处理资源、信息或责任。，可见结果为该记录为后续按具体谕旨、奏折、条约、报刊或地方档案拆分事件提供可回查入口。。这一节点虽可放入改革、外交或革命的宏观年表，却不能脱离省份财政、地方官、商会、学校、军队、家庭与交通网络来理解；同一上谕、条约或战事在不同地域的实际后果并不一致。

核心材料为《清德宗实录》光绪三十年相关条；军机处档、战事方略及地方奏报。它们可以支持年序、制度文本、机构设立或某些地点的行动，却不能自动代表全体居民的收入、教育、身份、性别关系或政治意愿。账本提示“此行仅确认相关文书链和可追查的行政动作；不得由其直接推断政策有效、地方服从、民众态度或单一因果。”。因此，官方统计、外交措辞、报刊论战、教会叙述和军报数字必须分别置于生成它们的立场与制度目的之中，不能互相替代。

在地方尺度上，这一事件应被看作资源和信息重新配置的过程：经费能否落实，取决于税源和中介；学校、警政或新军能否运行，取决于人员、土地和日常协作；家庭与社群对改革的经历，则可能因性别、职业、族属、居处与迁徙而不同。现存材料只照亮少数路径，叙事须保留被登记者与未被登记者之间的差距。

账本条目\texttt{EQ-1904-IMPERIAL-EXAMINATION-DEBATE}记录，光绪三十年（1904年），清发生“科举存废与新式教育的争论加速”。账本将其背景说明为军政危机使人才选拔制度受到质疑，可见结果为讨论、诏令和实际学堂替代之间存在时间差。这一节点虽可放入改革、外交或革命的宏观年表，却不能脱离省份财政、地方官、商会、学校、军队、家庭与交通网络来理解；同一上谕、条约或战事在不同地域的实际后果并不一致。

核心材料为《清德宗实录》光绪三十年相关条；内阁档、文集或报刊相关记录。它们可以支持年序、制度文本、机构设立或某些地点的行动，却不能自动代表全体居民的收入、教育、身份、性别关系或政治意愿。账本提示“以光绪三十年（1904年）的同期文书为定位起点；官修记录、奏报或外交文本服务特定机构立场，具体日期、规模、地方执行与对手视角须以独立材料复核。”。因此，官方统计、外交措辞、报刊论战、教会叙述和军报数字必须分别置于生成它们的立场与制度目的之中，不能互相替代。

在地方尺度上，这一事件应被看作资源和信息重新配置的过程：经费能否落实，取决于税源和中介；学校、警政或新军能否运行，取决于人员、土地和日常协作；家庭与社群对改革的经历，则可能因性别、职业、族属、居处与迁徙而不同。现存材料只照亮少数路径，叙事须保留被登记者与未被登记者之间的差距。

账本条目\texttt{EQ-1904-IMPERIAL-EXAMINATION-DEBATE-DOCUMENTS}记录，光绪三十年（1904年），清发生“围绕「科举存废与新式教育的争论加速」的禁令、编纂、教育或传播活动留下中央和地方文本记录”。账本将其背景说明为核心事件「科举存废与新式教育的争论加速」要求中央或地方通过文书处理资源、信息或责任。，可见结果为该记录为后续按具体谕旨、奏折、条约、报刊或地方档案拆分事件提供可回查入口。。这一节点虽可放入改革、外交或革命的宏观年表，却不能脱离省份财政、地方官、商会、学校、军队、家庭与交通网络来理解；同一上谕、条约或战事在不同地域的实际后果并不一致。

核心材料为《清德宗实录》光绪三十年相关条；内阁档、文集或报刊相关记录。它们可以支持年序、制度文本、机构设立或某些地点的行动，却不能自动代表全体居民的收入、教育、身份、性别关系或政治意愿。账本提示“此行仅确认相关文书链和可追查的行政动作；不得由其直接推断政策有效、地方服从、民众态度或单一因果。”。因此，官方统计、外交措辞、报刊论战、教会叙述和军报数字必须分别置于生成它们的立场与制度目的之中，不能互相替代。

在地方尺度上，这一事件应被看作资源和信息重新配置的过程：经费能否落实，取决于税源和中介；学校、警政或新军能否运行，取决于人员、土地和日常协作；家庭与社群对改革的经历，则可能因性别、职业、族属、居处与迁徙而不同。现存材料只照亮少数路径，叙事须保留被登记者与未被登记者之间的差距。

账本条目\texttt{EQ-1905-ABOLITION-OF-CIVIL-SERVICE-EXAM}记录，光绪三十一年（1905年），清发生“清廷废除科举，延续千余年的取士制度终结”。账本将其背景说明为新式教育、官僚改革与科举批评共同推动决策，可见结果为废科举不等于地方教育资源立即均衡或士人身份消失。这一节点虽可放入改革、外交或革命的宏观年表，却不能脱离省份财政、地方官、商会、学校、军队、家庭与交通网络来理解；同一上谕、条约或战事在不同地域的实际后果并不一致。

核心材料为《清德宗实录》光绪三十一年相关条；内阁档、文集或报刊相关记录。它们可以支持年序、制度文本、机构设立或某些地点的行动，却不能自动代表全体居民的收入、教育、身份、性别关系或政治意愿。账本提示“以光绪三十一年（1905年）的同期文书为定位起点；官修记录、奏报或外交文本服务特定机构立场，具体日期、规模、地方执行与对手视角须以独立材料复核。”。因此，官方统计、外交措辞、报刊论战、教会叙述和军报数字必须分别置于生成它们的立场与制度目的之中，不能互相替代。

在地方尺度上，这一事件应被看作资源和信息重新配置的过程：经费能否落实，取决于税源和中介；学校、警政或新军能否运行，取决于人员、土地和日常协作；家庭与社群对改革的经历，则可能因性别、职业、族属、居处与迁徙而不同。现存材料只照亮少数路径，叙事须保留被登记者与未被登记者之间的差距。

账本条目\texttt{EQ-1905-ABOLITION-OF-CIVIL-SERVICE-EXAM-DOCUMENTS}记录，光绪三十一年（1905年），清发生“围绕「清廷废除科举，延续千余年的取士制度终结」的禁令、编纂、教育或传播活动留下中央和地方文本记录”。账本将其背景说明为核心事件「清廷废除科举，延续千余年的取士制度终结」要求中央或地方通过文书处理资源、信息或责任。，可见结果为该记录为后续按具体谕旨、奏折、条约、报刊或地方档案拆分事件提供可回查入口。。这一节点虽可放入改革、外交或革命的宏观年表，却不能脱离省份财政、地方官、商会、学校、军队、家庭与交通网络来理解；同一上谕、条约或战事在不同地域的实际后果并不一致。

核心材料为《清德宗实录》光绪三十一年相关条；内阁档、文集或报刊相关记录。它们可以支持年序、制度文本、机构设立或某些地点的行动，却不能自动代表全体居民的收入、教育、身份、性别关系或政治意愿。账本提示“此行仅确认相关文书链和可追查的行政动作；不得由其直接推断政策有效、地方服从、民众态度或单一因果。”。因此，官方统计、外交措辞、报刊论战、教会叙述和军报数字必须分别置于生成它们的立场与制度目的之中，不能互相替代。

在地方尺度上，这一事件应被看作资源和信息重新配置的过程：经费能否落实，取决于税源和中介；学校、警政或新军能否运行，取决于人员、土地和日常协作；家庭与社群对改革的经历，则可能因性别、职业、族属、居处与迁徙而不同。现存材料只照亮少数路径，叙事须保留被登记者与未被登记者之间的差距。

账本条目\texttt{EQ-1905-TONGMENGHUI}记录，光绪三十一年（1905年），革命派与清发生“中国同盟会在东京成立，革命组织和宣传网络进一步整合”。账本将其背景说明为留日学生、海外华侨和反清团体形成跨境联系，可见结果为组织自称、成员规模与地方影响需要分别核验。这一节点虽可放入改革、外交或革命的宏观年表，却不能脱离省份财政、地方官、商会、学校、军队、家庭与交通网络来理解；同一上谕、条约或战事在不同地域的实际后果并不一致。

核心材料为《清德宗实录》光绪三十一年相关条；地方志、督抚奏折及地方档案。它们可以支持年序、制度文本、机构设立或某些地点的行动，却不能自动代表全体居民的收入、教育、身份、性别关系或政治意愿。账本提示“以光绪三十一年（1905年）的同期文书为定位起点；官修记录、奏报或外交文本服务特定机构立场，具体日期、规模、地方执行与对手视角须以独立材料复核。”。因此，官方统计、外交措辞、报刊论战、教会叙述和军报数字必须分别置于生成它们的立场与制度目的之中，不能互相替代。

在地方尺度上，这一事件应被看作资源和信息重新配置的过程：经费能否落实，取决于税源和中介；学校、警政或新军能否运行，取决于人员、土地和日常协作；家庭与社群对改革的经历，则可能因性别、职业、族属、居处与迁徙而不同。现存材料只照亮少数路径，叙事须保留被登记者与未被登记者之间的差距。

账本条目\texttt{EQ-1905-TONGMENGHUI-DOCUMENTS}记录，光绪三十一年（1905年），革命派与清发生“围绕「中国同盟会在东京成立，革命组织和宣传网络进一步整合」的地方呈报、案件或赈务记录将社会反应带入官府文书链”。账本将其背景说明为核心事件「中国同盟会在东京成立，革命组织和宣传网络进一步整合」要求中央或地方通过文书处理资源、信息或责任。，可见结果为该记录为后续按具体谕旨、奏折、条约、报刊或地方档案拆分事件提供可回查入口。。这一节点虽可放入改革、外交或革命的宏观年表，却不能脱离省份财政、地方官、商会、学校、军队、家庭与交通网络来理解；同一上谕、条约或战事在不同地域的实际后果并不一致。

核心材料为《清德宗实录》光绪三十一年相关条；地方志、督抚奏折及地方档案。它们可以支持年序、制度文本、机构设立或某些地点的行动，却不能自动代表全体居民的收入、教育、身份、性别关系或政治意愿。账本提示“此行仅确认相关文书链和可追查的行政动作；不得由其直接推断政策有效、地方服从、民众态度或单一因果。”。因此，官方统计、外交措辞、报刊论战、教会叙述和军报数字必须分别置于生成它们的立场与制度目的之中，不能互相替代。

在地方尺度上，这一事件应被看作资源和信息重新配置的过程：经费能否落实，取决于税源和中介；学校、警政或新军能否运行，取决于人员、土地和日常协作；家庭与社群对改革的经历，则可能因性别、职业、族属、居处与迁徙而不同。现存材料只照亮少数路径，叙事须保留被登记者与未被登记者之间的差距。

账本条目\texttt{EQ-1906-CONSTITUTIONAL-PROMISE}记录，光绪三十二年（1906年），清发生“清廷宣布预备立宪，着手考察各国制度和改革官制”。账本将其背景说明为新政需要重建统治合法性并回应立宪压力，可见结果为“预备”时间表和实际权力让渡之间存在重大距离。这一节点虽可放入改革、外交或革命的宏观年表，却不能脱离省份财政、地方官、商会、学校、军队、家庭与交通网络来理解；同一上谕、条约或战事在不同地域的实际后果并不一致。

核心材料为《清德宗实录》光绪三十二年相关条；宫中朱批奏折、内阁大库题本与《清史稿》相关志传。它们可以支持年序、制度文本、机构设立或某些地点的行动，却不能自动代表全体居民的收入、教育、身份、性别关系或政治意愿。账本提示“以光绪三十二年（1906年）的同期文书为定位起点；官修记录、奏报或外交文本服务特定机构立场，具体日期、规模、地方执行与对手视角须以独立材料复核。”。因此，官方统计、外交措辞、报刊论战、教会叙述和军报数字必须分别置于生成它们的立场与制度目的之中，不能互相替代。

在地方尺度上，这一事件应被看作资源和信息重新配置的过程：经费能否落实，取决于税源和中介；学校、警政或新军能否运行，取决于人员、土地和日常协作；家庭与社群对改革的经历，则可能因性别、职业、族属、居处与迁徙而不同。现存材料只照亮少数路径，叙事须保留被登记者与未被登记者之间的差距。

账本条目\texttt{EQ-1906-CONSTITUTIONAL-PROMISE-DOCUMENTS}记录，光绪三十二年（1906年），清发生“围绕「清廷宣布预备立宪，着手考察各国制度和改革官制」的上谕、奏折与部议在中央—地方行政链中流转”。账本将其背景说明为核心事件「清廷宣布预备立宪，着手考察各国制度和改革官制」要求中央或地方通过文书处理资源、信息或责任。，可见结果为该记录为后续按具体谕旨、奏折、条约、报刊或地方档案拆分事件提供可回查入口。。这一节点虽可放入改革、外交或革命的宏观年表，却不能脱离省份财政、地方官、商会、学校、军队、家庭与交通网络来理解；同一上谕、条约或战事在不同地域的实际后果并不一致。

核心材料为《清德宗实录》光绪三十二年相关条；宫中朱批奏折、内阁大库题本与《清史稿》相关志传。它们可以支持年序、制度文本、机构设立或某些地点的行动，却不能自动代表全体居民的收入、教育、身份、性别关系或政治意愿。账本提示“此行仅确认相关文书链和可追查的行政动作；不得由其直接推断政策有效、地方服从、民众态度或单一因果。”。因此，官方统计、外交措辞、报刊论战、教会叙述和军报数字必须分别置于生成它们的立场与制度目的之中，不能互相替代。

在地方尺度上，这一事件应被看作资源和信息重新配置的过程：经费能否落实，取决于税源和中介；学校、警政或新军能否运行，取决于人员、土地和日常协作；家庭与社群对改革的经历，则可能因性别、职业、族属、居处与迁徙而不同。现存材料只照亮少数路径，叙事须保留被登记者与未被登记者之间的差距。

账本条目\texttt{EQ-1906-NEW-ARMY-REFORM}记录，光绪三十二年（1906年），清发生“新军编练与军事学堂扩展，军队成为新式政治动员的重要空间”。账本将其背景说明为甲午、庚子和日俄战争显示旧军制弱点，可见结果为编制、装备和忠诚网络在各省差异很大。这一节点虽可放入改革、外交或革命的宏观年表，却不能脱离省份财政、地方官、商会、学校、军队、家庭与交通网络来理解；同一上谕、条约或战事在不同地域的实际后果并不一致。

核心材料为《清德宗实录》光绪三十二年相关条；军机处档、战事方略及地方奏报。它们可以支持年序、制度文本、机构设立或某些地点的行动，却不能自动代表全体居民的收入、教育、身份、性别关系或政治意愿。账本提示“以光绪三十二年（1906年）的同期文书为定位起点；官修记录、奏报或外交文本服务特定机构立场，具体日期、规模、地方执行与对手视角须以独立材料复核。”。因此，官方统计、外交措辞、报刊论战、教会叙述和军报数字必须分别置于生成它们的立场与制度目的之中，不能互相替代。

在地方尺度上，这一事件应被看作资源和信息重新配置的过程：经费能否落实，取决于税源和中介；学校、警政或新军能否运行，取决于人员、土地和日常协作；家庭与社群对改革的经历，则可能因性别、职业、族属、居处与迁徙而不同。现存材料只照亮少数路径，叙事须保留被登记者与未被登记者之间的差距。

账本条目\texttt{EQ-1906-NEW-ARMY-REFORM-DOCUMENTS}记录，光绪三十二年（1906年），清发生“清廷围绕「新军编练与军事学堂扩展，军队成为新式政治动员的重要空间」调遣军队、筹措饷械并处理战报，中央与地方的军政文书形成连续记录”。账本将其背景说明为核心事件「新军编练与军事学堂扩展，军队成为新式政治动员的重要空间」要求中央或地方通过文书处理资源、信息或责任。，可见结果为该记录为后续按具体谕旨、奏折、条约、报刊或地方档案拆分事件提供可回查入口。。这一节点虽可放入改革、外交或革命的宏观年表，却不能脱离省份财政、地方官、商会、学校、军队、家庭与交通网络来理解；同一上谕、条约或战事在不同地域的实际后果并不一致。

核心材料为《清德宗实录》光绪三十二年相关条；军机处档、战事方略及地方奏报。它们可以支持年序、制度文本、机构设立或某些地点的行动，却不能自动代表全体居民的收入、教育、身份、性别关系或政治意愿。账本提示“此行仅确认相关文书链和可追查的行政动作；不得由其直接推断政策有效、地方服从、民众态度或单一因果。”。因此，官方统计、外交措辞、报刊论战、教会叙述和军报数字必须分别置于生成它们的立场与制度目的之中，不能互相替代。

在地方尺度上，这一事件应被看作资源和信息重新配置的过程：经费能否落实，取决于税源和中介；学校、警政或新军能否运行，取决于人员、土地和日常协作；家庭与社群对改革的经历，则可能因性别、职业、族属、居处与迁徙而不同。现存材料只照亮少数路径，叙事须保留被登记者与未被登记者之间的差距。

账本条目\texttt{EQ-1907-MANCHURIAN-PROVINCES}记录，光绪三十三年（1907年），清与东北发生“清廷将东三省调整为行省体制，加强近代行政与边防经营”。账本将其背景说明为日俄战争后东北主权和铁路问题尖锐化，可见结果为省制设立与铁路沿线列强控制并存。这一节点虽可放入改革、外交或革命的宏观年表，却不能脱离省份财政、地方官、商会、学校、军队、家庭与交通网络来理解；同一上谕、条约或战事在不同地域的实际后果并不一致。

核心材料为《清德宗实录》光绪三十三年相关条；军机处档、理藩院档或相关方略。它们可以支持年序、制度文本、机构设立或某些地点的行动，却不能自动代表全体居民的收入、教育、身份、性别关系或政治意愿。账本提示“以光绪三十三年（1907年）的同期文书为定位起点；官修记录、奏报或外交文本服务特定机构立场，具体日期、规模、地方执行与对手视角须以独立材料复核。”。因此，官方统计、外交措辞、报刊论战、教会叙述和军报数字必须分别置于生成它们的立场与制度目的之中，不能互相替代。

在地方尺度上，这一事件应被看作资源和信息重新配置的过程：经费能否落实，取决于税源和中介；学校、警政或新军能否运行，取决于人员、土地和日常协作；家庭与社群对改革的经历，则可能因性别、职业、族属、居处与迁徙而不同。现存材料只照亮少数路径，叙事须保留被登记者与未被登记者之间的差距。

账本条目\texttt{EQ-1907-MANCHURIAN-PROVINCES-DOCUMENTS}记录，光绪三十三年（1907年），清与东北发生“清廷围绕「清廷将东三省调整为行省体制，加强近代行政与边防经营」处理盟旗、驻防、交通或地方呈报，相关边务文书进入中央决策链”。账本将其背景说明为核心事件「清廷将东三省调整为行省体制，加强近代行政与边防经营」要求中央或地方通过文书处理资源、信息或责任。，可见结果为该记录为后续按具体谕旨、奏折、条约、报刊或地方档案拆分事件提供可回查入口。。这一节点虽可放入改革、外交或革命的宏观年表，却不能脱离省份财政、地方官、商会、学校、军队、家庭与交通网络来理解；同一上谕、条约或战事在不同地域的实际后果并不一致。

核心材料为《清德宗实录》光绪三十三年相关条；军机处档、理藩院档或相关方略。它们可以支持年序、制度文本、机构设立或某些地点的行动，却不能自动代表全体居民的收入、教育、身份、性别关系或政治意愿。账本提示“此行仅确认相关文书链和可追查的行政动作；不得由其直接推断政策有效、地方服从、民众态度或单一因果。”。因此，官方统计、外交措辞、报刊论战、教会叙述和军报数字必须分别置于生成它们的立场与制度目的之中，不能互相替代。

在地方尺度上，这一事件应被看作资源和信息重新配置的过程：经费能否落实，取决于税源和中介；学校、警政或新军能否运行，取决于人员、土地和日常协作；家庭与社群对改革的经历，则可能因性别、职业、族属、居处与迁徙而不同。现存材料只照亮少数路径，叙事须保留被登记者与未被登记者之间的差距。

账本条目\texttt{EQ-1907-QING-LEGAL-REFORM}记录，光绪三十三年（1907年），清发生“清廷推进修律、审判和警政改革，尝试改造传统法律制度”。账本将其背景说明为领事裁判权和国内治理压力推动法制改革，可见结果为新法文本、试行和基层司法实践不能混同。这一节点虽可放入改革、外交或革命的宏观年表，却不能脱离省份财政、地方官、商会、学校、军队、家庭与交通网络来理解；同一上谕、条约或战事在不同地域的实际后果并不一致。

核心材料为《清德宗实录》光绪三十三年相关条；宫中朱批奏折、内阁大库题本与《清史稿》相关志传。它们可以支持年序、制度文本、机构设立或某些地点的行动，却不能自动代表全体居民的收入、教育、身份、性别关系或政治意愿。账本提示“以光绪三十三年（1907年）的同期文书为定位起点；官修记录、奏报或外交文本服务特定机构立场，具体日期、规模、地方执行与对手视角须以独立材料复核。”。因此，官方统计、外交措辞、报刊论战、教会叙述和军报数字必须分别置于生成它们的立场与制度目的之中，不能互相替代。

在地方尺度上，这一事件应被看作资源和信息重新配置的过程：经费能否落实，取决于税源和中介；学校、警政或新军能否运行，取决于人员、土地和日常协作；家庭与社群对改革的经历，则可能因性别、职业、族属、居处与迁徙而不同。现存材料只照亮少数路径，叙事须保留被登记者与未被登记者之间的差距。

账本条目\texttt{EQ-1907-QING-LEGAL-REFORM-DOCUMENTS}记录，光绪三十三年（1907年），清发生“围绕「清廷推进修律、审判和警政改革，尝试改造传统法律制度」的上谕、奏折与部议在中央—地方行政链中流转”。账本将其背景说明为核心事件「清廷推进修律、审判和警政改革，尝试改造传统法律制度」要求中央或地方通过文书处理资源、信息或责任。，可见结果为该记录为后续按具体谕旨、奏折、条约、报刊或地方档案拆分事件提供可回查入口。。这一节点虽可放入改革、外交或革命的宏观年表，却不能脱离省份财政、地方官、商会、学校、军队、家庭与交通网络来理解；同一上谕、条约或战事在不同地域的实际后果并不一致。

核心材料为《清德宗实录》光绪三十三年相关条；宫中朱批奏折、内阁大库题本与《清史稿》相关志传。它们可以支持年序、制度文本、机构设立或某些地点的行动，却不能自动代表全体居民的收入、教育、身份、性别关系或政治意愿。账本提示“此行仅确认相关文书链和可追查的行政动作；不得由其直接推断政策有效、地方服从、民众态度或单一因果。”。因此，官方统计、外交措辞、报刊论战、教会叙述和军报数字必须分别置于生成它们的立场与制度目的之中，不能互相替代。

在地方尺度上，这一事件应被看作资源和信息重新配置的过程：经费能否落实，取决于税源和中介；学校、警政或新军能否运行，取决于人员、土地和日常协作；家庭与社群对改革的经历，则可能因性别、职业、族属、居处与迁徙而不同。现存材料只照亮少数路径，叙事须保留被登记者与未被登记者之间的差距。

账本条目\texttt{EQ-1908-GUANGXU-CIXI-DEATHS}记录，光绪三十四年（1908年），清发生“光绪帝与慈禧太后相继去世，溥仪继位为宣统帝”。账本将其背景说明为晚清权力核心突然更替，可见结果为摄政王载沣与幼帝体制面临改革和危机双重压力。这一节点虽可放入改革、外交或革命的宏观年表，却不能脱离省份财政、地方官、商会、学校、军队、家庭与交通网络来理解；同一上谕、条约或战事在不同地域的实际后果并不一致。

核心材料为《清德宗实录》光绪三十四年相关条；宫中朱批奏折、内阁大库题本与《清史稿》相关志传。它们可以支持年序、制度文本、机构设立或某些地点的行动，却不能自动代表全体居民的收入、教育、身份、性别关系或政治意愿。账本提示“以光绪三十四年（1908年）的同期文书为定位起点；官修记录、奏报或外交文本服务特定机构立场，具体日期、规模、地方执行与对手视角须以独立材料复核。”。因此，官方统计、外交措辞、报刊论战、教会叙述和军报数字必须分别置于生成它们的立场与制度目的之中，不能互相替代。

在地方尺度上，这一事件应被看作资源和信息重新配置的过程：经费能否落实，取决于税源和中介；学校、警政或新军能否运行，取决于人员、土地和日常协作；家庭与社群对改革的经历，则可能因性别、职业、族属、居处与迁徙而不同。现存材料只照亮少数路径，叙事须保留被登记者与未被登记者之间的差距。

账本条目\texttt{EQ-1908-GUANGXU-CIXI-DEATHS-DOCUMENTS}记录，光绪三十四年（1908年），清发生“围绕「光绪帝与慈禧太后相继去世，溥仪继位为宣统帝」的上谕、奏折与部议在中央—地方行政链中流转”。账本将其背景说明为核心事件「光绪帝与慈禧太后相继去世，溥仪继位为宣统帝」要求中央或地方通过文书处理资源、信息或责任。，可见结果为该记录为后续按具体谕旨、奏折、条约、报刊或地方档案拆分事件提供可回查入口。。这一节点虽可放入改革、外交或革命的宏观年表，却不能脱离省份财政、地方官、商会、学校、军队、家庭与交通网络来理解；同一上谕、条约或战事在不同地域的实际后果并不一致。

核心材料为《清德宗实录》光绪三十四年相关条；宫中朱批奏折、内阁大库题本与《清史稿》相关志传。它们可以支持年序、制度文本、机构设立或某些地点的行动，却不能自动代表全体居民的收入、教育、身份、性别关系或政治意愿。账本提示“此行仅确认相关文书链和可追查的行政动作；不得由其直接推断政策有效、地方服从、民众态度或单一因果。”。因此，官方统计、外交措辞、报刊论战、教会叙述和军报数字必须分别置于生成它们的立场与制度目的之中，不能互相替代。

在地方尺度上，这一事件应被看作资源和信息重新配置的过程：经费能否落实，取决于税源和中介；学校、警政或新军能否运行，取决于人员、土地和日常协作；家庭与社群对改革的经历，则可能因性别、职业、族属、居处与迁徙而不同。现存材料只照亮少数路径，叙事须保留被登记者与未被登记者之间的差距。

账本条目\texttt{EQ-1908-CONSTITUTIONAL-OUTLINE}记录，光绪三十四年（1908年），清发生“清廷公布《钦定宪法大纲》等文件，明确预备立宪框架”。账本将其背景说明为立宪压力与中央保留权力的意图并存，可见结果为文本承诺与咨议局、资政院实际权限存在落差。这一节点虽可放入改革、外交或革命的宏观年表，却不能脱离省份财政、地方官、商会、学校、军队、家庭与交通网络来理解；同一上谕、条约或战事在不同地域的实际后果并不一致。

核心材料为《清德宗实录》光绪三十四年相关条；宫中朱批奏折、内阁大库题本与《清史稿》相关志传。它们可以支持年序、制度文本、机构设立或某些地点的行动，却不能自动代表全体居民的收入、教育、身份、性别关系或政治意愿。账本提示“以光绪三十四年（1908年）的同期文书为定位起点；官修记录、奏报或外交文本服务特定机构立场，具体日期、规模、地方执行与对手视角须以独立材料复核。”。因此，官方统计、外交措辞、报刊论战、教会叙述和军报数字必须分别置于生成它们的立场与制度目的之中，不能互相替代。

在地方尺度上，这一事件应被看作资源和信息重新配置的过程：经费能否落实，取决于税源和中介；学校、警政或新军能否运行，取决于人员、土地和日常协作；家庭与社群对改革的经历，则可能因性别、职业、族属、居处与迁徙而不同。现存材料只照亮少数路径，叙事须保留被登记者与未被登记者之间的差距。

账本条目\texttt{EQ-1908-CONSTITUTIONAL-OUTLINE-DOCUMENTS}记录，光绪三十四年（1908年），清发生“围绕「清廷公布《钦定宪法大纲》等文件，明确预备立宪框架」的上谕、奏折与部议在中央—地方行政链中流转”。账本将其背景说明为核心事件「清廷公布《钦定宪法大纲》等文件，明确预备立宪框架」要求中央或地方通过文书处理资源、信息或责任。，可见结果为该记录为后续按具体谕旨、奏折、条约、报刊或地方档案拆分事件提供可回查入口。。这一节点虽可放入改革、外交或革命的宏观年表，却不能脱离省份财政、地方官、商会、学校、军队、家庭与交通网络来理解；同一上谕、条约或战事在不同地域的实际后果并不一致。

核心材料为《清德宗实录》光绪三十四年相关条；宫中朱批奏折、内阁大库题本与《清史稿》相关志传。它们可以支持年序、制度文本、机构设立或某些地点的行动，却不能自动代表全体居民的收入、教育、身份、性别关系或政治意愿。账本提示“此行仅确认相关文书链和可追查的行政动作；不得由其直接推断政策有效、地方服从、民众态度或单一因果。”。因此，官方统计、外交措辞、报刊论战、教会叙述和军报数字必须分别置于生成它们的立场与制度目的之中，不能互相替代。

在地方尺度上，这一事件应被看作资源和信息重新配置的过程：经费能否落实，取决于税源和中介；学校、警政或新军能否运行，取决于人员、土地和日常协作；家庭与社群对改革的经历，则可能因性别、职业、族属、居处与迁徙而不同。现存材料只照亮少数路径，叙事须保留被登记者与未被登记者之间的差距。

账本条目\texttt{EQ-1909-PROVINCIAL-ASSEMBLIES}记录，宣统元年（1909年），清与各省发生“各省咨议局开办，地方精英获得新的公开政治场域”。账本将其背景说明为预备立宪需要有限度吸纳地方意见，可见结果为选举资格、代表性和实际权力均受限制。这一节点虽可放入改革、外交或革命的宏观年表，却不能脱离省份财政、地方官、商会、学校、军队、家庭与交通网络来理解；同一上谕、条约或战事在不同地域的实际后果并不一致。

核心材料为《宣统政纪》宣统元年相关条；宫中朱批奏折、内阁大库题本与《清史稿》相关志传。它们可以支持年序、制度文本、机构设立或某些地点的行动，却不能自动代表全体居民的收入、教育、身份、性别关系或政治意愿。账本提示“以宣统元年（1909年）的同期文书为定位起点；官修记录、奏报或外交文本服务特定机构立场，具体日期、规模、地方执行与对手视角须以独立材料复核。”。因此，官方统计、外交措辞、报刊论战、教会叙述和军报数字必须分别置于生成它们的立场与制度目的之中，不能互相替代。

在地方尺度上，这一事件应被看作资源和信息重新配置的过程：经费能否落实，取决于税源和中介；学校、警政或新军能否运行，取决于人员、土地和日常协作；家庭与社群对改革的经历，则可能因性别、职业、族属、居处与迁徙而不同。现存材料只照亮少数路径，叙事须保留被登记者与未被登记者之间的差距。

账本条目\texttt{EQ-1909-PROVINCIAL-ASSEMBLIES-DOCUMENTS}记录，宣统元年（1909年），清与各省发生“围绕「各省咨议局开办，地方精英获得新的公开政治场域」的上谕、奏折与部议在中央—地方行政链中流转”。账本将其背景说明为核心事件「各省咨议局开办，地方精英获得新的公开政治场域」要求中央或地方通过文书处理资源、信息或责任。，可见结果为该记录为后续按具体谕旨、奏折、条约、报刊或地方档案拆分事件提供可回查入口。。这一节点虽可放入改革、外交或革命的宏观年表，却不能脱离省份财政、地方官、商会、学校、军队、家庭与交通网络来理解；同一上谕、条约或战事在不同地域的实际后果并不一致。

核心材料为《宣统政纪》宣统元年相关条；宫中朱批奏折、内阁大库题本与《清史稿》相关志传。它们可以支持年序、制度文本、机构设立或某些地点的行动，却不能自动代表全体居民的收入、教育、身份、性别关系或政治意愿。账本提示“此行仅确认相关文书链和可追查的行政动作；不得由其直接推断政策有效、地方服从、民众态度或单一因果。”。因此，官方统计、外交措辞、报刊论战、教会叙述和军报数字必须分别置于生成它们的立场与制度目的之中，不能互相替代。

在地方尺度上，这一事件应被看作资源和信息重新配置的过程：经费能否落实，取决于税源和中介；学校、警政或新军能否运行，取决于人员、土地和日常协作；家庭与社群对改革的经历，则可能因性别、职业、族属、居处与迁徙而不同。现存材料只照亮少数路径，叙事须保留被登记者与未被登记者之间的差距。

账本条目\texttt{EQ-1909-RAILWAY-RECOVERY-MOVEMENT}记录，宣统元年（1909年），清与各省商绅发生“铁路国有化和路权问题引发各省商绅、公司和官府的争论”。账本将其背景说明为外债、地方集资与中央财政需求相冲突，可见结果为“收回路权”诉求包含不同地区和利益群体。这一节点虽可放入改革、外交或革命的宏观年表，却不能脱离省份财政、地方官、商会、学校、军队、家庭与交通网络来理解；同一上谕、条约或战事在不同地域的实际后果并不一致。

核心材料为《宣统政纪》宣统元年相关条；户部奏销、盐政、河工或海关档案。它们可以支持年序、制度文本、机构设立或某些地点的行动，却不能自动代表全体居民的收入、教育、身份、性别关系或政治意愿。账本提示“以宣统元年（1909年）的同期文书为定位起点；官修记录、奏报或外交文本服务特定机构立场，具体日期、规模、地方执行与对手视角须以独立材料复核。”。因此，官方统计、外交措辞、报刊论战、教会叙述和军报数字必须分别置于生成它们的立场与制度目的之中，不能互相替代。

在地方尺度上，这一事件应被看作资源和信息重新配置的过程：经费能否落实，取决于税源和中介；学校、警政或新军能否运行，取决于人员、土地和日常协作；家庭与社群对改革的经历，则可能因性别、职业、族属、居处与迁徙而不同。现存材料只照亮少数路径，叙事须保留被登记者与未被登记者之间的差距。

账本条目\texttt{EQ-1909-RAILWAY-RECOVERY-MOVEMENT-DOCUMENTS}记录，宣统元年（1909年），清与各省商绅发生“围绕「铁路国有化和路权问题引发各省商绅、公司和官府的争论」的经费、税收、赈济或工程呈报经由户部与地方衙门处理”。账本将其背景说明为核心事件「铁路国有化和路权问题引发各省商绅、公司和官府的争论」要求中央或地方通过文书处理资源、信息或责任。，可见结果为该记录为后续按具体谕旨、奏折、条约、报刊或地方档案拆分事件提供可回查入口。。这一节点虽可放入改革、外交或革命的宏观年表，却不能脱离省份财政、地方官、商会、学校、军队、家庭与交通网络来理解；同一上谕、条约或战事在不同地域的实际后果并不一致。

核心材料为《宣统政纪》宣统元年相关条；户部奏销、盐政、河工或海关档案。它们可以支持年序、制度文本、机构设立或某些地点的行动，却不能自动代表全体居民的收入、教育、身份、性别关系或政治意愿。账本提示“此行仅确认相关文书链和可追查的行政动作；不得由其直接推断政策有效、地方服从、民众态度或单一因果。”。因此，官方统计、外交措辞、报刊论战、教会叙述和军报数字必须分别置于生成它们的立场与制度目的之中，不能互相替代。

在地方尺度上，这一事件应被看作资源和信息重新配置的过程：经费能否落实，取决于税源和中介；学校、警政或新军能否运行，取决于人员、土地和日常协作；家庭与社群对改革的经历，则可能因性别、职业、族属、居处与迁徙而不同。现存材料只照亮少数路径，叙事须保留被登记者与未被登记者之间的差距。

账本条目\texttt{EQ-1910-ZIZHENG-YUAN}记录，宣统二年（1910年），清发生“资政院开院，中央层面的立宪咨询机构开始运作”。账本将其背景说明为朝廷试图展示宪政进程并控制代表参与，可见结果为咨询权与立法权的界限成为持续争议。这一节点虽可放入改革、外交或革命的宏观年表，却不能脱离省份财政、地方官、商会、学校、军队、家庭与交通网络来理解；同一上谕、条约或战事在不同地域的实际后果并不一致。

核心材料为《宣统政纪》宣统二年相关条；宫中朱批奏折、内阁大库题本与《清史稿》相关志传。它们可以支持年序、制度文本、机构设立或某些地点的行动，却不能自动代表全体居民的收入、教育、身份、性别关系或政治意愿。账本提示“以宣统二年（1910年）的同期文书为定位起点；官修记录、奏报或外交文本服务特定机构立场，具体日期、规模、地方执行与对手视角须以独立材料复核。”。因此，官方统计、外交措辞、报刊论战、教会叙述和军报数字必须分别置于生成它们的立场与制度目的之中，不能互相替代。

在地方尺度上，这一事件应被看作资源和信息重新配置的过程：经费能否落实，取决于税源和中介；学校、警政或新军能否运行，取决于人员、土地和日常协作；家庭与社群对改革的经历，则可能因性别、职业、族属、居处与迁徙而不同。现存材料只照亮少数路径，叙事须保留被登记者与未被登记者之间的差距。

账本条目\texttt{EQ-1910-ZIZHENG-YUAN-DOCUMENTS}记录，宣统二年（1910年），清发生“围绕「资政院开院，中央层面的立宪咨询机构开始运作」的上谕、奏折与部议在中央—地方行政链中流转”。账本将其背景说明为核心事件「资政院开院，中央层面的立宪咨询机构开始运作」要求中央或地方通过文书处理资源、信息或责任。，可见结果为该记录为后续按具体谕旨、奏折、条约、报刊或地方档案拆分事件提供可回查入口。。这一节点虽可放入改革、外交或革命的宏观年表，却不能脱离省份财政、地方官、商会、学校、军队、家庭与交通网络来理解；同一上谕、条约或战事在不同地域的实际后果并不一致。

核心材料为《宣统政纪》宣统二年相关条；宫中朱批奏折、内阁大库题本与《清史稿》相关志传。它们可以支持年序、制度文本、机构设立或某些地点的行动，却不能自动代表全体居民的收入、教育、身份、性别关系或政治意愿。账本提示“此行仅确认相关文书链和可追查的行政动作；不得由其直接推断政策有效、地方服从、民众态度或单一因果。”。因此，官方统计、外交措辞、报刊论战、教会叙述和军报数字必须分别置于生成它们的立场与制度目的之中，不能互相替代。

在地方尺度上，这一事件应被看作资源和信息重新配置的过程：经费能否落实，取决于税源和中介；学校、警政或新军能否运行，取决于人员、土地和日常协作；家庭与社群对改革的经历，则可能因性别、职业、族属、居处与迁徙而不同。现存材料只照亮少数路径，叙事须保留被登记者与未被登记者之间的差距。

账本条目\texttt{EQ-1910-FAMINE-AND-UNREST}记录，宣统二年（1910年），清与多地社会发生“灾荒、物价和财政压力持续影响晚清地方社会与政治动员”。账本将其背景说明为自然环境、市场波动和行政能力交互作用，可见结果为地方灾情与政治行动不可简单建立单因果关系。这一节点虽可放入改革、外交或革命的宏观年表，却不能脱离省份财政、地方官、商会、学校、军队、家庭与交通网络来理解；同一上谕、条约或战事在不同地域的实际后果并不一致。

核心材料为《宣统政纪》宣统二年相关条；地方志、督抚奏折及地方档案。它们可以支持年序、制度文本、机构设立或某些地点的行动，却不能自动代表全体居民的收入、教育、身份、性别关系或政治意愿。账本提示“以宣统二年（1910年）的同期文书为定位起点；官修记录、奏报或外交文本服务特定机构立场，具体日期、规模、地方执行与对手视角须以独立材料复核。”。因此，官方统计、外交措辞、报刊论战、教会叙述和军报数字必须分别置于生成它们的立场与制度目的之中，不能互相替代。

在地方尺度上，这一事件应被看作资源和信息重新配置的过程：经费能否落实，取决于税源和中介；学校、警政或新军能否运行，取决于人员、土地和日常协作；家庭与社群对改革的经历，则可能因性别、职业、族属、居处与迁徙而不同。现存材料只照亮少数路径，叙事须保留被登记者与未被登记者之间的差距。

\subsection{立宪、军政与革命前夜}

账本条目\texttt{EQ-1910-FAMINE-AND-UNREST-DOCUMENTS}记录，宣统二年（1910年），清与多地社会发生“围绕「灾荒、物价和财政压力持续影响晚清地方社会与政治动员」的地方呈报、案件或赈务记录将社会反应带入官府文书链”。账本将其背景说明为核心事件「灾荒、物价和财政压力持续影响晚清地方社会与政治动员」要求中央或地方通过文书处理资源、信息或责任。，可见结果为该记录为后续按具体谕旨、奏折、条约、报刊或地方档案拆分事件提供可回查入口。。这一节点虽可放入改革、外交或革命的宏观年表，却不能脱离省份财政、地方官、商会、学校、军队、家庭与交通网络来理解；同一上谕、条约或战事在不同地域的实际后果并不一致。

核心材料为《宣统政纪》宣统二年相关条；地方志、督抚奏折及地方档案。它们可以支持年序、制度文本、机构设立或某些地点的行动，却不能自动代表全体居民的收入、教育、身份、性别关系或政治意愿。账本提示“此行仅确认相关文书链和可追查的行政动作；不得由其直接推断政策有效、地方服从、民众态度或单一因果。”。因此，官方统计、外交措辞、报刊论战、教会叙述和军报数字必须分别置于生成它们的立场与制度目的之中，不能互相替代。

在地方尺度上，这一事件应被看作资源和信息重新配置的过程：经费能否落实，取决于税源和中介；学校、警政或新军能否运行，取决于人员、土地和日常协作；家庭与社群对改革的经历，则可能因性别、职业、族属、居处与迁徙而不同。现存材料只照亮少数路径，叙事须保留被登记者与未被登记者之间的差距。

账本条目\texttt{EQ-1911-RAILWAY-NATIONALIZATION}记录，宣统三年（1911年），清与四川、湖北等地发生“清廷宣布铁路干线国有化并以外债筹资，保路运动迅速扩大”。账本将其背景说明为中央财政、外债与地方商办铁路利益冲突，可见结果为政策公告、地方抗议和后续镇压应分开记录。这一节点虽可放入改革、外交或革命的宏观年表，却不能脱离省份财政、地方官、商会、学校、军队、家庭与交通网络来理解；同一上谕、条约或战事在不同地域的实际后果并不一致。

核心材料为《宣统政纪》宣统三年相关条；宫中朱批奏折、内阁大库题本与《清史稿》相关志传。它们可以支持年序、制度文本、机构设立或某些地点的行动，却不能自动代表全体居民的收入、教育、身份、性别关系或政治意愿。账本提示“以宣统三年（1911年）的同期文书为定位起点；官修记录、奏报或外交文本服务特定机构立场，具体日期、规模、地方执行与对手视角须以独立材料复核。”。因此，官方统计、外交措辞、报刊论战、教会叙述和军报数字必须分别置于生成它们的立场与制度目的之中，不能互相替代。

在地方尺度上，这一事件应被看作资源和信息重新配置的过程：经费能否落实，取决于税源和中介；学校、警政或新军能否运行，取决于人员、土地和日常协作；家庭与社群对改革的经历，则可能因性别、职业、族属、居处与迁徙而不同。现存材料只照亮少数路径，叙事须保留被登记者与未被登记者之间的差距。

账本条目\texttt{EQ-1911-RAILWAY-NATIONALIZATION-DOCUMENTS}记录，宣统三年（1911年），清与四川、湖北等地发生“围绕「清廷宣布铁路干线国有化并以外债筹资，保路运动迅速扩大」的上谕、奏折与部议在中央—地方行政链中流转”。账本将其背景说明为核心事件「清廷宣布铁路干线国有化并以外债筹资，保路运动迅速扩大」要求中央或地方通过文书处理资源、信息或责任。，可见结果为该记录为后续按具体谕旨、奏折、条约、报刊或地方档案拆分事件提供可回查入口。。这一节点虽可放入改革、外交或革命的宏观年表，却不能脱离省份财政、地方官、商会、学校、军队、家庭与交通网络来理解；同一上谕、条约或战事在不同地域的实际后果并不一致。

核心材料为《宣统政纪》宣统三年相关条；宫中朱批奏折、内阁大库题本与《清史稿》相关志传。它们可以支持年序、制度文本、机构设立或某些地点的行动，却不能自动代表全体居民的收入、教育、身份、性别关系或政治意愿。账本提示“此行仅确认相关文书链和可追查的行政动作；不得由其直接推断政策有效、地方服从、民众态度或单一因果。”。因此，官方统计、外交措辞、报刊论战、教会叙述和军报数字必须分别置于生成它们的立场与制度目的之中，不能互相替代。

在地方尺度上，这一事件应被看作资源和信息重新配置的过程：经费能否落实，取决于税源和中介；学校、警政或新军能否运行，取决于人员、土地和日常协作；家庭与社群对改革的经历，则可能因性别、职业、族属、居处与迁徙而不同。现存材料只照亮少数路径，叙事须保留被登记者与未被登记者之间的差距。

账本条目\texttt{EQ-1911-SICHUAN-RAILWAY-PROTEST}记录，宣统三年（1911年），清与四川保路运动发生“四川保路运动升级，赵尔丰处置引发更大政治危机”。账本将其背景说明为地方股东、商绅、学生和官府对路权与财政有不同诉求，可见结果为参与规模和暴力过程须依当地档案与报刊核验。这一节点虽可放入改革、外交或革命的宏观年表，却不能脱离省份财政、地方官、商会、学校、军队、家庭与交通网络来理解；同一上谕、条约或战事在不同地域的实际后果并不一致。

核心材料为《宣统政纪》宣统三年相关条；地方志、督抚奏折及地方档案。它们可以支持年序、制度文本、机构设立或某些地点的行动，却不能自动代表全体居民的收入、教育、身份、性别关系或政治意愿。账本提示“以宣统三年（1911年）的同期文书为定位起点；官修记录、奏报或外交文本服务特定机构立场，具体日期、规模、地方执行与对手视角须以独立材料复核。”。因此，官方统计、外交措辞、报刊论战、教会叙述和军报数字必须分别置于生成它们的立场与制度目的之中，不能互相替代。

在地方尺度上，这一事件应被看作资源和信息重新配置的过程：经费能否落实，取决于税源和中介；学校、警政或新军能否运行，取决于人员、土地和日常协作；家庭与社群对改革的经历，则可能因性别、职业、族属、居处与迁徙而不同。现存材料只照亮少数路径，叙事须保留被登记者与未被登记者之间的差距。

账本条目\texttt{EQ-1911-SICHUAN-RAILWAY-PROTEST-DOCUMENTS}记录，宣统三年（1911年），清与四川保路运动发生“围绕「四川保路运动升级，赵尔丰处置引发更大政治危机」的地方呈报、案件或赈务记录将社会反应带入官府文书链”。账本将其背景说明为核心事件「四川保路运动升级，赵尔丰处置引发更大政治危机」要求中央或地方通过文书处理资源、信息或责任。，可见结果为该记录为后续按具体谕旨、奏折、条约、报刊或地方档案拆分事件提供可回查入口。。这一节点虽可放入改革、外交或革命的宏观年表，却不能脱离省份财政、地方官、商会、学校、军队、家庭与交通网络来理解；同一上谕、条约或战事在不同地域的实际后果并不一致。

核心材料为《宣统政纪》宣统三年相关条；地方志、督抚奏折及地方档案。它们可以支持年序、制度文本、机构设立或某些地点的行动，却不能自动代表全体居民的收入、教育、身份、性别关系或政治意愿。账本提示“此行仅确认相关文书链和可追查的行政动作；不得由其直接推断政策有效、地方服从、民众态度或单一因果。”。因此，官方统计、外交措辞、报刊论战、教会叙述和军报数字必须分别置于生成它们的立场与制度目的之中，不能互相替代。

在地方尺度上，这一事件应被看作资源和信息重新配置的过程：经费能否落实，取决于税源和中介；学校、警政或新军能否运行，取决于人员、土地和日常协作；家庭与社群对改革的经历，则可能因性别、职业、族属、居处与迁徙而不同。现存材料只照亮少数路径，叙事须保留被登记者与未被登记者之间的差距。

账本条目\texttt{EQ-1911-WUCHANG-UPRISING}记录，宣统三年（1911年），革命军与清发生“武昌起义爆发，新军与革命团体控制武汉部分地区”。账本将其背景说明为新军组织、革命网络和铁路危机相互叠加，可见结果为起义的即时参与者、城市控制与后续省份响应需分层叙述。这一节点虽可放入改革、外交或革命的宏观年表，却不能脱离省份财政、地方官、商会、学校、军队、家庭与交通网络来理解；同一上谕、条约或战事在不同地域的实际后果并不一致。

核心材料为《宣统政纪》宣统三年相关条；军机处档、战事方略及地方奏报。它们可以支持年序、制度文本、机构设立或某些地点的行动，却不能自动代表全体居民的收入、教育、身份、性别关系或政治意愿。账本提示“以宣统三年（1911年）的同期文书为定位起点；官修记录、奏报或外交文本服务特定机构立场，具体日期、规模、地方执行与对手视角须以独立材料复核。”。因此，官方统计、外交措辞、报刊论战、教会叙述和军报数字必须分别置于生成它们的立场与制度目的之中，不能互相替代。

在地方尺度上，这一事件应被看作资源和信息重新配置的过程：经费能否落实，取决于税源和中介；学校、警政或新军能否运行，取决于人员、土地和日常协作；家庭与社群对改革的经历，则可能因性别、职业、族属、居处与迁徙而不同。现存材料只照亮少数路径，叙事须保留被登记者与未被登记者之间的差距。

账本条目\texttt{EQ-1911-WUCHANG-UPRISING-DOCUMENTS}记录，宣统三年（1911年），革命军与清发生“清廷围绕「武昌起义爆发，新军与革命团体控制武汉部分地区」调遣军队、筹措饷械并处理战报，中央与地方的军政文书形成连续记录”。账本将其背景说明为核心事件「武昌起义爆发，新军与革命团体控制武汉部分地区」要求中央或地方通过文书处理资源、信息或责任。，可见结果为该记录为后续按具体谕旨、奏折、条约、报刊或地方档案拆分事件提供可回查入口。。这一节点虽可放入改革、外交或革命的宏观年表，却不能脱离省份财政、地方官、商会、学校、军队、家庭与交通网络来理解；同一上谕、条约或战事在不同地域的实际后果并不一致。

核心材料为《宣统政纪》宣统三年相关条；军机处档、战事方略及地方奏报。它们可以支持年序、制度文本、机构设立或某些地点的行动，却不能自动代表全体居民的收入、教育、身份、性别关系或政治意愿。账本提示“此行仅确认相关文书链和可追查的行政动作；不得由其直接推断政策有效、地方服从、民众态度或单一因果。”。因此，官方统计、外交措辞、报刊论战、教会叙述和军报数字必须分别置于生成它们的立场与制度目的之中，不能互相替代。

在地方尺度上，这一事件应被看作资源和信息重新配置的过程：经费能否落实，取决于税源和中介；学校、警政或新军能否运行，取决于人员、土地和日常协作；家庭与社群对改革的经历，则可能因性别、职业、族属、居处与迁徙而不同。现存材料只照亮少数路径，叙事须保留被登记者与未被登记者之间的差距。

账本条目\texttt{EQ-1911-PROVINCIAL-INDEPENDENCE}记录，宣统三年（1911年），清与各省革命力量发生“多省宣布脱离清廷，帝国行政和财政网络快速分裂”。账本将其背景说明为武昌起义、地方军政力量和立宪派选择共同推动变化，可见结果为“独立”声明不等于各省形成稳定统一控制。这一节点虽可放入改革、外交或革命的宏观年表，却不能脱离省份财政、地方官、商会、学校、军队、家庭与交通网络来理解；同一上谕、条约或战事在不同地域的实际后果并不一致。

核心材料为《宣统政纪》宣统三年相关条；宫中朱批奏折、内阁大库题本与《清史稿》相关志传。它们可以支持年序、制度文本、机构设立或某些地点的行动，却不能自动代表全体居民的收入、教育、身份、性别关系或政治意愿。账本提示“以宣统三年（1911年）的同期文书为定位起点；官修记录、奏报或外交文本服务特定机构立场，具体日期、规模、地方执行与对手视角须以独立材料复核。”。因此，官方统计、外交措辞、报刊论战、教会叙述和军报数字必须分别置于生成它们的立场与制度目的之中，不能互相替代。

在地方尺度上，这一事件应被看作资源和信息重新配置的过程：经费能否落实，取决于税源和中介；学校、警政或新军能否运行，取决于人员、土地和日常协作；家庭与社群对改革的经历，则可能因性别、职业、族属、居处与迁徙而不同。现存材料只照亮少数路径，叙事须保留被登记者与未被登记者之间的差距。

账本条目\texttt{EQ-1911-PROVINCIAL-INDEPENDENCE-DOCUMENTS}记录，宣统三年（1911年），清与各省革命力量发生“围绕「多省宣布脱离清廷，帝国行政和财政网络快速分裂」的上谕、奏折与部议在中央—地方行政链中流转”。账本将其背景说明为核心事件「多省宣布脱离清廷，帝国行政和财政网络快速分裂」要求中央或地方通过文书处理资源、信息或责任。，可见结果为该记录为后续按具体谕旨、奏折、条约、报刊或地方档案拆分事件提供可回查入口。。这一节点虽可放入改革、外交或革命的宏观年表，却不能脱离省份财政、地方官、商会、学校、军队、家庭与交通网络来理解；同一上谕、条约或战事在不同地域的实际后果并不一致。

核心材料为《宣统政纪》宣统三年相关条；宫中朱批奏折、内阁大库题本与《清史稿》相关志传。它们可以支持年序、制度文本、机构设立或某些地点的行动，却不能自动代表全体居民的收入、教育、身份、性别关系或政治意愿。账本提示“此行仅确认相关文书链和可追查的行政动作；不得由其直接推断政策有效、地方服从、民众态度或单一因果。”。因此，官方统计、外交措辞、报刊论战、教会叙述和军报数字必须分别置于生成它们的立场与制度目的之中，不能互相替代。

在地方尺度上，这一事件应被看作资源和信息重新配置的过程：经费能否落实，取决于税源和中介；学校、警政或新军能否运行，取决于人员、土地和日常协作；家庭与社群对改革的经历，则可能因性别、职业、族属、居处与迁徙而不同。现存材料只照亮少数路径，叙事须保留被登记者与未被登记者之间的差距。

账本条目\texttt{EQ-1911-YUAN-SHIKAI-RETURN}记录，宣统三年（1911年），清、袁世凯与革命派发生“清廷起用袁世凯处理北方军政和议和事务”。账本将其背景说明为中央缺乏可依赖的军队与谈判人选，可见结果为袁在清廷、北洋与革命派之间的策略塑造退位进程。这一节点虽可放入改革、外交或革命的宏观年表，却不能脱离省份财政、地方官、商会、学校、军队、家庭与交通网络来理解；同一上谕、条约或战事在不同地域的实际后果并不一致。

核心材料为《宣统政纪》宣统三年相关条；宫中朱批奏折、内阁大库题本与《清史稿》相关志传。它们可以支持年序、制度文本、机构设立或某些地点的行动，却不能自动代表全体居民的收入、教育、身份、性别关系或政治意愿。账本提示“以宣统三年（1911年）的同期文书为定位起点；官修记录、奏报或外交文本服务特定机构立场，具体日期、规模、地方执行与对手视角须以独立材料复核。”。因此，官方统计、外交措辞、报刊论战、教会叙述和军报数字必须分别置于生成它们的立场与制度目的之中，不能互相替代。

在地方尺度上，这一事件应被看作资源和信息重新配置的过程：经费能否落实，取决于税源和中介；学校、警政或新军能否运行，取决于人员、土地和日常协作；家庭与社群对改革的经历，则可能因性别、职业、族属、居处与迁徙而不同。现存材料只照亮少数路径，叙事须保留被登记者与未被登记者之间的差距。

账本条目\texttt{EQ-1911-YUAN-SHIKAI-RETURN-DOCUMENTS}记录，宣统三年（1911年），清、袁世凯与革命派发生“围绕「清廷起用袁世凯处理北方军政和议和事务」的上谕、奏折与部议在中央—地方行政链中流转”。账本将其背景说明为核心事件「清廷起用袁世凯处理北方军政和议和事务」要求中央或地方通过文书处理资源、信息或责任。，可见结果为该记录为后续按具体谕旨、奏折、条约、报刊或地方档案拆分事件提供可回查入口。。这一节点虽可放入改革、外交或革命的宏观年表，却不能脱离省份财政、地方官、商会、学校、军队、家庭与交通网络来理解；同一上谕、条约或战事在不同地域的实际后果并不一致。

核心材料为《宣统政纪》宣统三年相关条；宫中朱批奏折、内阁大库题本与《清史稿》相关志传。它们可以支持年序、制度文本、机构设立或某些地点的行动，却不能自动代表全体居民的收入、教育、身份、性别关系或政治意愿。账本提示“此行仅确认相关文书链和可追查的行政动作；不得由其直接推断政策有效、地方服从、民众态度或单一因果。”。因此，官方统计、外交措辞、报刊论战、教会叙述和军报数字必须分别置于生成它们的立场与制度目的之中，不能互相替代。

在地方尺度上，这一事件应被看作资源和信息重新配置的过程：经费能否落实，取决于税源和中介；学校、警政或新军能否运行，取决于人员、土地和日常协作；家庭与社群对改革的经历，则可能因性别、职业、族属、居处与迁徙而不同。现存材料只照亮少数路径，叙事须保留被登记者与未被登记者之间的差距。

账本条目\texttt{EQ-1912-REPUBLIC-FOUNDING}记录，宣统四年（1912年），中华民国临时政府与清发生“中华民国临时政府在南京成立，与清廷退位谈判并行”。账本将其背景说明为各省独立与南北议和创造新的政权竞争格局，可见结果为临时政府的法理主张与实际控制范围需要区分。这一节点虽可放入改革、外交或革命的宏观年表，却不能脱离省份财政、地方官、商会、学校、军队、家庭与交通网络来理解；同一上谕、条约或战事在不同地域的实际后果并不一致。

核心材料为《宣统政纪》宣统四年相关条；宫中朱批奏折、内阁大库题本与《清史稿》相关志传。它们可以支持年序、制度文本、机构设立或某些地点的行动，却不能自动代表全体居民的收入、教育、身份、性别关系或政治意愿。账本提示“以宣统四年（1912年）的同期文书为定位起点；官修记录、奏报或外交文本服务特定机构立场，具体日期、规模、地方执行与对手视角须以独立材料复核。”。因此，官方统计、外交措辞、报刊论战、教会叙述和军报数字必须分别置于生成它们的立场与制度目的之中，不能互相替代。

在地方尺度上，这一事件应被看作资源和信息重新配置的过程：经费能否落实，取决于税源和中介；学校、警政或新军能否运行，取决于人员、土地和日常协作；家庭与社群对改革的经历，则可能因性别、职业、族属、居处与迁徙而不同。现存材料只照亮少数路径，叙事须保留被登记者与未被登记者之间的差距。

账本条目\texttt{EQ-1912-REPUBLIC-FOUNDING-DOCUMENTS}记录，宣统四年（1912年），中华民国临时政府与清发生“围绕「中华民国临时政府在南京成立，与清廷退位谈判并行」的上谕、奏折与部议在中央—地方行政链中流转”。账本将其背景说明为核心事件「中华民国临时政府在南京成立，与清廷退位谈判并行」要求中央或地方通过文书处理资源、信息或责任。，可见结果为该记录为后续按具体谕旨、奏折、条约、报刊或地方档案拆分事件提供可回查入口。。这一节点虽可放入改革、外交或革命的宏观年表，却不能脱离省份财政、地方官、商会、学校、军队、家庭与交通网络来理解；同一上谕、条约或战事在不同地域的实际后果并不一致。

核心材料为《宣统政纪》宣统四年相关条；宫中朱批奏折、内阁大库题本与《清史稿》相关志传。它们可以支持年序、制度文本、机构设立或某些地点的行动，却不能自动代表全体居民的收入、教育、身份、性别关系或政治意愿。账本提示“此行仅确认相关文书链和可追查的行政动作；不得由其直接推断政策有效、地方服从、民众态度或单一因果。”。因此，官方统计、外交措辞、报刊论战、教会叙述和军报数字必须分别置于生成它们的立场与制度目的之中，不能互相替代。

在地方尺度上，这一事件应被看作资源和信息重新配置的过程：经费能否落实，取决于税源和中介；学校、警政或新军能否运行，取决于人员、土地和日常协作；家庭与社群对改革的经历，则可能因性别、职业、族属、居处与迁徙而不同。现存材料只照亮少数路径，叙事须保留被登记者与未被登记者之间的差距。

账本条目\texttt{EQ-1912-QING-ABDICATION}记录，宣统四年（1912年），清、袁世凯与革命派发生“宣统帝颁布退位诏书，清帝国终结”。账本将其背景说明为南北谈判、财政军政压力与皇室待遇安排共同促成退位，可见结果为退位文本规定的优待与地方政治现实须分别追踪。这一节点虽可放入改革、外交或革命的宏观年表，却不能脱离省份财政、地方官、商会、学校、军队、家庭与交通网络来理解；同一上谕、条约或战事在不同地域的实际后果并不一致。

核心材料为《宣统政纪》宣统四年相关条；宫中朱批奏折、内阁大库题本与《清史稿》相关志传。它们可以支持年序、制度文本、机构设立或某些地点的行动，却不能自动代表全体居民的收入、教育、身份、性别关系或政治意愿。账本提示“以宣统四年（1912年）的同期文书为定位起点；官修记录、奏报或外交文本服务特定机构立场，具体日期、规模、地方执行与对手视角须以独立材料复核。”。因此，官方统计、外交措辞、报刊论战、教会叙述和军报数字必须分别置于生成它们的立场与制度目的之中，不能互相替代。

在地方尺度上，这一事件应被看作资源和信息重新配置的过程：经费能否落实，取决于税源和中介；学校、警政或新军能否运行，取决于人员、土地和日常协作；家庭与社群对改革的经历，则可能因性别、职业、族属、居处与迁徙而不同。现存材料只照亮少数路径，叙事须保留被登记者与未被登记者之间的差距。

账本条目\texttt{EQ-1912-QING-ABDICATION-DOCUMENTS}记录，宣统四年（1912年），清、袁世凯与革命派发生“围绕「宣统帝颁布退位诏书，清帝国终结」的上谕、奏折与部议在中央—地方行政链中流转”。账本将其背景说明为核心事件「宣统帝颁布退位诏书，清帝国终结」要求中央或地方通过文书处理资源、信息或责任。，可见结果为该记录为后续按具体谕旨、奏折、条约、报刊或地方档案拆分事件提供可回查入口。。这一节点虽可放入改革、外交或革命的宏观年表，却不能脱离省份财政、地方官、商会、学校、军队、家庭与交通网络来理解；同一上谕、条约或战事在不同地域的实际后果并不一致。

核心材料为《宣统政纪》宣统四年相关条；宫中朱批奏折、内阁大库题本与《清史稿》相关志传。它们可以支持年序、制度文本、机构设立或某些地点的行动，却不能自动代表全体居民的收入、教育、身份、性别关系或政治意愿。账本提示“此行仅确认相关文书链和可追查的行政动作；不得由其直接推断政策有效、地方服从、民众态度或单一因果。”。因此，官方统计、外交措辞、报刊论战、教会叙述和军报数字必须分别置于生成它们的立场与制度目的之中，不能互相替代。

在地方尺度上，这一事件应被看作资源和信息重新配置的过程：经费能否落实，取决于税源和中介；学校、警政或新军能否运行，取决于人员、土地和日常协作；家庭与社群对改革的经历，则可能因性别、职业、族属、居处与迁徙而不同。现存材料只照亮少数路径，叙事须保留被登记者与未被登记者之间的差距。

账本条目\texttt{EQ-1870-TIANJIN-MASSACRE-PRELUDE}记录，同治九年（1870年），清、法国与天津发生“「天津教案发生，法国领事、修女和中国民众死伤，引发外交危机」之前的条件、信息和争议已通过中央与地方材料显现，构成该事件可追踪的前史节点”。账本将其背景说明为教会活动、地方谣言和条约保护关系高度紧张，可见结果为为核心事件「天津教案发生，法国领事、修女和中国民众死伤，引发外交危机」提供时间、资源和行动者条件的检索入口，不能替代对核心事件本身的取证。。这一节点虽可放入改革、外交或革命的宏观年表，却不能脱离省份财政、地方官、商会、学校、军队、家庭与交通网络来理解；同一上谕、条约或战事在不同地域的实际后果并不一致。

核心材料为《清穆宗实录》同治九年相关条；《筹办夷务始末》相关年卷与中外照会、条约文本；同时期地方档案、地方志、日记或对方材料应按具体地区补检。它们可以支持年序、制度文本、机构设立或某些地点的行动，却不能自动代表全体居民的收入、教育、身份、性别关系或政治意愿。账本提示“官修编年和地方材料的记录密度与立场并不相同；本行不将条件、传闻、呈报或后见解释直接等同为确定因果。”。因此，官方统计、外交措辞、报刊论战、教会叙述和军报数字必须分别置于生成它们的立场与制度目的之中，不能互相替代。

在地方尺度上，这一事件应被看作资源和信息重新配置的过程：经费能否落实，取决于税源和中介；学校、警政或新军能否运行，取决于人员、土地和日常协作；家庭与社群对改革的经历，则可能因性别、职业、族属、居处与迁徙而不同。现存材料只照亮少数路径，叙事须保留被登记者与未被登记者之间的差距。

账本条目\texttt{EQ-1870-TIANJIN-MASSACRE-DECISION}记录，同治九年（1870年），清、法国与天津发生“围绕「天津教案发生，法国领事、修女和中国民众死伤，引发外交危机」的命令、调度、照会或呈报形成独立的中央—地方行动文书链”。账本将其背景说明为核心事件「天津教案发生，法国领事、修女和中国民众死伤，引发外交危机」要求相关机构处理人事、资源、军事、外交或司法事务。，可见结果为责任叙述在中法档案中差异显著，须避免单一归因。这一节点虽可放入改革、外交或革命的宏观年表，却不能脱离省份财政、地方官、商会、学校、军队、家庭与交通网络来理解；同一上谕、条约或战事在不同地域的实际后果并不一致。

核心材料为《清穆宗实录》同治九年相关条；《筹办夷务始末》相关年卷与中外照会、条约文本。它们可以支持年序、制度文本、机构设立或某些地点的行动，却不能自动代表全体居民的收入、教育、身份、性别关系或政治意愿。账本提示“奏折、上谕、部议、军机档或条约可证明行政动作和机构立场，不自动证明所有地区、群体或对手按其意图行动。”。因此，官方统计、外交措辞、报刊论战、教会叙述和军报数字必须分别置于生成它们的立场与制度目的之中，不能互相替代。

在地方尺度上，这一事件应被看作资源和信息重新配置的过程：经费能否落实，取决于税源和中介；学校、警政或新军能否运行，取决于人员、土地和日常协作；家庭与社群对改革的经历，则可能因性别、职业、族属、居处与迁徙而不同。现存材料只照亮少数路径，叙事须保留被登记者与未被登记者之间的差距。

账本条目\texttt{EQ-1870-TIANJIN-MASSACRE-AFTERMATH}记录，同治九年（1870年），清、法国与天津发生“「天津教案发生，法国领事、修女和中国民众死伤，引发外交危机」引发的善后、执行或地方回应构成可由不同材料追踪的后续节点”。账本将其背景说明为核心事件「天津教案发生，法国领事、修女和中国民众死伤，引发外交危机」改变了相关行动者面对的资源、秩序或风险。，可见结果为责任叙述在中法档案中差异显著，须避免单一归因。这一节点虽可放入改革、外交或革命的宏观年表，却不能脱离省份财政、地方官、商会、学校、军队、家庭与交通网络来理解；同一上谕、条约或战事在不同地域的实际后果并不一致。

核心材料为《清穆宗实录》同治九年相关条；《筹办夷务始末》相关年卷与中外照会、条约文本；相关地区的档案、志书、契约、报刊或对方材料。它们可以支持年序、制度文本、机构设立或某些地点的行动，却不能自动代表全体居民的收入、教育、身份、性别关系或政治意愿。账本提示“战后、改革后或条约后的记忆、地方记录和官方总结常具有事后选择性；后续影响须按地点、群体和时间分层，不作整体化推断。”。因此，官方统计、外交措辞、报刊论战、教会叙述和军报数字必须分别置于生成它们的立场与制度目的之中，不能互相替代。

在地方尺度上，这一事件应被看作资源和信息重新配置的过程：经费能否落实，取决于税源和中介；学校、警政或新军能否运行，取决于人员、土地和日常协作；家庭与社群对改革的经历，则可能因性别、职业、族属、居处与迁徙而不同。现存材料只照亮少数路径，叙事须保留被登记者与未被登记者之间的差距。

账本条目\texttt{EQ-1871-MUDAN-INCIDENT-PRELUDE}记录，同治十年（1871年），清、琉球、日本与台湾原住民发生“「琉球船民在台湾牡丹社一带遇害，事件成为日本对台政策借口」之前的条件、信息和争议已通过中央与地方材料显现，构成该事件可追踪的前史节点”。账本将其背景说明为海难、原住民地域与清朝地方治理边界交错，可见结果为为核心事件「琉球船民在台湾牡丹社一带遇害，事件成为日本对台政策借口」提供时间、资源和行动者条件的检索入口，不能替代对核心事件本身的取证。。这一节点虽可放入改革、外交或革命的宏观年表，却不能脱离省份财政、地方官、商会、学校、军队、家庭与交通网络来理解；同一上谕、条约或战事在不同地域的实际后果并不一致。

核心材料为《清穆宗实录》同治十年相关条；《筹办夷务始末》相关年卷与中外照会、条约文本；同时期地方档案、地方志、日记或对方材料应按具体地区补检。它们可以支持年序、制度文本、机构设立或某些地点的行动，却不能自动代表全体居民的收入、教育、身份、性别关系或政治意愿。账本提示“官修编年和地方材料的记录密度与立场并不相同；本行不将条件、传闻、呈报或后见解释直接等同为确定因果。”。因此，官方统计、外交措辞、报刊论战、教会叙述和军报数字必须分别置于生成它们的立场与制度目的之中，不能互相替代。

在地方尺度上，这一事件应被看作资源和信息重新配置的过程：经费能否落实，取决于税源和中介；学校、警政或新军能否运行，取决于人员、土地和日常协作；家庭与社群对改革的经历，则可能因性别、职业、族属、居处与迁徙而不同。现存材料只照亮少数路径，叙事须保留被登记者与未被登记者之间的差距。

账本条目\texttt{EQ-1871-MUDAN-INCIDENT-DECISION}记录，同治十年（1871年），清、琉球、日本与台湾原住民发生“围绕「琉球船民在台湾牡丹社一带遇害，事件成为日本对台政策借口」的命令、调度、照会或呈报形成独立的中央—地方行动文书链”。账本将其背景说明为核心事件「琉球船民在台湾牡丹社一带遇害，事件成为日本对台政策借口」要求相关机构处理人事、资源、军事、外交或司法事务。，可见结果为事件叙述须区分原住民行动、清方管辖与日方外交利用。这一节点虽可放入改革、外交或革命的宏观年表，却不能脱离省份财政、地方官、商会、学校、军队、家庭与交通网络来理解；同一上谕、条约或战事在不同地域的实际后果并不一致。

核心材料为《清穆宗实录》同治十年相关条；《筹办夷务始末》相关年卷与中外照会、条约文本。它们可以支持年序、制度文本、机构设立或某些地点的行动，却不能自动代表全体居民的收入、教育、身份、性别关系或政治意愿。账本提示“奏折、上谕、部议、军机档或条约可证明行政动作和机构立场，不自动证明所有地区、群体或对手按其意图行动。”。因此，官方统计、外交措辞、报刊论战、教会叙述和军报数字必须分别置于生成它们的立场与制度目的之中，不能互相替代。

在地方尺度上，这一事件应被看作资源和信息重新配置的过程：经费能否落实，取决于税源和中介；学校、警政或新军能否运行，取决于人员、土地和日常协作；家庭与社群对改革的经历，则可能因性别、职业、族属、居处与迁徙而不同。现存材料只照亮少数路径，叙事须保留被登记者与未被登记者之间的差距。

账本条目\texttt{EQ-1871-MUDAN-INCIDENT-AFTERMATH}记录，同治十年（1871年），清、琉球、日本与台湾原住民发生“「琉球船民在台湾牡丹社一带遇害，事件成为日本对台政策借口」引发的善后、执行或地方回应构成可由不同材料追踪的后续节点”。账本将其背景说明为核心事件「琉球船民在台湾牡丹社一带遇害，事件成为日本对台政策借口」改变了相关行动者面对的资源、秩序或风险。，可见结果为事件叙述须区分原住民行动、清方管辖与日方外交利用。这一节点虽可放入改革、外交或革命的宏观年表，却不能脱离省份财政、地方官、商会、学校、军队、家庭与交通网络来理解；同一上谕、条约或战事在不同地域的实际后果并不一致。

核心材料为《清穆宗实录》同治十年相关条；《筹办夷务始末》相关年卷与中外照会、条约文本；相关地区的档案、志书、契约、报刊或对方材料。它们可以支持年序、制度文本、机构设立或某些地点的行动，却不能自动代表全体居民的收入、教育、身份、性别关系或政治意愿。账本提示“战后、改革后或条约后的记忆、地方记录和官方总结常具有事后选择性；后续影响须按地点、群体和时间分层，不作整体化推断。”。因此，官方统计、外交措辞、报刊论战、教会叙述和军报数字必须分别置于生成它们的立场与制度目的之中，不能互相替代。

在地方尺度上，这一事件应被看作资源和信息重新配置的过程：经费能否落实，取决于税源和中介；学校、警政或新军能否运行，取决于人员、土地和日常协作；家庭与社群对改革的经历，则可能因性别、职业、族属、居处与迁徙而不同。现存材料只照亮少数路径，叙事须保留被登记者与未被登记者之间的差距。

账本条目\texttt{EQ-1871-QING-STUDENTS-ABROAD-PRELUDE}记录，同治十年（1871年），清与美国发生“「清廷派遣幼童赴美留学的计划启动，洋务教育进入海外培训阶段」之前的条件、信息和争议已通过中央与地方材料显现，构成该事件可追踪的前史节点”。账本将其背景说明为技术、语言与外交人才需求上升，可见结果为为核心事件「清廷派遣幼童赴美留学的计划启动，洋务教育进入海外培训阶段」提供时间、资源和行动者条件的检索入口，不能替代对核心事件本身的取证。。这一节点虽可放入改革、外交或革命的宏观年表，却不能脱离省份财政、地方官、商会、学校、军队、家庭与交通网络来理解；同一上谕、条约或战事在不同地域的实际后果并不一致。

核心材料为《清穆宗实录》同治十年相关条；内阁档、文集或报刊相关记录；同时期地方档案、地方志、日记或对方材料应按具体地区补检。它们可以支持年序、制度文本、机构设立或某些地点的行动，却不能自动代表全体居民的收入、教育、身份、性别关系或政治意愿。账本提示“官修编年和地方材料的记录密度与立场并不相同；本行不将条件、传闻、呈报或后见解释直接等同为确定因果。”。因此，官方统计、外交措辞、报刊论战、教会叙述和军报数字必须分别置于生成它们的立场与制度目的之中，不能互相替代。

在地方尺度上，这一事件应被看作资源和信息重新配置的过程：经费能否落实，取决于税源和中介；学校、警政或新军能否运行，取决于人员、土地和日常协作；家庭与社群对改革的经历，则可能因性别、职业、族属、居处与迁徙而不同。现存材料只照亮少数路径，叙事须保留被登记者与未被登记者之间的差距。

账本条目\texttt{EQ-1871-QING-STUDENTS-ABROAD-DECISION}记录，同治十年（1871年），清与美国发生“围绕「清廷派遣幼童赴美留学的计划启动，洋务教育进入海外培训阶段」的命令、调度、照会或呈报形成独立的中央—地方行动文书链”。账本将其背景说明为核心事件「清廷派遣幼童赴美留学的计划启动，洋务教育进入海外培训阶段」要求相关机构处理人事、资源、军事、外交或司法事务。，可见结果为选派和留学经历不代表全国教育制度已转型。这一节点虽可放入改革、外交或革命的宏观年表，却不能脱离省份财政、地方官、商会、学校、军队、家庭与交通网络来理解；同一上谕、条约或战事在不同地域的实际后果并不一致。

核心材料为《清穆宗实录》同治十年相关条；内阁档、文集或报刊相关记录。它们可以支持年序、制度文本、机构设立或某些地点的行动，却不能自动代表全体居民的收入、教育、身份、性别关系或政治意愿。账本提示“奏折、上谕、部议、军机档或条约可证明行政动作和机构立场，不自动证明所有地区、群体或对手按其意图行动。”。因此，官方统计、外交措辞、报刊论战、教会叙述和军报数字必须分别置于生成它们的立场与制度目的之中，不能互相替代。

在地方尺度上，这一事件应被看作资源和信息重新配置的过程：经费能否落实，取决于税源和中介；学校、警政或新军能否运行，取决于人员、土地和日常协作；家庭与社群对改革的经历，则可能因性别、职业、族属、居处与迁徙而不同。现存材料只照亮少数路径，叙事须保留被登记者与未被登记者之间的差距。

账本条目\texttt{EQ-1871-QING-STUDENTS-ABROAD-AFTERMATH}记录，同治十年（1871年），清与美国发生“「清廷派遣幼童赴美留学的计划启动，洋务教育进入海外培训阶段」引发的善后、执行或地方回应构成可由不同材料追踪的后续节点”。账本将其背景说明为核心事件「清廷派遣幼童赴美留学的计划启动，洋务教育进入海外培训阶段」改变了相关行动者面对的资源、秩序或风险。，可见结果为选派和留学经历不代表全国教育制度已转型。这一节点虽可放入改革、外交或革命的宏观年表，却不能脱离省份财政、地方官、商会、学校、军队、家庭与交通网络来理解；同一上谕、条约或战事在不同地域的实际后果并不一致。

核心材料为《清穆宗实录》同治十年相关条；内阁档、文集或报刊相关记录；相关地区的档案、志书、契约、报刊或对方材料。它们可以支持年序、制度文本、机构设立或某些地点的行动，却不能自动代表全体居民的收入、教育、身份、性别关系或政治意愿。账本提示“战后、改革后或条约后的记忆、地方记录和官方总结常具有事后选择性；后续影响须按地点、群体和时间分层，不作整体化推断。”。因此，官方统计、外交措辞、报刊论战、教会叙述和军报数字必须分别置于生成它们的立场与制度目的之中，不能互相替代。

在地方尺度上，这一事件应被看作资源和信息重新配置的过程：经费能否落实，取决于税源和中介；学校、警政或新军能否运行，取决于人员、土地和日常协作；家庭与社群对改革的经历，则可能因性别、职业、族属、居处与迁徙而不同。现存材料只照亮少数路径，叙事须保留被登记者与未被登记者之间的差距。

账本条目\texttt{EQ-1872-CHINA-MERCHANTS-STEAM-NAVIGATION-PRELUDE}记录，同治十一年（1872年），清与上海商人发生“「轮船招商局创办，试图以官督商办方式发展近代航运」之前的条件、信息和争议已通过中央与地方材料显现，构成该事件可追踪的前史节点”。账本将其背景说明为外轮竞争和沿海、长江运输需求推动新企业，可见结果为为核心事件「轮船招商局创办，试图以官督商办方式发展近代航运」提供时间、资源和行动者条件的检索入口，不能替代对核心事件本身的取证。。这一节点虽可放入改革、外交或革命的宏观年表，却不能脱离省份财政、地方官、商会、学校、军队、家庭与交通网络来理解；同一上谕、条约或战事在不同地域的实际后果并不一致。

核心材料为《清穆宗实录》同治十一年相关条；户部奏销、盐政、河工或海关档案；同时期地方档案、地方志、日记或对方材料应按具体地区补检。它们可以支持年序、制度文本、机构设立或某些地点的行动，却不能自动代表全体居民的收入、教育、身份、性别关系或政治意愿。账本提示“官修编年和地方材料的记录密度与立场并不相同；本行不将条件、传闻、呈报或后见解释直接等同为确定因果。”。因此，官方统计、外交措辞、报刊论战、教会叙述和军报数字必须分别置于生成它们的立场与制度目的之中，不能互相替代。

在地方尺度上，这一事件应被看作资源和信息重新配置的过程：经费能否落实，取决于税源和中介；学校、警政或新军能否运行，取决于人员、土地和日常协作；家庭与社群对改革的经历，则可能因性别、职业、族属、居处与迁徙而不同。现存材料只照亮少数路径，叙事须保留被登记者与未被登记者之间的差距。

账本条目\texttt{EQ-1872-CHINA-MERCHANTS-STEAM-NAVIGATION-DECISION}记录，同治十一年（1872年），清与上海商人发生“围绕「轮船招商局创办，试图以官督商办方式发展近代航运」的命令、调度、照会或呈报形成独立的中央—地方行动文书链”。账本将其背景说明为核心事件「轮船招商局创办，试图以官督商办方式发展近代航运」要求相关机构处理人事、资源、军事、外交或司法事务。，可见结果为企业性质、资本来源和经营绩效需依账册分析。这一节点虽可放入改革、外交或革命的宏观年表，却不能脱离省份财政、地方官、商会、学校、军队、家庭与交通网络来理解；同一上谕、条约或战事在不同地域的实际后果并不一致。

核心材料为《清穆宗实录》同治十一年相关条；户部奏销、盐政、河工或海关档案。它们可以支持年序、制度文本、机构设立或某些地点的行动，却不能自动代表全体居民的收入、教育、身份、性别关系或政治意愿。账本提示“奏折、上谕、部议、军机档或条约可证明行政动作和机构立场，不自动证明所有地区、群体或对手按其意图行动。”。因此，官方统计、外交措辞、报刊论战、教会叙述和军报数字必须分别置于生成它们的立场与制度目的之中，不能互相替代。

在地方尺度上，这一事件应被看作资源和信息重新配置的过程：经费能否落实，取决于税源和中介；学校、警政或新军能否运行，取决于人员、土地和日常协作；家庭与社群对改革的经历，则可能因性别、职业、族属、居处与迁徙而不同。现存材料只照亮少数路径，叙事须保留被登记者与未被登记者之间的差距。

账本条目\texttt{EQ-1872-CHINA-MERCHANTS-STEAM-NAVIGATION-AFTERMATH}记录，同治十一年（1872年），清与上海商人发生“「轮船招商局创办，试图以官督商办方式发展近代航运」引发的善后、执行或地方回应构成可由不同材料追踪的后续节点”。账本将其背景说明为核心事件「轮船招商局创办，试图以官督商办方式发展近代航运」改变了相关行动者面对的资源、秩序或风险。，可见结果为企业性质、资本来源和经营绩效需依账册分析。这一节点虽可放入改革、外交或革命的宏观年表，却不能脱离省份财政、地方官、商会、学校、军队、家庭与交通网络来理解；同一上谕、条约或战事在不同地域的实际后果并不一致。

核心材料为《清穆宗实录》同治十一年相关条；户部奏销、盐政、河工或海关档案；相关地区的档案、志书、契约、报刊或对方材料。它们可以支持年序、制度文本、机构设立或某些地点的行动，却不能自动代表全体居民的收入、教育、身份、性别关系或政治意愿。账本提示“战后、改革后或条约后的记忆、地方记录和官方总结常具有事后选择性；后续影响须按地点、群体和时间分层，不作整体化推断。”。因此，官方统计、外交措辞、报刊论战、教会叙述和军报数字必须分别置于生成它们的立场与制度目的之中，不能互相替代。

在地方尺度上，这一事件应被看作资源和信息重新配置的过程：经费能否落实，取决于税源和中介；学校、警政或新军能否运行，取决于人员、土地和日常协作；家庭与社群对改革的经历，则可能因性别、职业、族属、居处与迁徙而不同。现存材料只照亮少数路径，叙事须保留被登记者与未被登记者之间的差距。

账本条目\texttt{EQ-1873-JAPAN-RYUKYU-DISPUTE-PRELUDE}记录，同治十二年（1873年），清、日本与琉球发生“「日本围绕琉球与台湾事务加强对清交涉」之前的条件、信息和争议已通过中央与地方材料显现，构成该事件可追踪的前史节点”。账本将其背景说明为日本明治国家扩张与东亚朝贡关系重组，可见结果为为核心事件「日本围绕琉球与台湾事务加强对清交涉」提供时间、资源和行动者条件的检索入口，不能替代对核心事件本身的取证。。这一节点虽可放入改革、外交或革命的宏观年表，却不能脱离省份财政、地方官、商会、学校、军队、家庭与交通网络来理解；同一上谕、条约或战事在不同地域的实际后果并不一致。

核心材料为《清穆宗实录》同治十二年相关条；《筹办夷务始末》相关年卷与中外照会、条约文本；同时期地方档案、地方志、日记或对方材料应按具体地区补检。它们可以支持年序、制度文本、机构设立或某些地点的行动，却不能自动代表全体居民的收入、教育、身份、性别关系或政治意愿。账本提示“官修编年和地方材料的记录密度与立场并不相同；本行不将条件、传闻、呈报或后见解释直接等同为确定因果。”。因此，官方统计、外交措辞、报刊论战、教会叙述和军报数字必须分别置于生成它们的立场与制度目的之中，不能互相替代。

在地方尺度上，这一事件应被看作资源和信息重新配置的过程：经费能否落实，取决于税源和中介；学校、警政或新军能否运行，取决于人员、土地和日常协作；家庭与社群对改革的经历，则可能因性别、职业、族属、居处与迁徙而不同。现存材料只照亮少数路径，叙事须保留被登记者与未被登记者之间的差距。

账本条目\texttt{EQ-1873-JAPAN-RYUKYU-DISPUTE-DECISION}记录，同治十二年（1873年），清、日本与琉球发生“围绕「日本围绕琉球与台湾事务加强对清交涉」的命令、调度、照会或呈报形成独立的中央—地方行动文书链”。账本将其背景说明为核心事件「日本围绕琉球与台湾事务加强对清交涉」要求相关机构处理人事、资源、军事、外交或司法事务。，可见结果为册封、属国和主权词汇应置于各方文本语境中。这一节点虽可放入改革、外交或革命的宏观年表，却不能脱离省份财政、地方官、商会、学校、军队、家庭与交通网络来理解；同一上谕、条约或战事在不同地域的实际后果并不一致。

核心材料为《清穆宗实录》同治十二年相关条；《筹办夷务始末》相关年卷与中外照会、条约文本。它们可以支持年序、制度文本、机构设立或某些地点的行动，却不能自动代表全体居民的收入、教育、身份、性别关系或政治意愿。账本提示“奏折、上谕、部议、军机档或条约可证明行政动作和机构立场，不自动证明所有地区、群体或对手按其意图行动。”。因此，官方统计、外交措辞、报刊论战、教会叙述和军报数字必须分别置于生成它们的立场与制度目的之中，不能互相替代。

在地方尺度上，这一事件应被看作资源和信息重新配置的过程：经费能否落实，取决于税源和中介；学校、警政或新军能否运行，取决于人员、土地和日常协作；家庭与社群对改革的经历，则可能因性别、职业、族属、居处与迁徙而不同。现存材料只照亮少数路径，叙事须保留被登记者与未被登记者之间的差距。

账本条目\texttt{EQ-1873-JAPAN-RYUKYU-DISPUTE-AFTERMATH}记录，同治十二年（1873年），清、日本与琉球发生“「日本围绕琉球与台湾事务加强对清交涉」引发的善后、执行或地方回应构成可由不同材料追踪的后续节点”。账本将其背景说明为核心事件「日本围绕琉球与台湾事务加强对清交涉」改变了相关行动者面对的资源、秩序或风险。，可见结果为册封、属国和主权词汇应置于各方文本语境中。这一节点虽可放入改革、外交或革命的宏观年表，却不能脱离省份财政、地方官、商会、学校、军队、家庭与交通网络来理解；同一上谕、条约或战事在不同地域的实际后果并不一致。

核心材料为《清穆宗实录》同治十二年相关条；《筹办夷务始末》相关年卷与中外照会、条约文本；相关地区的档案、志书、契约、报刊或对方材料。它们可以支持年序、制度文本、机构设立或某些地点的行动，却不能自动代表全体居民的收入、教育、身份、性别关系或政治意愿。账本提示“战后、改革后或条约后的记忆、地方记录和官方总结常具有事后选择性；后续影响须按地点、群体和时间分层，不作整体化推断。”。因此，官方统计、外交措辞、报刊论战、教会叙述和军报数字必须分别置于生成它们的立场与制度目的之中，不能互相替代。

在地方尺度上，这一事件应被看作资源和信息重新配置的过程：经费能否落实，取决于税源和中介；学校、警政或新军能否运行，取决于人员、土地和日常协作；家庭与社群对改革的经历，则可能因性别、职业、族属、居处与迁徙而不同。现存材料只照亮少数路径，叙事须保留被登记者与未被登记者之间的差距。

账本条目\texttt{EQ-1874-JAPANESE-TAIWAN-EXPEDITION-PRELUDE}记录，同治十三年（1874年），清与日本发生“「日本以牡丹社事件为由出兵台湾南部，清廷被迫应对」之前的条件、信息和争议已通过中央与地方材料显现，构成该事件可追踪的前史节点”。账本将其背景说明为清廷对台湾东部和原住民地区的治理能力有限，可见结果为为核心事件「日本以牡丹社事件为由出兵台湾南部，清廷被迫应对」提供时间、资源和行动者条件的检索入口，不能替代对核心事件本身的取证。。这一节点虽可放入改革、外交或革命的宏观年表，却不能脱离省份财政、地方官、商会、学校、军队、家庭与交通网络来理解；同一上谕、条约或战事在不同地域的实际后果并不一致。

核心材料为《清穆宗实录》同治十三年相关条；军机处档、战事方略及地方奏报；同时期地方档案、地方志、日记或对方材料应按具体地区补检。它们可以支持年序、制度文本、机构设立或某些地点的行动，却不能自动代表全体居民的收入、教育、身份、性别关系或政治意愿。账本提示“官修编年和地方材料的记录密度与立场并不相同；本行不将条件、传闻、呈报或后见解释直接等同为确定因果。”。因此，官方统计、外交措辞、报刊论战、教会叙述和军报数字必须分别置于生成它们的立场与制度目的之中，不能互相替代。

在地方尺度上，这一事件应被看作资源和信息重新配置的过程：经费能否落实，取决于税源和中介；学校、警政或新军能否运行，取决于人员、土地和日常协作；家庭与社群对改革的经历，则可能因性别、职业、族属、居处与迁徙而不同。现存材料只照亮少数路径，叙事须保留被登记者与未被登记者之间的差距。

账本条目\texttt{EQ-1874-JAPANESE-TAIWAN-EXPEDITION-DECISION}记录，同治十三年（1874年），清与日本发生“围绕「日本以牡丹社事件为由出兵台湾南部，清廷被迫应对」的命令、调度、照会或呈报形成独立的中央—地方行动文书链”。账本将其背景说明为核心事件「日本以牡丹社事件为由出兵台湾南部，清廷被迫应对」要求相关机构处理人事、资源、军事、外交或司法事务。，可见结果为战后赔款与外交安排刺激清廷加强台湾经营。这一节点虽可放入改革、外交或革命的宏观年表，却不能脱离省份财政、地方官、商会、学校、军队、家庭与交通网络来理解；同一上谕、条约或战事在不同地域的实际后果并不一致。

核心材料为《清穆宗实录》同治十三年相关条；军机处档、战事方略及地方奏报。它们可以支持年序、制度文本、机构设立或某些地点的行动，却不能自动代表全体居民的收入、教育、身份、性别关系或政治意愿。账本提示“奏折、上谕、部议、军机档或条约可证明行政动作和机构立场，不自动证明所有地区、群体或对手按其意图行动。”。因此，官方统计、外交措辞、报刊论战、教会叙述和军报数字必须分别置于生成它们的立场与制度目的之中，不能互相替代。

在地方尺度上，这一事件应被看作资源和信息重新配置的过程：经费能否落实，取决于税源和中介；学校、警政或新军能否运行，取决于人员、土地和日常协作；家庭与社群对改革的经历，则可能因性别、职业、族属、居处与迁徙而不同。现存材料只照亮少数路径，叙事须保留被登记者与未被登记者之间的差距。

账本条目\texttt{EQ-1874-JAPANESE-TAIWAN-EXPEDITION-AFTERMATH}记录，同治十三年（1874年），清与日本发生“「日本以牡丹社事件为由出兵台湾南部，清廷被迫应对」引发的善后、执行或地方回应构成可由不同材料追踪的后续节点”。账本将其背景说明为核心事件「日本以牡丹社事件为由出兵台湾南部，清廷被迫应对」改变了相关行动者面对的资源、秩序或风险。，可见结果为战后赔款与外交安排刺激清廷加强台湾经营。这一节点虽可放入改革、外交或革命的宏观年表，却不能脱离省份财政、地方官、商会、学校、军队、家庭与交通网络来理解；同一上谕、条约或战事在不同地域的实际后果并不一致。

核心材料为《清穆宗实录》同治十三年相关条；军机处档、战事方略及地方奏报；相关地区的档案、志书、契约、报刊或对方材料。它们可以支持年序、制度文本、机构设立或某些地点的行动，却不能自动代表全体居民的收入、教育、身份、性别关系或政治意愿。账本提示“战后、改革后或条约后的记忆、地方记录和官方总结常具有事后选择性；后续影响须按地点、群体和时间分层，不作整体化推断。”。因此，官方统计、外交措辞、报刊论战、教会叙述和军报数字必须分别置于生成它们的立场与制度目的之中，不能互相替代。

在地方尺度上，这一事件应被看作资源和信息重新配置的过程：经费能否落实，取决于税源和中介；学校、警政或新军能否运行，取决于人员、土地和日常协作；家庭与社群对改革的经历，则可能因性别、职业、族属、居处与迁徙而不同。现存材料只照亮少数路径，叙事须保留被登记者与未被登记者之间的差距。

账本条目\texttt{EQ-1874-TONGZHI-DEATH-PRELUDE}记录，同治十三年（1874年），清发生“「同治帝去世，慈禧选择光绪帝继位，再次形成幼主政治」之前的条件、信息和争议已通过中央与地方材料显现，构成该事件可追踪的前史节点”。账本将其背景说明为同治无嗣使宗法与宫廷权力问题交织，可见结果为为核心事件「同治帝去世，慈禧选择光绪帝继位，再次形成幼主政治」提供时间、资源和行动者条件的检索入口，不能替代对核心事件本身的取证。。这一节点虽可放入改革、外交或革命的宏观年表，却不能脱离省份财政、地方官、商会、学校、军队、家庭与交通网络来理解；同一上谕、条约或战事在不同地域的实际后果并不一致。

核心材料为《清穆宗实录》同治十三年相关条；宫中朱批奏折、内阁大库题本与《清史稿》相关志传；同时期地方档案、地方志、日记或对方材料应按具体地区补检。它们可以支持年序、制度文本、机构设立或某些地点的行动，却不能自动代表全体居民的收入、教育、身份、性别关系或政治意愿。账本提示“官修编年和地方材料的记录密度与立场并不相同；本行不将条件、传闻、呈报或后见解释直接等同为确定因果。”。因此，官方统计、外交措辞、报刊论战、教会叙述和军报数字必须分别置于生成它们的立场与制度目的之中，不能互相替代。

在地方尺度上，这一事件应被看作资源和信息重新配置的过程：经费能否落实，取决于税源和中介；学校、警政或新军能否运行，取决于人员、土地和日常协作；家庭与社群对改革的经历，则可能因性别、职业、族属、居处与迁徙而不同。现存材料只照亮少数路径，叙事须保留被登记者与未被登记者之间的差距。

账本条目\texttt{EQ-1874-TONGZHI-DEATH-DECISION}记录，同治十三年（1874年），清发生“围绕「同治帝去世，慈禧选择光绪帝继位，再次形成幼主政治」的命令、调度、照会或呈报形成独立的中央—地方行动文书链”。账本将其背景说明为核心事件「同治帝去世，慈禧选择光绪帝继位，再次形成幼主政治」要求相关机构处理人事、资源、军事、外交或司法事务。，可见结果为慈禧重新垂帘，晚清中枢进入新阶段。这一节点虽可放入改革、外交或革命的宏观年表，却不能脱离省份财政、地方官、商会、学校、军队、家庭与交通网络来理解；同一上谕、条约或战事在不同地域的实际后果并不一致。

核心材料为《清穆宗实录》同治十三年相关条；宫中朱批奏折、内阁大库题本与《清史稿》相关志传。它们可以支持年序、制度文本、机构设立或某些地点的行动，却不能自动代表全体居民的收入、教育、身份、性别关系或政治意愿。账本提示“奏折、上谕、部议、军机档或条约可证明行政动作和机构立场，不自动证明所有地区、群体或对手按其意图行动。”。因此，官方统计、外交措辞、报刊论战、教会叙述和军报数字必须分别置于生成它们的立场与制度目的之中，不能互相替代。

在地方尺度上，这一事件应被看作资源和信息重新配置的过程：经费能否落实，取决于税源和中介；学校、警政或新军能否运行，取决于人员、土地和日常协作；家庭与社群对改革的经历，则可能因性别、职业、族属、居处与迁徙而不同。现存材料只照亮少数路径，叙事须保留被登记者与未被登记者之间的差距。

账本条目\texttt{EQ-1874-TONGZHI-DEATH-AFTERMATH}记录，同治十三年（1874年），清发生“「同治帝去世，慈禧选择光绪帝继位，再次形成幼主政治」引发的善后、执行或地方回应构成可由不同材料追踪的后续节点”。账本将其背景说明为核心事件「同治帝去世，慈禧选择光绪帝继位，再次形成幼主政治」改变了相关行动者面对的资源、秩序或风险。，可见结果为慈禧重新垂帘，晚清中枢进入新阶段。这一节点虽可放入改革、外交或革命的宏观年表，却不能脱离省份财政、地方官、商会、学校、军队、家庭与交通网络来理解；同一上谕、条约或战事在不同地域的实际后果并不一致。

核心材料为《清穆宗实录》同治十三年相关条；宫中朱批奏折、内阁大库题本与《清史稿》相关志传；相关地区的档案、志书、契约、报刊或对方材料。它们可以支持年序、制度文本、机构设立或某些地点的行动，却不能自动代表全体居民的收入、教育、身份、性别关系或政治意愿。账本提示“战后、改革后或条约后的记忆、地方记录和官方总结常具有事后选择性；后续影响须按地点、群体和时间分层，不作整体化推断。”。因此，官方统计、外交措辞、报刊论战、教会叙述和军报数字必须分别置于生成它们的立场与制度目的之中，不能互相替代。

在地方尺度上，这一事件应被看作资源和信息重新配置的过程：经费能否落实，取决于税源和中介；学校、警政或新军能否运行，取决于人员、土地和日常协作；家庭与社群对改革的经历，则可能因性别、职业、族属、居处与迁徙而不同。现存材料只照亮少数路径，叙事须保留被登记者与未被登记者之间的差距。

账本条目\texttt{EQ-1875-YUNNAN-MUSLIM-WAR-PRELUDE}记录，光绪元年（1875年），清与云南回民社群发生“「杜文秀政权覆亡前后，云南回民战争进入惨烈收束阶段」之前的条件、信息和争议已通过中央与地方材料显现，构成该事件可追踪的前史节点”。账本将其背景说明为地方军政、矿业贸易与族群冲突长期累积，可见结果为为核心事件「杜文秀政权覆亡前后，云南回民战争进入惨烈收束阶段」提供时间、资源和行动者条件的检索入口，不能替代对核心事件本身的取证。。这一节点虽可放入改革、外交或革命的宏观年表，却不能脱离省份财政、地方官、商会、学校、军队、家庭与交通网络来理解；同一上谕、条约或战事在不同地域的实际后果并不一致。

核心材料为《清德宗实录》光绪元年相关条；军机处档、战事方略及地方奏报；同时期地方档案、地方志、日记或对方材料应按具体地区补检。它们可以支持年序、制度文本、机构设立或某些地点的行动，却不能自动代表全体居民的收入、教育、身份、性别关系或政治意愿。账本提示“官修编年和地方材料的记录密度与立场并不相同；本行不将条件、传闻、呈报或后见解释直接等同为确定因果。”。因此，官方统计、外交措辞、报刊论战、教会叙述和军报数字必须分别置于生成它们的立场与制度目的之中，不能互相替代。

在地方尺度上，这一事件应被看作资源和信息重新配置的过程：经费能否落实，取决于税源和中介；学校、警政或新军能否运行，取决于人员、土地和日常协作；家庭与社群对改革的经历，则可能因性别、职业、族属、居处与迁徙而不同。现存材料只照亮少数路径，叙事须保留被登记者与未被登记者之间的差距。

账本条目\texttt{EQ-1875-YUNNAN-MUSLIM-WAR-DECISION}记录，光绪元年（1875年），清与云南回民社群发生“围绕「杜文秀政权覆亡前后，云南回民战争进入惨烈收束阶段」的命令、调度、照会或呈报形成独立的中央—地方行动文书链”。账本将其背景说明为核心事件「杜文秀政权覆亡前后，云南回民战争进入惨烈收束阶段」要求相关机构处理人事、资源、军事、外交或司法事务。，可见结果为死亡、迁徙与复垦规模的估计须保持不确定性。这一节点虽可放入改革、外交或革命的宏观年表，却不能脱离省份财政、地方官、商会、学校、军队、家庭与交通网络来理解；同一上谕、条约或战事在不同地域的实际后果并不一致。

核心材料为《清德宗实录》光绪元年相关条；军机处档、战事方略及地方奏报。它们可以支持年序、制度文本、机构设立或某些地点的行动，却不能自动代表全体居民的收入、教育、身份、性别关系或政治意愿。账本提示“奏折、上谕、部议、军机档或条约可证明行政动作和机构立场，不自动证明所有地区、群体或对手按其意图行动。”。因此，官方统计、外交措辞、报刊论战、教会叙述和军报数字必须分别置于生成它们的立场与制度目的之中，不能互相替代。

在地方尺度上，这一事件应被看作资源和信息重新配置的过程：经费能否落实，取决于税源和中介；学校、警政或新军能否运行，取决于人员、土地和日常协作；家庭与社群对改革的经历，则可能因性别、职业、族属、居处与迁徙而不同。现存材料只照亮少数路径，叙事须保留被登记者与未被登记者之间的差距。

账本条目\texttt{EQ-1875-YUNNAN-MUSLIM-WAR-AFTERMATH}记录，光绪元年（1875年），清与云南回民社群发生“「杜文秀政权覆亡前后，云南回民战争进入惨烈收束阶段」引发的善后、执行或地方回应构成可由不同材料追踪的后续节点”。账本将其背景说明为核心事件「杜文秀政权覆亡前后，云南回民战争进入惨烈收束阶段」改变了相关行动者面对的资源、秩序或风险。，可见结果为死亡、迁徙与复垦规模的估计须保持不确定性。这一节点虽可放入改革、外交或革命的宏观年表，却不能脱离省份财政、地方官、商会、学校、军队、家庭与交通网络来理解；同一上谕、条约或战事在不同地域的实际后果并不一致。

核心材料为《清德宗实录》光绪元年相关条；军机处档、战事方略及地方奏报；相关地区的档案、志书、契约、报刊或对方材料。它们可以支持年序、制度文本、机构设立或某些地点的行动，却不能自动代表全体居民的收入、教育、身份、性别关系或政治意愿。账本提示“战后、改革后或条约后的记忆、地方记录和官方总结常具有事后选择性；后续影响须按地点、群体和时间分层，不作整体化推断。”。因此，官方统计、外交措辞、报刊论战、教会叙述和军报数字必须分别置于生成它们的立场与制度目的之中，不能互相替代。

在地方尺度上，这一事件应被看作资源和信息重新配置的过程：经费能否落实，取决于税源和中介；学校、警政或新军能否运行，取决于人员、土地和日常协作；家庭与社群对改革的经历，则可能因性别、职业、族属、居处与迁徙而不同。现存材料只照亮少数路径，叙事须保留被登记者与未被登记者之间的差距。

账本条目\texttt{EQ-1876-CHEFOO-CONVENTION-PRELUDE}记录，光绪二年（1876年），清与英国发生“「《烟台条约》签订，扩大英国在华通商和内地活动权利」之前的条件、信息和争议已通过中央与地方材料显现，构成该事件可追踪的前史节点”。账本将其背景说明为马嘉理事件与英国外交压力推动谈判，可见结果为为核心事件「《烟台条约》签订，扩大英国在华通商和内地活动权利」提供时间、资源和行动者条件的检索入口，不能替代对核心事件本身的取证。。这一节点虽可放入改革、外交或革命的宏观年表，却不能脱离省份财政、地方官、商会、学校、军队、家庭与交通网络来理解；同一上谕、条约或战事在不同地域的实际后果并不一致。

核心材料为《清德宗实录》光绪二年相关条；《筹办夷务始末》相关年卷与中外照会、条约文本；同时期地方档案、地方志、日记或对方材料应按具体地区补检。它们可以支持年序、制度文本、机构设立或某些地点的行动，却不能自动代表全体居民的收入、教育、身份、性别关系或政治意愿。账本提示“官修编年和地方材料的记录密度与立场并不相同；本行不将条件、传闻、呈报或后见解释直接等同为确定因果。”。因此，官方统计、外交措辞、报刊论战、教会叙述和军报数字必须分别置于生成它们的立场与制度目的之中，不能互相替代。

在地方尺度上，这一事件应被看作资源和信息重新配置的过程：经费能否落实，取决于税源和中介；学校、警政或新军能否运行，取决于人员、土地和日常协作；家庭与社群对改革的经历，则可能因性别、职业、族属、居处与迁徙而不同。现存材料只照亮少数路径，叙事须保留被登记者与未被登记者之间的差距。

账本条目\texttt{EQ-1876-CHEFOO-CONVENTION-DECISION}记录，光绪二年（1876年），清与英国发生“围绕「《烟台条约》签订，扩大英国在华通商和内地活动权利」的命令、调度、照会或呈报形成独立的中央—地方行动文书链”。账本将其背景说明为核心事件「《烟台条约》签订，扩大英国在华通商和内地活动权利」要求相关机构处理人事、资源、军事、外交或司法事务。，可见结果为条约权利的地区落实和地方反应需要逐案追踪。这一节点虽可放入改革、外交或革命的宏观年表，却不能脱离省份财政、地方官、商会、学校、军队、家庭与交通网络来理解；同一上谕、条约或战事在不同地域的实际后果并不一致。

核心材料为《清德宗实录》光绪二年相关条；《筹办夷务始末》相关年卷与中外照会、条约文本。它们可以支持年序、制度文本、机构设立或某些地点的行动，却不能自动代表全体居民的收入、教育、身份、性别关系或政治意愿。账本提示“奏折、上谕、部议、军机档或条约可证明行政动作和机构立场，不自动证明所有地区、群体或对手按其意图行动。”。因此，官方统计、外交措辞、报刊论战、教会叙述和军报数字必须分别置于生成它们的立场与制度目的之中，不能互相替代。

在地方尺度上，这一事件应被看作资源和信息重新配置的过程：经费能否落实，取决于税源和中介；学校、警政或新军能否运行，取决于人员、土地和日常协作；家庭与社群对改革的经历，则可能因性别、职业、族属、居处与迁徙而不同。现存材料只照亮少数路径，叙事须保留被登记者与未被登记者之间的差距。

账本条目\texttt{EQ-1876-CHEFOO-CONVENTION-AFTERMATH}记录，光绪二年（1876年），清与英国发生“「《烟台条约》签订，扩大英国在华通商和内地活动权利」引发的善后、执行或地方回应构成可由不同材料追踪的后续节点”。账本将其背景说明为核心事件「《烟台条约》签订，扩大英国在华通商和内地活动权利」改变了相关行动者面对的资源、秩序或风险。，可见结果为条约权利的地区落实和地方反应需要逐案追踪。这一节点虽可放入改革、外交或革命的宏观年表，却不能脱离省份财政、地方官、商会、学校、军队、家庭与交通网络来理解；同一上谕、条约或战事在不同地域的实际后果并不一致。

核心材料为《清德宗实录》光绪二年相关条；《筹办夷务始末》相关年卷与中外照会、条约文本；相关地区的档案、志书、契约、报刊或对方材料。它们可以支持年序、制度文本、机构设立或某些地点的行动，却不能自动代表全体居民的收入、教育、身份、性别关系或政治意愿。账本提示“战后、改革后或条约后的记忆、地方记录和官方总结常具有事后选择性；后续影响须按地点、群体和时间分层，不作整体化推断。”。因此，官方统计、外交措辞、报刊论战、教会叙述和军报数字必须分别置于生成它们的立场与制度目的之中，不能互相替代。

在地方尺度上，这一事件应被看作资源和信息重新配置的过程：经费能否落实，取决于税源和中介；学校、警政或新军能否运行，取决于人员、土地和日常协作；家庭与社群对改革的经历，则可能因性别、职业、族属、居处与迁徙而不同。现存材料只照亮少数路径，叙事须保留被登记者与未被登记者之间的差距。

账本条目\texttt{EQ-1876-NORTH-CHINA-FAMINE-PRELUDE}记录，光绪二年（1876年），清与华北发生“「丁戊奇荒在华北及西北造成广泛饥荒、迁徙和救济行动」之前的条件、信息和争议已通过中央与地方材料显现，构成该事件可追踪的前史节点”。账本将其背景说明为旱灾、粮价、交通和行政能力共同塑造灾情，可见结果为为核心事件「丁戊奇荒在华北及西北造成广泛饥荒、迁徙和救济行动」提供时间、资源和行动者条件的检索入口，不能替代对核心事件本身的取证。。这一节点虽可放入改革、外交或革命的宏观年表，却不能脱离省份财政、地方官、商会、学校、军队、家庭与交通网络来理解；同一上谕、条约或战事在不同地域的实际后果并不一致。

核心材料为《清德宗实录》光绪二年相关条；地方志、督抚奏折及地方档案；同时期地方档案、地方志、日记或对方材料应按具体地区补检。它们可以支持年序、制度文本、机构设立或某些地点的行动，却不能自动代表全体居民的收入、教育、身份、性别关系或政治意愿。账本提示“官修编年和地方材料的记录密度与立场并不相同；本行不将条件、传闻、呈报或后见解释直接等同为确定因果。”。因此，官方统计、外交措辞、报刊论战、教会叙述和军报数字必须分别置于生成它们的立场与制度目的之中，不能互相替代。

在地方尺度上，这一事件应被看作资源和信息重新配置的过程：经费能否落实，取决于税源和中介；学校、警政或新军能否运行，取决于人员、土地和日常协作；家庭与社群对改革的经历，则可能因性别、职业、族属、居处与迁徙而不同。现存材料只照亮少数路径，叙事须保留被登记者与未被登记者之间的差距。

账本条目\texttt{EQ-1876-NORTH-CHINA-FAMINE-DECISION}记录，光绪二年（1876年），清与华北发生“围绕「丁戊奇荒在华北及西北造成广泛饥荒、迁徙和救济行动」的命令、调度、照会或呈报形成独立的中央—地方行动文书链”。账本将其背景说明为核心事件「丁戊奇荒在华北及西北造成广泛饥荒、迁徙和救济行动」要求相关机构处理人事、资源、军事、外交或司法事务。，可见结果为伤亡总数争议极大，应按地区、材料与估算口径说明。这一节点虽可放入改革、外交或革命的宏观年表，却不能脱离省份财政、地方官、商会、学校、军队、家庭与交通网络来理解；同一上谕、条约或战事在不同地域的实际后果并不一致。

核心材料为《清德宗实录》光绪二年相关条；地方志、督抚奏折及地方档案。它们可以支持年序、制度文本、机构设立或某些地点的行动，却不能自动代表全体居民的收入、教育、身份、性别关系或政治意愿。账本提示“奏折、上谕、部议、军机档或条约可证明行政动作和机构立场，不自动证明所有地区、群体或对手按其意图行动。”。因此，官方统计、外交措辞、报刊论战、教会叙述和军报数字必须分别置于生成它们的立场与制度目的之中，不能互相替代。

在地方尺度上，这一事件应被看作资源和信息重新配置的过程：经费能否落实，取决于税源和中介；学校、警政或新军能否运行，取决于人员、土地和日常协作；家庭与社群对改革的经历，则可能因性别、职业、族属、居处与迁徙而不同。现存材料只照亮少数路径，叙事须保留被登记者与未被登记者之间的差距。

账本条目\texttt{EQ-1876-NORTH-CHINA-FAMINE-AFTERMATH}记录，光绪二年（1876年），清与华北发生“「丁戊奇荒在华北及西北造成广泛饥荒、迁徙和救济行动」引发的善后、执行或地方回应构成可由不同材料追踪的后续节点”。账本将其背景说明为核心事件「丁戊奇荒在华北及西北造成广泛饥荒、迁徙和救济行动」改变了相关行动者面对的资源、秩序或风险。，可见结果为伤亡总数争议极大，应按地区、材料与估算口径说明。这一节点虽可放入改革、外交或革命的宏观年表，却不能脱离省份财政、地方官、商会、学校、军队、家庭与交通网络来理解；同一上谕、条约或战事在不同地域的实际后果并不一致。

核心材料为《清德宗实录》光绪二年相关条；地方志、督抚奏折及地方档案；相关地区的档案、志书、契约、报刊或对方材料。它们可以支持年序、制度文本、机构设立或某些地点的行动，却不能自动代表全体居民的收入、教育、身份、性别关系或政治意愿。账本提示“战后、改革后或条约后的记忆、地方记录和官方总结常具有事后选择性；后续影响须按地点、群体和时间分层，不作整体化推断。”。因此，官方统计、外交措辞、报刊论战、教会叙述和军报数字必须分别置于生成它们的立场与制度目的之中，不能互相替代。

在地方尺度上，这一事件应被看作资源和信息重新配置的过程：经费能否落实，取决于税源和中介；学校、警政或新军能否运行，取决于人员、土地和日常协作；家庭与社群对改革的经历，则可能因性别、职业、族属、居处与迁徙而不同。现存材料只照亮少数路径，叙事须保留被登记者与未被登记者之间的差距。
